\documentclass[11pt]{article}
\usepackage[a4paper,margin=27mm]{geometry}
\usepackage[T1]{fontenc}
\usepackage{lmodern,microtype,amsmath,amssymb,amsthm,mathtools,enumitem}
\usepackage[hidelinks,hypertexnames=false]{hyperref}
\newtheorem{theorem}{Theorem}
\newtheorem{lemma}[theorem]{Lemma}
\newtheorem{proposition}[theorem]{Proposition}
\newtheorem{corollary}[theorem]{Corollary}
\theoremstyle{definition}\newtheorem{definition}[theorem]{Definition}
\newtheorem{example}[theorem]{Example}
\theoremstyle{remark}\newtheorem{remark}[theorem]{Remark}
\newcommand{\R}{\mathbb R}
\newcommand{\Dscr}{\mathcal D}
\newcommand{\supp}{\operatorname{supp}}
\newcommand{\debt}{\mathrm{debt}}
\newcommand{\cred}{\mathrm{credit}}
\newcommand{\sil}{\mathrm{silence}}
\newcommand{\Jid}{\mathcal J}
\newcommand{\Kid}{\mathcal K}
\newcommand{\Nid}{\mathcal N}
\newcommand{\Iid}{\mathcal I}
\newcommand{\Band}{\mathrm B}
\newcommand{\Uf}{\mathcal U}
\newcommand{\Bfam}{\mathcal B}
\newcommand{\Hset}{\mathcal H}
\newcommand{\Lo}[3]{\mathrm L^{#1}_{#2}(#3)}
\newcommand{\Hi}[3]{\mathrm H^{#1}_{#2}(#3)}
\newcommand{\pw}{\mathrm{pw}}
\newcommand{\Tw}{\textup{(T1$^\circ_{\mathrm w}$)}}
\newcommand{\To}{\textup{(T1$^\circ$)}}
\newcommand{\Tos}{\textup{(T1$^{\circ*}$)}}
\newcommand{\Ts}{\textup{(T1$^{\mathrm s}$)}}
\newcommand{\Tcc}{\textup{(T1$^{\mathrm{cc}}$)}}
\newcommand{\Tcut}{\textup{(T1$^{\mathrm{cut}}$)}}
\title{Incomplete Rankings and Coverage}
\author{}\date{11 October 2026}
\begin{document}\maketitle

\noindent
This note continues TAW (\emph{Two Additive Worlds Without Completeness})
and the three-world note (\emph{Three Complete Rankings Cannot Cover}). It
uses results of CCF (\emph{Closed Choice Forces Fineness}), FSE
(\emph{Fineness Spreads Along Edges}), CBOI (\emph{Closed Choice Beyond
Ordinary Integrals}) and EDL (\emph{Edge Dichotomy and Leverage}). There are three worlds, every
ranking is a preorder satisfying (G), and the selection axiom (S) holds;
completeness is not assumed unless it is stated. We ask what forces the
vacuous model $V$ to be dominated when rankings may be incomplete, and what
a configuration in which $V$ is undominated would need.

\section{Summary}\label{sec:summary}

Two results frame the note. By the three-world note (Theorem~12 and
Corollary~13), if all three rankings are complete and satisfy (G) and (IB),
and (S) holds, then $V$ is dominated: a configuration in which $V$ is
undominated must contain incomparabilities. By TAW (Theorem~71), if two
worlds have one-sided weak extensions (W) of CBOI's finitely additive form
and the third world is any preorder satisfying (IB), then $V$ is dominated;
no completeness is used anywhere. A world with such an extension is called a
\emph{(W)-world} below. The results of this note are the following.

\begin{enumerate}[leftmargin=*,itemsep=3pt]
\item \emph{The endpoint property at a (W)-world}
(Theorem~\ref{thm:EP}). Let world~$2$ be a (W)-world, the bound (W$_2$)
holding at every coherent undominated model, $V$ included; let worlds $1$
and $3$ be arbitrary preorders satisfying (IB); and suppose $V$ is
undominated. Then both edges at $e_2$ have the endpoint property. No
completeness is used, and no (IB) at world~$2$.
\item \emph{One (W)-world and one complete world}
(Theorem~\ref{thm:onecomplete}). Let world~$2$ be a (W)-world as in item~1,
let world~$3$ be complete with (IB), and let world~$1$ be any preorder with
(IB). Then $V$ is dominated. Consequently
(Corollary~\ref{cor:Wforces}), if every world is a preorder with (IB) and
$V$ is undominated, a (W)-world forces both other worlds to be incomplete.
\item \emph{Simple completeness} (Theorems~\ref{thm:SCext}
and~\ref{thm:SC}). CCF's Theorem~15 holds for every world that satisfies
(IB), the condition (SC) of Definition~\ref{def:SC} (its debt order and its
credit order are complete, and every pair is SI or weakly AC), and the
endpoint property on both of its edges. In item~2, completeness of world~$3$
may be replaced by (IB), (SC) and a pointwise edge-openness condition along
$[e_3,e_1]$.
\item \emph{Theorems A and B} (Theorems~\ref{thm:A} and~\ref{thm:B}). Let
every ranking be a preorder with (G) and (IB), assume (S), and assume the
pointwise edge-openness condition \To{} of Definition~\ref{def:T1o} at every
world; it holds automatically at complete worlds. If $V$ is undominated,
then every edge has the endpoint property and every world is fine. If
moreover at most one world is incomplete, then $V$ is dominated.
\item \emph{The shape of a configuration with $V$ undominated}
(Corollary~\ref{cor:C}). If every ranking is a preorder with (G) and (IB),
(S) holds and $V$ is undominated, then either (R-a) \To{} fails at some
world, which is then incomplete, or (R-b) every edge has the endpoint
property, every world is fine, and at least two worlds violate (SC) and
have no one-sided weak extension of CBOI's form (equivalently, by EDL's
Theorem~2, carry leverage sequences).
\item \emph{One incomplete world without \To}
(Proposition~\ref{prop:F4}). If worlds $2$ and $3$ are complete, world~$1$
is any preorder with (IB), and $V$ is undominated, then neither complete
world has the endpoint property on both of its edges.
\item \emph{Exact cancellations} (Theorem~\ref{thm:NC},
Corollaries~\ref{cor:NCext} and~\ref{cor:NCclasses},
Lemma~\ref{lem:smallcanc}, Examples~\ref{ex:K4} and~\ref{ex:half5}). Encode the change between
two regular models at world $v$ by the change of the indicators of the
events ``rank $\ge1$'' and ``rank $=2$'', a function on $H'\times\{1,2\}$. At a fine world with an SI
pair with non-null debt part and a weakly AC pair with non-null credit part,
a weak extension of CBOI's form valid for all regular models exists if and
only if (NC): no nonnegative combination of accepted change vectors is
$\le-1$. The proof is a minimax argument on finite systems together with
Tychonoff's theorem; it uses no completeness, no (IB), no Savage step and no
exchange rate. At such a world (NC) holds under (SC) with the endpoint
property, under a generalised cancellation condition, for filter-Pareto
orders, and if the world is contained in a complete preorder with (G) and
(IB). The order axioms do not imply (NC) (a ten-vector example) or the
stronger (NC$^+$) (a four-vector example); neither example is fine. A
cancellation inside one layer needs at least four accepted changes.
\item \emph{Two worlds violating (NC)} (Theorem~\ref{thm:twoNC},
Corollary~\ref{cor:twoNC}). Under the hypotheses of Theorem~A, if $V$ is
undominated, at least two worlds violate (NC); so two worlds satisfying (NC)
force domination. More generally, if $V$ is undominated, every edge has the
endpoint property and every world satisfies a weak form \Tcc{} of \To, then
at most one world with a fine credit order satisfies (NC).
\item \emph{(S) does not imply \To} (Theorem~\ref{thm:conf},
Proposition~\ref{prop:confT1}). In the configurations
$\mathrm{Conf}(\theta)$ worlds $2$ and $3$ are complete and additive,
world~$1$ is the Pareto order of its expected loss and of the same loss with
silence penalised by $\theta$, and the selections are precise. All axioms
hold, and \To{} fails at world~$1$ along $[e_1,e_j]$ if and only if
$\{\theta=1\}\cap D_{1j}$ is non-null; $V$ is dominated there.
\item \emph{What the proof of Theorem~A uses} (Lemma~\ref{lem:SACo},
Proposition~\ref{prop:audit}, Theorems~\ref{thm:Astar}
and~\ref{thm:Bstar}). Theorem~A needs \To{} only through (E2), a property of
the debt order; through \Tcc, \To{} for approaches whose selections give no
credit beyond that of the improved model; through \Tcut{} at one world along
one edge, to close a cut; and, for the fineness of the credit orders,
through the instances of CCF's Lemma~6. FSE's Lemma~11 is replaced by a
completeness-free lemma. Hence $V$ is dominated if at most one world is
incomplete and that world satisfies the stripped condition \Ts, which
the K-world shows to be weaker than \To.
\item \emph{Crossing cuts and the refined shape}
(Theorem~\ref{thm:crossing}, Corollary~\ref{cor:refined}). If $V$ is
undominated, (E2) and \Tcc{} hold everywhere and some edge lacks the
endpoint property, then every
ordered pair of worlds has a cut approached only by \emph{crossing}
selections, which give both more credit and more debt than the selection at
the cut. Consequently, if $V$ is undominated, one of the following holds:
(R-a1) an incomplete world violates (E2) or \Tcc; (R-a2$''$) crossing cuts
for all six ordered pairs, at which \Tcut{} fails, with all worlds
incomplete; (R-b$'$) the endpoint property on every edge and at least two
worlds that are not (W)-worlds, which violate (NC) when the credit orders
are fine.
\end{enumerate}

\emph{What is settled.} Incomparabilities are necessary for $V$ to be
undominated (the three-world note). Two (W)-worlds force domination (TAW),
and so do a single (W)-world and a single complete world (item~2). At most
one incomplete world forces domination if that world satisfies \Ts, in
particular if it satisfies \To{} (items~4 and~10). (S), (G) and (IB) do not
imply \To{} (item~9). At a fine world with an SI pair and a weakly AC pair,
a weak extension exists exactly when (NC) holds (item~7).

\emph{What remains open.} The open core is the three alternatives of
Corollary~\ref{cor:refined}. (R-a1): whether $V$ undominated forces (E2)
and \Tcc{} at the incomplete worlds. All failures of \To{} known under (S)
occur at world~$1$ of some $\mathrm{Conf}(\theta)$, a (W)-world; they never
violate \Tcc{} or \Tcut, and in the K-world not (E2) either; so the only
known mechanism that could produce (R-a1) is a failure of (E2).
(R-a2$''$): whether any configuration satisfying (S) has a crossing cut;
none is known. (R-b$'$): whether there can be two worlds with the endpoint
property on their edges, (E2), \Tcc{} and fine debt orders that are not
(W)-worlds. By Theorem~\ref{thm:twoNC}, at each such world either the credit
order is not fine or (NC) fails; if the credit orders are fine (for instance
under \Tos), the question is whether fine worlds can violate (NC), and we
know no fine world with (IB) that does. This is the extension problem (EXT)
of Section~\ref{sec:EXT}. Further: exactly one incomplete world violating (E2)
or \Tcc{} (item~6 is the part we can prove without any form of \To), and a
(W)-world whose two partners are both incomplete
(Remark~\ref{rem:OWgeneral}).

The results of Sections~\ref{sec:onew}--\ref{sec:shape} come from working
notes that were checked by independent referees; their corrections are
incorporated, and results taken from those notes are proved in full here.
The proof of Theorem~\ref{thm:rich} follows the route of the three-world
note's Proposition~6, and the audit of Section~\ref{sec:T1s} is carried out
for the proof of Theorem~\ref{thm:A} given here.

\section{Setting}\label{sec:setting}

\subsection{Problems, models and axioms}\label{sec:axioms}

We use the setting and notation of FSE, Section~1, and of the three-world
note, Section~2, with $n=3$ worlds. $\Delta\coloneqq\{p\in\R^3:p\ge0,\
\sum_ip_i=1\}$ has vertices $e_1,e_2,e_3$, and $\supp p\coloneqq
\{i:p_i>0\}$. Binary problems $\{g,-g\}$ are represented by $g\in
H\coloneqq\{g\in\R^3:g_1>0\}$, with normalized Gaussian measure~$\mu$;
coordinate hyperplanes are null, and $H'$ is the set of mixed problems in
$H$. Models are measurable and identified modulo null sets. A model is
\emph{regular} if it is never contradictory and correct on unanimous
problems. Coherent models $\Dscr_K$, $K\subseteq\Delta$ nonempty, compact
and convex, are regular; $V\coloneqq\Dscr_\Delta$, and $D_i$ is the ideal
model of world~$i$, which recommends the option good at~$i$ on every
problem. All models below are regular.

On a mixed problem a regular model gives the option good at world~$v$
(\emph{credit} at $v$, rank~$0$), silence (rank~$1$), or the option bad at
$v$ (\emph{debt} at $v$, rank~$2$). A regular model $M$ is determined by its
world-$v$ rank function $r_v(\cdot,M)\colon H'\to\{0,1,2\}$, and every
measurable function $H'\to\{0,1,2\}$ is the rank function of a regular
model. For measurable $Q\subseteq H'$, the models $M+\debt_v(Q)$,
$M+\cred_v(Q)$ and $M+\sil(Q)$ agree with $M$ off $Q$ and have world-$v$
rank $2$, $0$, $1$ on $Q$; $M+\sil(Q)$ is also called $M$ \emph{silenced on}
$Q$. $V$ is silent on $H'$, and $D_v$ has world-$v$ rank $0$ on $H'$. We
write $X\le^v_\pw Y$ if $r_v(\cdot,X)\le r_v(\cdot,Y)$ on $H'$, and call $Y$
a \emph{strict pointwise $v$-improvement} of $C$ if $Y\le^v_\pw C$ and
$r_v(\cdot,Y)<r_v(\cdot,C)$ on a non-null set.

At each world, $\preceq_i$ is a preorder on models; smaller is better.
$X\prec_iY$ means that $X\preceq_iY$ and not $Y\preceq_iX$, and
$X\parallel_iY$ means that neither $X\preceq_iY$ nor $Y\preceq_iX$. In a
preorder, $X\preceq_iY\prec_iZ$ and $X\prec_iY\preceq_iZ$ both imply
$X\prec_iZ$. The ranking of world~$i$ is \emph{complete} if no two models
are incomparable. (G) at world $i$ is guidance monotonicity: $X\le^i_\pw Y$
implies $X\preceq_iY$ (\emph{weak (G)}), and $X\prec_iY$ if moreover $X$ is a
strict pointwise $i$-improvement of $Y$ (\emph{strict (G)}). (IB) at
world~$i$ is independence under replacement of shared advice; we use it in
the form of the three-world note, Section~2.1:
\begin{enumerate}[label=(IB$'$),leftmargin=*]
\item If $X,Y,X',Y'$ are regular, $\{X\ne Y\}=\{X'\ne Y'\}$, and $X=X'$ and
$Y=Y'$ on this set, then $X\preceq_iY\iff X'\preceq_iY'$ and
$X\prec_iY\iff X'\prec_iY'$.
\end{enumerate}

(S) supplies coherent undominated models $C_p$, $p\in\Delta$, competitive on
$\supp p$, with closed sections $A_i(X)\coloneqq\{p:C_p\preceq_iX\}$ for
every model $X$ and every world $i$. A model $C$ is \emph{competitive} on a
set $\mathcal W$ of worlds if for every model $X$ some $i\in\mathcal W$ has
$C\preceq_iX$. A model $M$ is \emph{dominated} if some model $X$ has
$X\preceq_iM$ at every world and $X\prec_iM$ at some world. By CBOI's
Lemma~1 every model $X$ has a selected $C_q$ with $C_q\preceq_iX$ at every
world; its proof uses Brouwer's theorem and no completeness. So a dominated
model is dominated by a selected coherent model. By the face lemma (FSE,
Section~1.1) the credal set of $C_p$ lies in the face spanned by $\supp p$;
in particular $C_{e_i}=D_i$. For $q$ on an edge $[e_v,e_w]$, $C_q$ is a
\emph{selection on the edge}; if $q$ lies in the open edge $(e_v,e_w)$, then
$C_q$ is competitive on $\{v,w\}$. If $V$ is undominated, it is competitive
on $\{1,2,3\}$ (CCF, Section~1). We use the consequences (T1) and (T2) of
(S) (FSE, Section~1.1; the three-world note, Section~2.1):
\begin{enumerate}[label=(T\arabic*),leftmargin=*,itemsep=2pt]
\item If world $k$ is complete, $Y$ is regular and $Y\prec_kC_{p^*}$, then
$Y\prec_kC_p$ for all $p$ near $p^*$.
\item If $p^*\in[e_v,e_w]$, $X$ is regular and $X\prec_wC_{p^*}$, then
$C_p\preceq_vX$ for all $p$ in the open edge $(e_v,e_w)$ near $p^*$.
\end{enumerate}
(T1) holds because $A_k(Y)$ is closed and misses $p^*$, so that
$C_p\not\preceq_kY$ near $p^*$, and completeness turns this into
$Y\prec_kC_p$. (T2) holds because $A_w(X)$ is closed and misses $p^*$, and
$C_p$ is competitive on $\{v,w\}$. Only (T1) uses completeness. In place of
completeness we shall sometimes assume the following.

\begin{definition}\label{def:T1o}
World $k$ satisfies \emph{\To{} along the edge $[e_k,e_j]$} if for every
$p^*\in[e_k,e_j]$ and every strict pointwise $k$-improvement $Y$ of
$C_{p^*}$ there is a neighbourhood $N$ of $p^*$ in $[e_k,e_j]$ such that
$Y\prec_kC_p$ for every $p\in N$. It satisfies \emph{\Tw{} along
$[e_k,e_j]$} if the same holds with $Y\preceq_kC_p$ in place of
$Y\prec_kC_p$. \Tos{} means that every world satisfies \To{} along both of
its edges.
\end{definition}

By strict (G), $Y\prec_kC_{p^*}$ in Definition~\ref{def:T1o}; so (T1) gives
\To{} at a complete world, and \To{} implies \Tw. Section~\ref{sec:T1o}
discusses the condition further.

\subsection{Debt, credit and pairs}\label{sec:pairs}

For measurable $B,C\subseteq H'$ (FSE, Section~1.2),
\begin{gather*}
 B\preceq^v_dC\ :\Longleftrightarrow\ V+\debt_v(B)\preceq_vV+\debt_v(C),\\
 B\preceq^v_cC\ :\Longleftrightarrow\ V+\cred_v(C)\preceq_vV+\cred_v(B).
\end{gather*}
So $B\preceq^v_dC$ says that debt on $B$ is no worse than debt on $C$, and
$B\preceq^v_cC$ says that credit on $B$ is worth no more than credit on $C$.
The strict parts are $\prec^v_d$ and $\prec^v_c$. Both relations are
transitive. By (IB$'$), $V$ may be replaced by any regular model silent on
$B\cup C$, and the comparison depends only on $B\setminus C$ and
$C\setminus B$; overlapping sets may be compared. By weak (G),
$B\subseteq C$ implies $B\preceq^v_dC$ and $B\preceq^v_cC$, and by strict (G)
both are strict if $C\setminus B$ is non-null.

A \emph{pair} $(A,T)$ at world $v$ consists of disjoint measurable
$A,T\subseteq H'$ (debt part and credit part). It is \emph{SI} if
$V+\debt_v(A)+\cred_v(T)\prec_vV$, \emph{weakly SI} if
$V+\debt_v(A)+\cred_v(T)\preceq_vV$, \emph{weakly AC} if
$V\preceq_vV+\debt_v(A)+\cred_v(T)$, and \emph{strictly AC} if
$V\prec_vV+\debt_v(A)+\cred_v(T)$. An SI pair is weakly SI and not weakly
AC; a strictly AC pair is weakly AC and not weakly SI. At a complete world a
pair that is not SI is weakly AC, and a pair that is not weakly SI is
strictly AC; without completeness a pair may be neither weakly SI nor
weakly AC. By (IB$'$), $V$ may
be replaced by any regular model silent on $A\cup T$. By strict (G), a pair
$(\varnothing,T)$ with $T$ non-null is SI and a pair $(A,\varnothing)$ with
$A$ non-null is strictly AC.

For non-null $A,T\subseteq H'$, $\Jid^v_A$ ($\Kid^v_T$) is the family of
measurable sets contained in a finite union of measurable sets each
$\preceq^v_dA$ (each $\preceq^v_cT$). These are ideals of the measure
algebra of $H'$ (FSE, Definition~1 and Lemma~2). World $v$ is \emph{fine}
if $H'\in\Jid^v_A\cap\Kid^v_T$ for all non-null $A,T$. Ideals and
ultrafilters below are those of the measure algebra of $H'$, so membership
is read modulo null sets; an ultrafilter \emph{avoids} an ideal if they are
disjoint. We use FSE's Lemma~5 in the following form.
\begin{enumerate}[label=(U\arabic*),leftmargin=*,itemsep=2pt]
\item\label{U:small} (FSE, Lemma~5(a)) Every ultrafilter contains sets of
arbitrarily small measure. If a finite union of measurable sets is in an
ultrafilter, one of them is.
\item\label{U:avoid} (FSE, proof of Lemma~5(b); the three-world note,
(F1)) If $\Iid$ is an ideal and $R\notin\Iid$, some ultrafilter contains $R$
and avoids $\Iid$.
\item\label{U:strict} (FSE, Lemma~5(c)) If $\preceq^v_c$ is complete and an
ultrafilter $\Uf$ avoids $\Kid^v_T$, then $T\prec^v_cB$ for every
$B\in\Uf$. Indeed $B\notin\Kid^v_T$ gives $B\not\preceq^v_cT$, and
completeness of $\preceq^v_c$ gives $T\preceq^v_cB$.
\end{enumerate}
FSE states (U3) for a complete world; its proof uses completeness only of
$\preceq^v_c$, as shown.

\subsection{Regions, frames and edges}\label{sec:frames}

Put
\[
 P_1\coloneqq\{g_2<0,\ g_3<0\},\qquad P_2\coloneqq\{g_2<0<g_3\},\qquad
 P_3\coloneqq\{g_3<0<g_2\}
\]
inside $H$; $P_x$ is the set of mixed problems at which world $x$ is the odd
one out, and $H'=P_1\cup P_2\cup P_3$ up to null sets. For distinct worlds
$v,w$ put $D_{vw}\coloneqq\{g\in H':g_vg_w<0\}=P_v\cup P_w$ and
$t^v_w\coloneqq|g_v|/|g_w|$ on $D_{vw}$, so $t^w_v=1/t^v_w$. On $D_{vw}$
the option good at $v$ is the option bad at $w$, so the world-$v$ and
world-$w$ ranks of every regular model add up to $2$ there (FSE,
equation~(1)):
\begin{equation}
 \cred_v=\debt_w\quad\text{and}\quad\debt_v=\cred_w\qquad\text{on }D_{vw};
 \label{eq:swap}
\end{equation}
off $D_{vw}$ the ranks at $v$ and $w$ are equal. As in FSE, Section~1.4,
$o_K$ is the option good exactly at the worlds in $K$:
\begin{center}
\begin{tabular}{c|ccc|ccc}
 & \multicolumn{3}{c|}{accept $g$} & \multicolumn{3}{c}{accept $-g$}\\
 & option & credit at & debt at & option & credit at & debt at\\\hline
$P_1$ & $o_1$ & $1$ & $2,3$ & $o_{23}$ & $2,3$ & $1$\\
$P_2$ & $o_{13}$ & $1,3$ & $2$ & $o_2$ & $2$ & $1,3$\\
$P_3$ & $o_{12}$ & $1,2$ & $3$ & $o_3$ & $3$ & $1,2$
\end{tabular}
\end{center}
For $M$ regular and measurable $Q\subseteq P_x$, $M+o(Q)$ gives the option
$o$ on $Q$ and agrees with $M$ elsewhere. On each $P_x$ the joint law of two
ratios $|g_a|/|g_b|$ and $|g_c|/|g_d|$ with $\{a,b\}\ne\{c,d\}$ has a
positive density on $(0,\infty)^2$, and the law of $t^v_w$ on $P_v$ and on
$P_w$ has a positive density on $(0,\infty)$; level sets of these ratios
are null. So every subset of $P_x$ defined by a nonempty open condition on
one or two such ratios is non-null; we use this without comment. A
\emph{band} of the edge $[e_v,e_w]$ is a set $D_{vw}\cap\{s<t^v_w<s'\}$,
$0\le s<s'\le\infty$. For $0<a,b<\infty$,
\[
 \Lo vwa\coloneqq D_{vw}\cap\{t^v_w<a\},\qquad
 \Hi vwb\coloneqq D_{vw}\cap\{t^v_w>b\}
\]
are the \emph{low cone} and the \emph{high tail} of world $v$ on
$[e_v,e_w]$.

\emph{Edge models} (FSE, Section~1.5). By the face lemma a selection on
$[e_v,e_w]$ has a credal interval on that edge, and there are
$0\le\beta\le\alpha\le\infty$ such that it gives credit at $v$ on
$D_{vw}\cap\{t^v_w>\alpha\}$, silence on $D_{vw}\cap\{\beta<t^v_w<\alpha\}$,
debt at $v$ on $D_{vw}\cap\{t^v_w<\beta\}$, and the option good at $v$ and
$w$ off $D_{vw}$. It is written $(\alpha,\beta)=(\alpha^{(v)},\beta^{(v)})$
in \emph{world $v$'s frame}; in world $w$'s frame it is
$(1/\beta^{(v)},1/\alpha^{(v)})$. Its credal set contains $e_v$ if and only
if $\beta=0$, and $e_w$ if and only if $\alpha=\infty$; $(0,0)=D_v$ and
$(\infty,\infty)=D_w$. A selection is \emph{two-sided} if
$0<\beta\le\alpha<\infty$; a selection whose credal set contains neither
endpoint of its edge is two-sided and lies in the open edge. In world $v$'s
frame we use the endpoint families
\[
 E_w(b)\coloneqq(\infty,b),\qquad E^w_v(a)\coloneqq(a,0),
 \qquad V_{vw}\coloneqq E_w(0)=E^w_v(\infty)=\Dscr_{[e_v,e_w]},
\]
with $E_w(\infty)=D_w$ and $E^w_v(0)=D_v$. For $0\le b<b'\le\infty$ the
models $E_w(b)$ and $E_w(b')$ differ exactly on $D_{vw}\cap\{b<t^v_w<b'\}$,
where the first is silent and the second gives debt at $v$, that is, credit
at $w$; so $E_w(b)\prec_vE_w(b')$ and $E_w(b')\prec_wE_w(b)$ by strict (G)
(CCF, equation~(2)).

The edge $[e_v,e_w]$ has the \emph{endpoint property} (EP$_{vw}$) if every
selection on it has $e_v$ or $e_w$ in its credal set (CCF,
Definition~2). We use the following results, stated in CCF for world~$1$
and an edge $[e_1,e_j]$.
\begin{enumerate}[label=(P\arabic*),leftmargin=*,itemsep=2pt]
\item\label{P:param} (CCF, Lemma~3) Assume (G) at worlds $v$ and $w$, (S)
and (EP$_{vw}$). Put $M_s\coloneqq E^w_v(-1/s)$ for $s<0$,
$M_{-\infty}\coloneqq D_v$ and $M_s\coloneqq E_w(s)$ for $s\ge0$. There is a
continuous map $s\colon[e_v,e_w]\to[-\infty,\infty]$ with $C_q=M_{s(q)}$,
$s(e_v)=-\infty$ and $s(e_w)=\infty$. Every $M_s$ is selected; if $s$ is
finite, only at points of the open edge. In particular $E_w(b)$ and
$E^w_v(a)$, $0<a,b<\infty$, and $V_{vw}$ are competitive on $\{v,w\}$.
\item\label{P:approach} (CCF, Lemma~4) Under the hypotheses of
\ref{P:param}, let $s^*\in[-\infty,\infty]$ and $s_n\to s^*$ with all
$s_n<s^*$ or all $s_n>s^*$. Then there is $q^*$ with $C_{q^*}=M_{s^*}$ such
that every neighbourhood $N$ of $q^*$ has the property: for some $n_0$,
every $s$ between $s_{n_0}$ and $s^*$ (inclusive) is $s(q)$ for some $q\in
N$.
\item\label{P:ends} (FSE, Lemma~6) Assume (G) at worlds $v$ and $w$ and
(S), and write $C_q=(\alpha_q,\beta_q)$ in world $v$'s frame. (a) If
$C_{q^*}=D_w$ and $q_m\to q^*$ along the edge, then $\beta_{q_m}\to\infty$;
if $C_{q^*}=D_v$, then $\alpha_{q_m}\to0$. (b) $\{q\in[e_v,e_w]:C_q=D_w\}$
is closed, contains $e_w$ and not $e_v$. (c) $\{q:C_q=D_v\}$ is closed,
contains $e_v$ and not $e_w$, and it has a boundary point $q_0$ that is a
limit of points $q_m$ of the open edge with $C_{q_m}\ne D_v$.
\end{enumerate}
For a \emph{frame} $(k,l)$, that is, an ordered pair of distinct worlds,
the \emph{far end} of the frame is the point $q_e$ of
$\{q\in[e_k,e_l]:C_q=D_l\}$ nearest to $e_k$ (\ref{P:ends}(b)); every $q$
between $e_k$ and $q_e$ has $C_q\ne D_l$. The far end is \emph{rich} if
there are $q_m\to q_e$ strictly between $e_k$ and $q_e$ with
$\alpha^{(k)}_{q_m}<\infty$ (the three-world note, Definition~5). In world
$l$'s frame this says $\beta^{(l)}_{q_m}>0$: the selections $C_{q_m}$ give
debt at $l$ on a non-null set.

\emph{Relabelling.} The axioms are symmetric in the labels of the worlds,
and every object above is a set of problems. So every statement made in
CCF, FSE, TAW or the three-world note for particular worlds holds for any
permutation of the labels; the change of representing half-space is
harmless, as in the proof of CCF's Theorem~16 (FSE, Section~1.3; the
three-world note, Section~2.4). We use this without further comment.

\subsection{Losses of finitely additive form}\label{sec:losses}

A loss of \emph{CBOI's form} at world $k$ is $L_k(M)\coloneqq
\sum_{l=1}^3\nu_{kl}\{g:r_k(g,M)\ge l\}$, where the $\nu_{kl}$ are finite
nonnegative finitely additive measures (charges) on the measurable subsets
of $H$ with
\begin{equation}
 \nu_{kl}(A)=0\text{ if }\mu(A)=0,\qquad\nu_{kl}(A)>0\text{ if }\mu(A)>0.
 \label{eq:positive}
\end{equation}
For regular $M$, $L_k(M)=\nu_{k1}(R^k_1(M))+\nu_{k2}(R^k_2(M))$ with
$R^k_l(M)\coloneqq\{g\in H':r_k(g,M)\ge l\}$, since regular models have rank
$0$ on unanimous problems and never rank $3$. The one-sided bound is
\begin{equation}
 B\text{ coherent undominated},\ X\text{ regular},\ B\preceq_kX
 \ \Longrightarrow\ L_k(B)\le L_k(X).
 \tag{W$_k$}\label{eq:W}
\end{equation}
It is required for \emph{every} coherent undominated $B$; when $V$ is
undominated this includes $B=V$, which need not be a selected model. A world
with such a loss is a \emph{(W)-world}. By EDL's Theorem~2, world $k$ is a
(W)-world if and only if it carries no leverage sequence in EDL's sense.

If regular $M,M'$ agree off a measurable $Q\subseteq H'$, then by finite
additivity (TAW, Lemma~2)
\[
 L_k(M')-L_k(M)=\sum_{l=1}^2\bigl[\nu_{kl}(Q\cap R^k_l(M'))
 -\nu_{kl}(Q\cap R^k_l(M))\bigr].
\]
In particular, if $M$ is silent on $Q$, changing its advice on $Q$ to credit
at $k$ lowers $L_k$ by $\nu_{k1}(Q)$, and changing it to debt at $k$ raises
$L_k$ by $\nu_{k2}(Q)$; a change on several disjoint sets adds up. Every
non-null $R$ contains non-null sets of arbitrarily small $\nu_{kl}$-mass:
split $R$ into $N$ disjoint non-null sets ($\mu$ is atomless); one of them
has mass at most $\nu_{kl}(R)/N$. Finally (CBOI, Lemma~3, as in the proof of
TAW's Lemma~3, which treats each charge separately; \cite{YH}): for each
$l$ there are an integrable $c_{kl}>0$ a.e.\ and measurable sets
$Z_1\supseteq Z_2\supseteq\cdots$ with $\mu(Z_m)\to0$ such that
$\nu_{kl}(Q)=\int_Qc_{kl}\,d\mu$ for every measurable $Q\subseteq
H\setminus Z_m$ and every $m$.

\subsection{Results used}\label{sec:inputs}

Besides \ref{U:small}--\ref{U:strict} and \ref{P:param}--\ref{P:ends} we
use the following, with the numbering of the sources.
\begin{enumerate}[label=(R\arabic*),leftmargin=*,itemsep=3pt]
\item\label{R:E1} (FSE, Lemma~7 (E1)) Assume (G) at worlds $v$ and $w$,
(IB) at $v$, and (S). For every $c>0$ and every non-null $Q\subseteq
D_{vw}\cap\{t^v_w>c\}$ there is $a>0$ with $\Lo vwa\preceq^v_cQ$.
\item\label{R:yardc} (FSE, Lemma~8(b)) Assume (G) at every world, (IB) at
world~$v$, and (S); let $\{v,w,u\}=\{1,2,3\}$. For every non-null
$T\subseteq H'$ there is $a>0$ such that every measurable subset of $\Lo
vwa$ and every measurable subset of $\Lo vua$ is $\preceq^v_cT$.
\item\label{R:XOR} (FSE, Lemma~25) Assume (IB) at worlds $1$ and $2$. Let
$E=(\alpha_E,\beta_E)$ (world~$1$'s frame) be a selection on the open edge
$(e_1,e_2)$, and put $T_0(E)\coloneqq P_2\cap\{t^1_2>\alpha_E\}$ and
$B_0(E)\coloneqq P_1\cap\{t^1_2<\beta_E\}$; $E$ gives $o_{13}$ on $T_0(E)$
and $o_{23}$ on $B_0(E)$. For $T\subseteq T_0(E)$ and $B\subseteq B_0(E)$
let $V+\tau_{T,B}$ be $V$ with $o_{13}$ on $T$ and $o_{23}$ on $B$. Then
$V+\tau_{T,B}\preceq_1V$ or $V+\tau_{T,B}\preceq_2V$. If (G) holds at
world~$3$, $V$ is undominated and $T\cup B$ is non-null, then not both.
\item\label{R:CCF} (CCF, Theorems~7 and~15) Assume (G) at every world and
(S), and let world $v$ be complete with (IB) and have (EP) on both of its
edges. Then world $v$ is fine, and it has a loss $L_v$ of CBOI's form with
\eqref{eq:positive} such that $D\preceq_vE$ implies $L_v(D)\le L_v(E)$ for
all regular $D,E$; in particular world $v$ is a (W)-world.
\item\label{R:TAW} (TAW, Theorem~71 and Corollary~72) Assume (G) at every
world and (S). If worlds $2$ and $3$ are (W)-worlds and world~$1$ is a
preorder with (IB), then $V$ is dominated. Hence, if every world is a
preorder with (IB) and $V$ is undominated, at most one world is a
(W)-world. Completeness is used nowhere, and (IB) only at world~$1$.
\item\label{R:cut} (The three-world note, Theorem~3) Let world $v$ be
complete, and assume (G) at worlds $v$ and $w$ and (S). For every
band-dense family of measurable subsets of $D_{vw}$ there is a connecting
chain from $D_v$ to $D_w$ (Section~\ref{sec:count}); every such chain has
length at least $2$.
\item\label{R:count} (The three-world note, Lemma~4) The counting lemma,
stated in Section~\ref{sec:count}.
\end{enumerate}
Results of FSE and CCF that are stated for complete worlds and use
completeness only through (T1) are extended to worlds satisfying \To{} in
Lemma~\ref{lem:T1otools}.

\section{Tools without completeness}\label{sec:tools}

\subsection{Connecting chains and counting}\label{sec:count}

Let $v\ne w$ be worlds and $D\coloneqq D_{vw}$. Following the three-world
note (Definition~1), let $\Bfam$ be a family of measurable subsets of $D$.
Two regular models are \emph{$\Bfam$-neighbours} if they agree off some
member of $\Bfam$. A \emph{connecting chain} of length $\ell\ge1$ from $D_v$
to a regular model $X$ is a sequence
\[
 D_v=Y_0,\ X_1,\ Y_1,\ X_2,\ \dots,\ Y_{\ell-1},\ X_\ell=X
\]
of regular models such that $X_i\preceq_vY_{i-1}$ for $1\le i\le\ell$ and,
for $1\le i\le\ell-1$, $Y_i$ agrees with $X_i$ off a set $B_i\in\Bfam$, the
$i$-th \emph{perturbation set}. $\Bfam$ is \emph{band-dense} if every band of
$[e_v,e_w]$ contains a non-null member of $\Bfam$. The set $\mathcal R$ of
regular models reachable from $D_v$ by connecting chains contains $D_v$, is
\emph{down-closed} ($Z\preceq_vX\in\mathcal R$ implies $Z\in\mathcal R$) and
\emph{$\Bfam$-saturated} ($\Bfam$-neighbours of members are members) (the
three-world note, Lemma~2).

\ref{R:count} \emph{Counting lemma} (the three-world note, Lemma~4). Let
$\preceq_v$ be a preorder that satisfies weak (G) and (IB) at world $v$. Let
$D_v=Y_0,X_1,\dots,Y_{\ell-1},X_\ell=D_w$ be a connecting chain of length
$\ell\ge2$ whose perturbation sets $B_1,\dots,B_{\ell-1}$ are contained in a
measurable set $F$. Then there are no pairwise disjoint measurable sets
$K_1,\dots,K_\ell\subseteq D\setminus F$ such that
\begin{equation}
 D_v+\debt_v(B_i)\prec_vD_v+\debt_v(K_l)\qquad
 \text{for all }1\le i\le\ell-1\text{ and }1\le l\le\ell.
 \label{eq:heavy}
\end{equation}
It uses only transitivity, weak (G) and (IB) at world $v$: no completeness,
no (S), and nothing about the other worlds.

\begin{remark}\label{rem:bounded}
\emph{Counting to a bounded target.} The induction in the proof of the
three-world note's Lemma~4 shows $Y_k\prec_v\Pi_{k+1}\coloneqq
D_v+\debt_v(K_1\cup\dots\cup K_{k+1})$ for $0\le k\le\ell-1$. It uses
$\ell\ge2$ (in the step $k=0$, where \eqref{eq:heavy} is applied with
$i=l=1$) and $B_i\cap K_l=\varnothing$, but neither the endpoint of the
chain nor $K_l\subseteq D$. Hence, for a connecting chain of length
$\ell\ge2$ from $D_v$ to an arbitrary regular $X$ with perturbation sets in
$F$, there are no pairwise disjoint measurable
$K_1,\dots,K_\ell\subseteq\{g\in H':r_v(g,X)=2\}\setminus F$ satisfying
\eqref{eq:heavy}: otherwise $X=X_\ell\preceq_vY_{\ell-1}\prec_v\Pi_\ell$,
while $\Pi_\ell\le^v_\pw X$, since $\Pi_\ell$ has rank $2$ only on
$K_1\cup\dots\cup K_\ell$ and rank $0$ elsewhere on $H'$; so
$X\prec_vX$. The hypothesis $\ell\ge2$ cannot be dropped: for $\ell=1$
condition \eqref{eq:heavy} is vacuous, and $K_1=\varnothing$ satisfies it.
\end{remark}

\subsection{Connectivity}\label{sec:cut}

\begin{theorem}\label{thm:cut}
Assume (G) at worlds $v$ and $w$ and (S), and let world $v$ satisfy \Tw{}
along $[e_v,e_w]$. For every band-dense family $\Bfam$ of measurable subsets
of $D$ there is a connecting chain from $D_v$ to $D_w$. Every connecting
chain from $D_v$ to $D_w$ has length at least~$2$.
\end{theorem}
\begin{proof}
This is the proof of the three-world note's Theorem~3, which uses the
completeness of world $v$ only to obtain the comparison in the second
paragraph below. A chain of length $1$ would give
$D_w=X_1\preceq_vY_0=D_v$, whereas $D_v\prec_vD_w$ by strict (G).

Suppose there is no chain from $D_v$ to $D_w$, and let $\mathcal R$ be the
set of models reachable from $D_v$; so $D_w\notin\mathcal R$. Put
$\Sigma\coloneqq\{q\in[e_v,e_w]:C_q\in\mathcal R\}$. By the face lemma,
$e_v\in\Sigma$ and $e_w\notin\Sigma$. Write $C_q=(\alpha_q,\beta_q)$ in
world $v$'s frame and $t\coloneqq t^v_w$.

\emph{The complement of $\Sigma$ is open.} Let $C_q\notin\mathcal R$. Then
$C_q\ne D_v$, so $\alpha_q>0$. The band $D\cap\{0<t<\alpha_q\}$ contains a
non-null $B\in\Bfam$, on which $C_q$ gives silence or debt at $v$. So
$M^-\coloneqq C_q+\cred_v(B)$ is a strict pointwise $v$-improvement of
$C_q$, and \Tw{} gives a neighbourhood $N$ of $q$ in the edge with
$M^-\preceq_vC_{q'}$ for every $q'\in N$. Since $M^-$ is a $\Bfam$-neighbour
of $C_q\notin\mathcal R$, $M^-\notin\mathcal R$. If some $q'\in N$ had
$C_{q'}\in\mathcal R$, down-closure would give $M^-\in\mathcal R$. So
$N\cap\Sigma=\varnothing$.

\emph{$\Sigma$ is open.} Let $C_q\in\mathcal R$. Then $C_q\ne D_w$, so
$\beta_q<\infty$ and $q\ne e_w$. The band $D\cap\{\beta_q<t<\infty\}$
contains a non-null $B\in\Bfam$, on which $C_q$ gives silence or credit at
$v$, that is, by \eqref{eq:swap}, silence or debt at $w$. The model
$M^+\coloneqq C_q+\debt_v(B)$ gives credit at $w$ on $B$; it lies in
$\mathcal R$ by saturation, and $M^+\prec_wC_q$ by strict (G) at world $w$.
By (T2), $C_{q'}\preceq_vM^+$ for all $q'$ in the open edge near $q$, and
these $C_{q'}$ lie in $\mathcal R$ by down-closure. The points of the edge
outside the open edge are $e_v\in\Sigma$ and $e_w\ne q$. So all points of
the edge near $q$ lie in $\Sigma$.

$\Sigma$ is open and closed in the connected edge, contains $e_v$ and not
$e_w$: a contradiction.
\end{proof}

The proof uses \Tw{} at $v$, strict (G) at $v$ and $w$, and (S) through
(T2) and the face lemma; no (IB), and nothing about $V$ or the third world.
If world $v$ is complete, (T1) gives \Tw{}, and Theorem~\ref{thm:cut} is
\ref{R:cut}.

\subsection{Sets light in both layers}\label{sec:light}

\begin{definition}\label{def:light}
For measurable $W,T\subseteq H'$ and a world $v$, $W$ is \emph{$T$-light at
$v$} if $W\cap T=\varnothing$ and $D_v+\debt_v(W)\preceq_vD_v+\sil(T)$.
\end{definition}

This is the three-world note's Definition~7, at world $v$ instead of
world~$2$: the full change from credit to debt at $v$ on $W$ costs no more
than the loss of credit on $T$. By weak (G), measurable subsets of
$T$-light sets are $T$-light.

\begin{lemma}\label{lem:light}
Let world $v$ satisfy (G) and (IB).
\begin{enumerate}[label=(\alph*),leftmargin=*,itemsep=2pt]
\item Let $W_1,\dots,W_N$ be pairwise disjoint and $T_1,\dots,T_N$ pairwise
disjoint measurable sets with $T_i\cap W_k=\varnothing$ for all $i,k$. If
each $W_i$ is $T_i$-light at $v$, then $\bigcup_iW_i$ is
$\bigcup_iT_i$-light at $v$.
\item If $W$ is $T$-light at $v$, then $W\preceq^v_cT$.
\item If $W$ is $T$-light at $v$, $T\preceq^v_cP$, and $W,T,P$ are pairwise
disjoint, then $D_v+\debt_v(W)\preceq_vD_v+\debt_v(P)$, and
$D_v+\debt_v(W)\prec_vD_v+\debt_v(P)$ if $P$ is non-null.
\item If $W,T_a,T_b$ are pairwise disjoint, $(W,T_a)$ is a weakly SI pair of
world $v$, and $W\preceq^v_cT_b$, then $W$ is $(T_a\cup T_b)$-light at $v$.
\end{enumerate}
\end{lemma}
\begin{proof}
(a) is the three-world note's Lemma~9, whose proof uses only (IB) and
transitivity at the world concerned.

(b) By weak (G) and lightness,
$D_v+\sil(W)\preceq_vD_v+\debt_v(W)\preceq_vD_v+\sil(T)$. The outer two
models differ exactly on $W\cup T$, with ranks $1$ and $0$ on $W$, and $0$
and $1$ on $T$; so do $V+\cred_v(T)$ and $V+\cred_v(W)$. By (IB$'$),
$V+\cred_v(T)\preceq_vV+\cred_v(W)$, that is, $W\preceq^v_cT$.

(c) We have $D_v+\debt_v(W)\preceq_vD_v+\sil(T)\preceq_vD_v+\sil(P)
\preceq_vD_v+\debt_v(P)$. The first step is lightness. In the second, the
two models differ exactly on $T\cup P$, with ranks $1$ and $0$ on $T$, and
$0$ and $1$ on $P$, as do $V+\cred_v(P)$ and $V+\cred_v(T)$; so (IB$'$)
turns $T\preceq^v_cP$ into the second step. The third step is weak (G), and
strict (G) if $P$ is non-null.

(d) Put $T\coloneqq T_a\cup T_b$. The models $D_v+\debt_v(W)$ and
$D_v+\sil(T)$ differ exactly on $W\cup T$, with ranks $2$ and $0$ on $W$,
and $0$ and $1$ on $T$; so do $V+\debt_v(W)+\cred_v(T)$ and $V+\cred_v(W)$.
By (IB$'$) it suffices to show
$V+\debt_v(W)+\cred_v(T)\preceq_vV+\cred_v(W)$. The context
$V+\cred_v(T_b)$ is silent on $W\cup T_a$, so (IB$'$) transfers the weakly
SI pair $(W,T_a)$ to
$V+\debt_v(W)+\cred_v(T_a)+\cred_v(T_b)\preceq_vV+\cred_v(T_b)$; and
$V+\cred_v(T_b)\preceq_vV+\cred_v(W)$ is $W\preceq^v_cT_b$.
\end{proof}

\begin{lemma}\label{lem:near}
Let $v\ne w$, assume (G) at worlds $v$ and $w$ and (S), and suppose the far
end of the frame $(w,v)$ is rich. Then for every $c>0$ and every non-null
$T\subseteq D_{vw}\cap\{t^v_w>c\}$ there is $b>0$ such that every measurable
$W\subseteq D_{vw}\cap\{t^v_w<b\}$ is $T$-light at $v$.
\end{lemma}
\begin{proof}
This is the three-world note's Lemma~8, with the frame $(1,2)$ replaced by
$(w,v)$; its proof uses Proposition~6 there only to obtain richness. Let
$q_e$ be the far end of $(w,v)$ and $q_m\to q_e$ in the open edge with
$\alpha^{(w)}_{q_m}<\infty$. In world $v$'s frame $C_{q_m}$ has thresholds
$\alpha_m\coloneqq\alpha^{(v)}_{q_m}=1/\beta^{(w)}_{q_m}$ and
$\beta_m\coloneqq\beta^{(v)}_{q_m}=1/\alpha^{(w)}_{q_m}>0$. Since
$C_{q_e}=D_v$, \ref{P:ends}(a) in world $w$'s frame gives
$\beta^{(w)}_{q_m}\to\infty$, so $\alpha_m\to0$.

Let $M_T\coloneqq D_v+\sil(T)$. On $T\subseteq D_{vw}$, $D_v$ gives credit
at $v$, which is debt at $w$, while $M_T$ is silent; elsewhere the two
models agree. So $M_T\prec_wD_v=C_{q_e}$ by strict (G) at $w$, and (T2)
gives $C_q\preceq_vM_T$ for all $q$ in the open edge near $q_e$. Fix $m$
with $C_{q_m}\preceq_vM_T$ and $\alpha_m<c$, and put $b\coloneqq\beta_m$;
then $0<b\le\alpha_m<c$. For measurable $W\subseteq
D_{vw}\cap\{t^v_w<b\}$, $C_{q_m}$ gives debt at $v$ on $W$, so
$D_v+\debt_v(W)\le^v_\pw C_{q_m}$, and weak (G) gives
$D_v+\debt_v(W)\preceq_vC_{q_m}\preceq_vM_T$. Finally $W\cap
T=\varnothing$, since $b<c$.
\end{proof}

\subsection{Forced moves at an edge of a (W)-world}\label{sec:forced}

\begin{lemma}\label{lem:forced}
Let world $k$ be a (W)-world and $x\ne k$ a world with (IB). Assume (G) at
worlds $x$ and $k$, (S) and (EP$_{xk}$). Put $D\coloneqq D_{xk}$ and
$t\coloneqq t^x_k$, and use world $x$'s frame on $[e_x,e_k]$.
\begin{enumerate}[label=(\alph*),leftmargin=*,itemsep=2pt]
\item If $B$ is coherent, undominated and competitive on $\{x,k\}$, and $X$
is regular with $L_k(X)<L_k(B)$, then $B\preceq_xX$.
\item Let $0<s<\infty$, $Q\subseteq D\cap\{t<s\}$ and $W\subseteq
D\cap\{t\ge s\}$ be measurable with $\nu_{k1}(Q)<\nu_{k1}(W)$. Then
$Q\preceq^x_dW$.
\item Let $0<b<\infty$ and let $Z\subseteq D\cap\{t\ge b\}$ be non-null.
There is $\delta>0$ such that for every $s\in(b-\delta,b)$
\[
 V\not\preceq_xV+\debt_x(D\cap\{s\le t<b\})+\cred_x(Z).
\]
\end{enumerate}
Part (a) does not use (EP$_{xk}$), and part (c) does not use \eqref{eq:W}.
\end{lemma}
\begin{proof}
(a) By \eqref{eq:W} in contrapositive form, $B\not\preceq_kX$; by
competitiveness, $B\preceq_xX$.

(b) This is TAW's Lemma~8 (M1), with worlds $1,2$ replaced by $x,k$. By
\ref{P:param}, $E_k(s)$ is selected at a point of the open edge, so it is
coherent, undominated and competitive on $\{x,k\}$. It gives debt at $x$ on
$Q$ and is silent on $W$. Let $X\coloneqq E_k(s)+\sil(Q)+\debt_x(W)$. By
\eqref{eq:swap}, debt at $x$ is credit at $k$ on $D$; so $X$ changes the
world-$k$ rank from $0$ to $1$ on $Q$ and from $1$ to $0$ on $W$, and
$L_k(X)-L_k(E_k(s))=\nu_{k1}(Q)-\nu_{k1}(W)<0$
(Section~\ref{sec:losses}). By (a), $E_k(s)\preceq_xX$. The pairs
$(E_k(s),X)$ and $(V+\debt_x(Q),V+\debt_x(W))$ differ exactly on $Q\cup W$,
with the same ranks there; (IB$'$) gives $Q\preceq^x_dW$.

(c) This is TAW's Lemma~10(c). Let $Y\coloneqq E_k(b)+\cred_x(Z)$; $E_k(b)$
is silent on $Z$ (up to the null set $\{t=b\}$), so $Y\prec_xE_k(b)$ by
strict (G). Apply \ref{P:approach} with $s^*=b$ and $s_n\uparrow b$,
$s_n>0$; it gives $q^*$ with $C_{q^*}=E_k(b)$. Since $E_k(b)\not\preceq_xY$,
$q^*\notin A_x(Y)$, and as $A_x(Y)$ is closed some neighbourhood $N$ of
$q^*$ misses it. By \ref{P:approach} there is $n_0$ such that every
$s\in[s_{n_0},b)$ is $s(q)$ for some $q\in N$; so $E_k(s)\not\preceq_xY$.
For such $s$, the models $E_k(s)$ and $Y$ differ exactly on
$D\cap\{s\le t<b\}$ (silence and debt at $x$) and on $Z$ (silence and
credit), as do $V$ and $V+\debt_x(D\cap\{s\le t<b\})+\cred_x(Z)$. By
(IB$'$) the claim holds with $\delta\coloneqq b-s_{n_0}$.
\end{proof}

TAW's forced moves (M3) and (M5) hold in the same generality; they are not
needed below.

\subsection{The chain lemma and sweeps}\label{sec:sweep}

\begin{lemma}\label{lem:chain}
Let $\{u,v,w\}=\{1,2,3\}$ and assume (IB) and strict (G) at world $u$. Let
$V$ be competitive on $\{1,2,3\}$ and let $C$ be competitive on $\{u,v\}$.
Let $A\subseteq D_{uv}$ and $S\subseteq P_w$ be measurable, $S$ non-null,
$X'\coloneqq V+o_{uv}(S)$, $X$ the model $X'$ with $C$'s advice on $A$, and
$Y$ the model $C$ silenced on $A$. Then $V\not\preceq_vX$,
$V\not\preceq_wX$ and $C\not\preceq_vY$ do not all hold. If world $v$
satisfies (IB), then $C\not\preceq_vY$ is equivalent to $X\not\preceq_vX'$.
\end{lemma}
\begin{proof}
This is FSE's Lemma~37, with the non-comparison
$C\not\preceq_vY$ taken as a hypothesis; then no (IB) is needed at $v$.
Suppose all three hold. Competitiveness of $V$ gives $V\preceq_uX$, and
competitiveness of $C$ gives $C\preceq_uY$. Since $S\cap A=\varnothing$,
$X'$ is silent on $A$; so the pairs $(C,Y)$ and $(X,X')$ differ exactly on
the set where $C$ is not silent inside $A$, with $C$'s advice against
silence. By (IB$'$) at $u$, $X\preceq_uX'$. But $o_{uv}$ is good at $u$ on
$P_w$, so $X'\prec_uV$ by strict (G), and $V\preceq_uX\preceq_uX'\prec_uV$
is impossible. The last statement is (IB$'$) at $v$ applied to the same two
pairs.
\end{proof}

For a non-null measurable $S$ and a world $x$, let $\Nid^x_S$ be the family
of measurable sets contained in a finite union of measurable sets $B$ with
$S\not\preceq^x_cB$. If $B'\subseteq B$ and $S\preceq^x_cB'$, then
$S\preceq^x_cB$ by weak (G); and $S\not\preceq^x_c\varnothing$, since
$S\preceq^x_c\varnothing$ would say $V\preceq_xV+\cred_x(S)$, against strict
(G). So $\Nid^x_S$ is an ideal, and an ultrafilter $\Uf$ that avoids it
satisfies $S\preceq^x_cP$ for \emph{every} $P\in\Uf$. This is a positive
comparison, which is what the counting lemma needs; the ideals generated
by non-comparisons give only non-comparisons for their avoiders.

\begin{lemma}\label{lem:sweep}
Let world $x$ satisfy (G) and (IB). Let $S,T,Q\subseteq H'$ be pairwise
disjoint measurable sets, $S$ and $T$ non-null, such that the pair $(T,Q)$
is not weakly AC at $x$. Suppose that for every measurable
$\Sigma\subseteq Q$ the model
$X_\Sigma\coloneqq V+\cred_x(S)+\debt_x(T)+\cred_x(\Sigma)$ satisfies
\begin{equation}
 X_\Sigma\preceq_xV+\cred_x(S)\qquad\text{or}\qquad V\preceq_xX_\Sigma.
 \label{eq:cover}
\end{equation}
If $Q$ is the union of pairwise disjoint non-null measurable sets
$Q^1,\dots,Q^r$, then $S\preceq^x_cQ^i$ for some $i$. In particular
$Q\notin\Nid^x_S$, and some ultrafilter contains $Q$ and satisfies
$S\preceq^x_cP$ for all its members $P$.
\end{lemma}
\begin{proof}
$Q$ is non-null, since $(T,\varnothing)$ is weakly AC by weak (G). Fix a
measurable $\psi\colon H\to\R$ whose level sets are null, for instance
$\psi(g)\coloneqq g_1/|g|$, and let $F_i(y)\coloneqq\mu(Q^i\cap\{\psi\le
y\})/\mu(Q^i)$, a continuous distribution function. Then
$\mu(Q^i\cap\{F_i(\psi)\le\lambda\})=\lambda\mu(Q^i)$ for $\lambda\in[0,1]$.
For $\theta\in[i-1,i]$ put
\[
 \Sigma_\theta\coloneqq Q^1\cup\dots\cup Q^{i-1}\cup
 \bigl(Q^i\cap\{F_i(\psi)\le\theta-i+1\}\bigr);
\]
up to null sets this is consistent at integers, $\Sigma_0=\varnothing$,
$\Sigma_r=Q$, and for $i-1\le\theta<\theta'\le i$ the set
$\Sigma_{\theta'}\setminus\Sigma_\theta$ is a non-null subset of $Q^i$. Put
$X'\coloneqq V+\cred_x(S)$, $X_\theta\coloneqq X_{\Sigma_\theta}$,
\[
 \mathcal U\coloneqq\{\theta\in[0,r]:X_\theta\preceq_xX'\},\qquad
 \mathcal D\coloneqq\{\theta\in[0,r]:V\preceq_xX_\theta\}.
\]
If $\theta\le\theta'$, then $X_{\theta'}\le^x_\pw X_\theta$, so
$X_{\theta'}\preceq_xX_\theta$; hence $\mathcal U$ is an up-set and
$\mathcal D$ a down-set. They are disjoint, since
$V\preceq_xX_\theta\preceq_xX'$ would contradict $X'\prec_xV$ (strict (G),
$S$ non-null). $0\notin\mathcal U$, since $X_0=X'+\debt_x(T)$ and
$X'\prec_xX_0$ by strict (G). $r\notin\mathcal D$, since
$X_r\le^x_\pw V+\debt_x(T)+\cred_x(Q)$ and $(T,Q)$ is not weakly AC. By
\eqref{eq:cover}, $\mathcal U\cup\mathcal D=[0,r]$; so $0\in\mathcal D$ and
$r\in\mathcal U$. Let $\theta^*\coloneqq\sup\mathcal D$.

If $\theta^*\in\mathcal D$, then $\theta^*<r$ and every $\theta>\theta^*$
lies in $\mathcal U$. Choose $i$ with $\theta^*\in[i-1,i)$ and
$\theta\in(\theta^*,i]$, and put $\Delta\coloneqq\Sigma_\theta\setminus
\Sigma_{\theta^*}$, a non-null subset of $Q^i$, and
$N\coloneqq V+\debt_x(T)+\cred_x(\Sigma_{\theta^*})$, which is silent on
$S\cup\Delta$. The pairs $(X_\theta,X')$ and $(N+\cred_x(\Delta),V)$ differ
exactly on $T\cup\Sigma_\theta$, with the same advice there; so
$X_\theta\preceq_xX'$ gives $N+\cred_x(\Delta)\preceq_xV$ by (IB$'$).
Moreover $V\preceq_xX_{\theta^*}=N+\cred_x(S)$. Hence
$N+\cred_x(\Delta)\preceq_xN+\cred_x(S)$, and (IB$'$) in the context $N$
gives $V+\cred_x(\Delta)\preceq_xV+\cred_x(S)$, that is,
$S\preceq^x_c\Delta$. If $\theta^*\notin\mathcal D$, then
$\theta^*\in\mathcal U$, $\theta^*>0$, and every $\theta<\theta^*$ lies in
$\mathcal D$. Choose $i$ with $\theta^*\in(i-1,i]$ and
$\theta\in[i-1,\theta^*)$, and put $\Delta\coloneqq\Sigma_{\theta^*}
\setminus\Sigma_\theta\subseteq Q^i$ and $N\coloneqq V+\debt_x(T)+
\cred_x(\Sigma_\theta)$. The same computation, with
$X_{\theta^*}\preceq_xX'$ and $V\preceq_xX_\theta$, gives
$S\preceq^x_c\Delta$. In both cases $\Delta\subseteq Q^i$ gives
$\Delta\preceq^x_cQ^i$, so $S\preceq^x_cQ^i$.

If $Q\in\Nid^x_S$, then $Q\subseteq B^1\cup\dots\cup B^r$ up to a null set,
with $S\not\preceq^x_cB^j$. The non-null sets among
$(Q\cap B^j)\setminus(B^1\cup\dots\cup B^{j-1})$, with the null remainder of
$Q$ added to one of them, partition $Q$, and each lies in some $B^j$, so
$S$ is $\preceq^x_c$ none of them. This contradicts the first part. The
ultrafilter is given by \ref{U:avoid}.
\end{proof}

\section{One additive world}\label{sec:onew}

In this section world~$2$ is a (W)-world. The standing hypothesis is:
\begin{enumerate}[label=(OW),leftmargin=*]
\item $n=3$; (G) at every world; (S); world~$2$ is a preorder with a loss
$L_2$ of CBOI's form satisfying \eqref{eq:positive} and (W$_2$) for every
coherent undominated model; worlds $1$ and $3$ are preorders satisfying
(IB).
\end{enumerate}
Neither completeness nor (IB) is assumed at world~$2$, and completeness is
assumed at worlds $1$ and $3$ only where stated. For $x\in\{1,3\}$, by
\eqref{eq:swap}, changing silence to debt at $x$ on a set $Q\subseteq
D_{x2}$ lowers $L_2$ by $\nu_{21}(Q)$, and changing silence to credit at $x$
raises it by $\nu_{22}(Q)$ (Section~\ref{sec:losses}).

\subsection{The endpoint property at the (W)-world}\label{sec:EP}

\begin{lemma}\label{lem:nwAC}
Assume (OW) and let $x\in\{1,3\}$; write $t\coloneqq t^x_2$ on $D_{x2}$.
\begin{enumerate}[label=(\alph*),leftmargin=*,itemsep=2pt]
\item If (EP$_{x2}$) holds, then for every non-null $T\subseteq
P_2\cap\{t\ge1\}$ there is a non-null $S\subseteq P_x$ such that the pair
$(S,T)$ is not weakly AC at world $x$.
\item If $C'$ is a two-sided selection on $(e_x,e_2)$, with thresholds
$(\alpha',\beta')$ in world $x$'s frame, then for every non-null
$T\subseteq P_2\cap\{t>\alpha'\}$ there is a non-null $S\subseteq
P_x\cap\{t<\beta'\}$ such that $(S,T)$ is SI at world $x$.
\end{enumerate}
No completeness is used, and $V$ is not assumed to be undominated.
\end{lemma}
\begin{proof}
(b) $C'$ is selected at a point of the open edge, so it is coherent,
undominated and competitive on $\{x,2\}$. It gives credit at $x$ on $T$ and
debt at $x$ on every subset of $P_x\cap\{t<\beta'\}$. Since
$\nu_{22}(T)>0$ by \eqref{eq:positive}, there is a non-null $Q\subseteq
P_x\cap\{t<\beta'\}$ with $\nu_{21}(Q)<\nu_{22}(T)$; choose a non-null
$S\subseteq Q$ with $Q\setminus S$ non-null. Let $X\coloneqq
C'+\sil(Q\cup T)$. Then $L_2(X)-L_2(C')=\nu_{21}(Q)-\nu_{22}(T)<0$, so
Lemma~\ref{lem:forced}(a) gives $C'\preceq_xX$. As
$C'=X+\debt_x(Q)+\cred_x(T)$ and $X$ is silent on $Q\cup T$, (IB$'$) gives
$V+\debt_x(Q)+\cred_x(T)\preceq_xV$. Removing the debt on the non-null
$Q\setminus S$ is a strict improvement by strict (G), so
$V+\debt_x(S)+\cred_x(T)\prec_xV$.

(a) Lemma~\ref{lem:forced}(c), with $k=2$, $b=1$ and $Z=T$, gives
$\delta>0$. Fix $s\in(\max(0,1-\delta),1)$ and put $\Xi\coloneqq
D_{x2}\cap\{s\le t<1\}$, so that
$V\not\preceq_xV+\debt_x(\Xi)+\cred_x(T)$. The sets $W\coloneqq\Xi\cap P_2$
and $S_1\coloneqq\Xi\cap P_x$ are non-null. Choose a non-null $Q\subseteq
P_x\cap\{t<s\}$ with $\nu_{21}(Q)<\nu_{21}(W)$. Lemma~\ref{lem:forced}(b),
with threshold $s$, gives $Q\preceq^x_dW$. Put $S\coloneqq Q\cup S_1\subseteq
P_x$. Then $S\setminus\Xi=Q$ and $\Xi\setminus S=W$, so $S\preceq^x_d\Xi$.
The context $V+\cred_x(T)$ is silent on $S\cup\Xi$, since $T\subseteq
P_2\cap\{t\ge1\}$; so (IB$'$) gives
$V+\debt_x(S)+\cred_x(T)\preceq_xV+\debt_x(\Xi)+\cred_x(T)$. If $(S,T)$ were
weakly AC, transitivity would give $V\preceq_xV+\debt_x(\Xi)+\cred_x(T)$,
which is false.
\end{proof}

\begin{theorem}\label{thm:EP}
Assume (OW) and that $V$ is undominated. Then (EP$_{12}$) and (EP$_{32}$)
hold: every selection on $[e_1,e_2]$ or on $[e_3,e_2]$ has an endpoint of
its edge in its credal set.
\end{theorem}
\begin{proof}
Let $\{x,u\}=\{1,3\}$; we prove (EP$_{x2}$). $V$ is competitive on
$\{1,2,3\}$. Suppose a selection $C$ on $[e_x,e_2]$ has neither $e_x$ nor
$e_2$ in its credal set. Then $C$ is two-sided and selected at a point of
the open edge, so it is coherent, undominated and competitive on $\{x,2\}$.
Write $C=(\alpha,\beta)$ in world $x$'s frame, $0<\beta\le\alpha<\infty$,
$t\coloneqq t^x_2$, and
\[
 T_0\coloneqq P_2\cap\{t>\alpha\},\qquad B_0\coloneqq P_x\cap\{t<\beta\}.
\]
$C$ gives credit at $x$ on $T_0$, which on $P_2$ is the option $o_{xu}$
(credit at $x$ and $u$, debt at $2$), and debt at $x$ on $B_0$, which on
$P_x$ is $o_{u2}$ (debt at $x$, credit at $u$ and $2$).

\emph{Step 0.} If (EP$_{u2}$) holds, put $K\coloneqq P_2\cap\{t^u_2\ge1\}$.
Otherwise some selection on $[e_u,e_2]$ omits both endpoints; it is
two-sided, with thresholds $(\alpha',\beta')$ in world $u$'s frame, and we
put $K\coloneqq P_2\cap\{t^u_2>\alpha'\}$. By Lemma~\ref{lem:nwAC} at
world $u$, for every non-null $T\subseteq K$ there is a non-null
$S\subseteq P_u$ such that $(S,T)$ is not weakly AC at $u$. The set
$T_0\cap K$ is non-null (two ratios on $P_2$).

\emph{Step 1.} Let $c_{21}$ and $Z_m$ be the density and the sets of
Section~\ref{sec:losses} for $\nu_{21}$. Fix $m$ with $\mu(Z_m)<\mu(B_0)$
and put $B'_0\coloneqq B_0\setminus Z_m$, a non-null set on whose measurable
subsets $\nu_{21}=\int c_{21}\,d\mu$. Choose a non-null $T\subseteq T_0\cap
K$ with $\nu_{22}(T)<\nu_{21}(B'_0)$.

\emph{Step 2.} By Step~0 there is a non-null $S\subseteq P_u$ with
$V\not\preceq_uV+\debt_u(S)+\cred_u(T)$. By \eqref{eq:positive},
$\nu_{21}(S)>0$.

\emph{Step 3.} Let $\psi(g)\coloneqq g_1/|g|$, whose level sets are null.
The map $y\mapsto\nu_{21}(B'_0\cap\{\psi\le y\})=\int_{B'_0}
c_{21}\mathbf 1\{\psi\le y\}\,d\mu$ is continuous by dominated convergence,
and increases from $0$ to $\nu_{21}(B'_0)>\nu_{22}(T)$ as $y$ runs through
$[0,1]$. So there is $y$ with
$0<\nu_{22}(T)-\nu_{21}(\Sigma)<\nu_{21}(S)$ for
$\Sigma\coloneqq B'_0\cap\{\psi\le y\}$.

\emph{Step 4.} Apply Lemma~\ref{lem:chain} with $(u,v,w)$ there equal to
$(x,2,u)$, with $C$, $A\coloneqq T\cup\Sigma\subseteq D_{x2}$ and
$S\subseteq P_u$. Then $X'=V+o_{x2}(S)$,
$X=V+o_{x2}(S)+o_{xu}(T)+o_{u2}(\Sigma)$, and $Y$ is $C$ silenced on $A$.
On $P_u$ the option $o_{x2}$ is credit at $2$ and debt at $u$. So
\[
 L_2(X)-L_2(V)=-\nu_{21}(S)+\nu_{22}(T)-\nu_{21}(\Sigma)<0,\qquad
 L_2(Y)-L_2(C)=-\nu_{22}(T)+\nu_{21}(\Sigma)<0.
\]
By (W$_2$) at the coherent undominated models $V$ and $C$, in contrapositive
form, $V\not\preceq_2X$ and $C\not\preceq_2Y$. At world $u$, $o_{x2}$ is
debt on $P_u$, $o_{xu}$ is credit on $P_2$ and $o_{u2}$ is credit on $P_x$;
so $X\le^u_\pw V+\debt_u(S)+\cred_u(T)$, and $V\preceq_uX$ would give
$V\preceq_uV+\debt_u(S)+\cred_u(T)$, contrary to Step~2. So
$V\not\preceq_uX$. World $x$ satisfies (IB) and strict (G), $V$ is
competitive on $\{1,2,3\}$ and $C$ on $\{x,2\}$; Lemma~\ref{lem:chain} gives
a contradiction.
\end{proof}

\begin{remark}\label{rem:EP}
(a) The proof uses (IB) at worlds $1$ and $3$, (G) at every world,
transitivity, closedness of sections, and (W$_2$) only in contrapositive
form, at selected models ($C'$ in Lemma~\ref{lem:nwAC}(b), $E_2(s)$ in
world $u$'s frame inside Lemma~\ref{lem:forced}(b), and $C$) and at $V$. It
uses no completeness at any world and no (IB) at world~$2$. Every passage
from a non-comparison to a comparison is competitiveness of a selected
model or of $V$.

(b) The use of (W$_2$) at $V$ is essential. $V$ is tied at every world with
a selected model (CBOI, Lemma~1), but (W$_2$) at that selected model gives
only an upper bound on its loss, not $V\not\preceq_2X$. Nor would a bound
local to the edges of world~$2$ do: $V$ and $X$ in Step~4 differ on
$D_{x2}\cup P_u$.

(c) CCF's Proposition~17 obtains the endpoint property on the two edges of
a world from (W)-bounds at the \emph{other} two worlds. Theorem~\ref{thm:EP}
obtains it on the two edges of the single (W)-world.
\end{remark}

\subsection{One (W)-world and one complete world}\label{sec:onecomplete}

On $D_{31}=P_3\cup P_1$ put $\sigma\coloneqq t^3_1$, and on
$D_{32}=P_3\cup P_2$ put $\tau\coloneqq t^3_2$; on $P_3$ the pair
$(\sigma,\tau)$ has a positive joint density. In world~$3$'s frame on
$[e_3,e_2]$ write $F(b)\coloneqq E_2(b)$: debt at $3$ on
$D_{32}\cap\{\tau<b\}$, silence on the rest of $D_{32}$.

\begin{lemma}\label{lem:SI3}
Assume (G) at worlds $2$ and $3$, (S), (IB) at world~$3$, (EP$_{32}$), and
that world~$3$ is complete. For $0<b<\infty$ and non-null $Z\subseteq
D_{32}\cap\{\tau>b\}$ there is $s\in(0,b)$ such that
$(D_{32}\cap\{s\le\tau<b\},Z)$ is SI at world~$3$.
\end{lemma}
\begin{proof}
This is the SI half of CCF's Lemma~11, at the edge $[e_3,e_2]$. Let
$Y\coloneqq F(b)+\cred_3(Z)$; $F(b)$ is silent on $Z$, so $Y\prec_3F(b)$ by
strict (G). By \ref{P:approach} with $s^*=b$ and $s_n\uparrow b$, $s_n>0$,
there is $q^*$ with $C_{q^*}=F(b)$ such that every neighbourhood of $q^*$
contains points $q$ with $C_q=F(s)$ for all $s$ in some $[s_{n_0},b)$. By
(T1), which uses the completeness of world~$3$, $Y\prec_3C_q$ for $q$ near
$q^*$; so $Y\prec_3F(s)$ for some $s\in(0,b)$. Since
$Y=F(s)+\debt_3(D_{32}\cap\{s\le\tau<b\})+\cred_3(Z)$, (IB$'$) gives the
claim.
\end{proof}

\begin{lemma}\label{lem:fulllight}
Assume (OW), (EP$_{32}$), and that world~$3$ is complete. For every non-null
$T_0\subseteq H'$ there are $0<c_0<a$ such that
$T\coloneqq P_2\cap\{c_0<\tau<a\}$ satisfies $T\preceq^3_cT_0$, and such
that for every non-null $T'\subseteq T$ there are $\varepsilon>0$ and $\delta>0$ for
which every measurable $W\subseteq P_3\cap\{\tau<\varepsilon\}$ with
$\nu_{21}(W)<\delta$ is $T'$-light at world~$3$.
\end{lemma}
\begin{proof}
By \ref{R:yardc} at world~$3$ there is $a>0$ such that every measurable
subset of $\Lo32a$ is $\preceq^3_cT_0$. Put $c_0\coloneqq a/2$; then
$T\subseteq\Lo32a$ is non-null and $T\preceq^3_cT_0$. Let $T'\subseteq T$
be non-null and split it into disjoint non-null $T'_a,T'_b$.
Lemma~\ref{lem:SI3} with $b\coloneqq c_0$ and $Z\coloneqq T'_a$ gives
$s\in(0,c_0)$ such that $(\Xi,T'_a)$ is SI at~$3$, where
$\Xi\coloneqq D_{32}\cap\{s\le\tau<c_0\}$. By \ref{R:E1} at world~$3$, with
$c\coloneqq c_0$ and $Q\coloneqq T'_b$, there is $a'>0$ with
$\Lo32{a'}\preceq^3_cT'_b$. Put $\varepsilon\coloneqq\min(s,a')$ and
$\delta\coloneqq\nu_{21}(\Xi)>0$.

Let $W\subseteq P_3\cap\{\tau<\varepsilon\}$ with $\nu_{21}(W)<\delta$. Since
$\Xi\subseteq\{\tau\ge \varepsilon\}$, Lemma~\ref{lem:forced}(b) at $(x,k)=(3,2)$, with
threshold $\varepsilon$, gives $W\preceq^3_d\Xi$. The context $V+\cred_3(T'_a)$ is
silent on $W\cup\Xi$, so
$V+\debt_3(W)+\cred_3(T'_a)\preceq_3V+\debt_3(\Xi)+\cred_3(T'_a)\prec_3V$,
and $(W,T'_a)$ is SI, in particular weakly SI. As $W\subseteq\Lo32{a'}$,
$W\preceq^3_cT'_b$. The sets $W\subseteq P_3$ and $T'_a,T'_b\subseteq P_2$
are pairwise disjoint, and Lemma~\ref{lem:light}(d) shows that $W$ is
$T'$-light at~$3$.
\end{proof}

For non-null $T_0\subseteq H'$ put
\[
 \Hset(T_0)\coloneqq\{\Uf:\ D_{31}\in\Uf\text{ and }T_0\prec^3_cB\text{ for
 every }B\in\Uf\}.
\]

\begin{proposition}\label{prop:count3}
Assume (OW), (EP$_{32}$), and that world~$3$ is complete. Then
$\Hset(T_0)$ is finite for every non-null $T_0$, and
$\Hset\coloneqq\bigcup_{T_0}\Hset(T_0)$ is countable.
\end{proposition}
\begin{proof}
Fix $T_0$, and let $T$, $c_0$, $a$ be as in Lemma~\ref{lem:fulllight}.
Split $T$ into pairwise disjoint non-null sets $T_n$, $n\ge1$, and let
$\varepsilon_n,\delta_n$ be given by Lemma~\ref{lem:fulllight} for $T'\coloneqq T_n$.
Let $(b_n)_{n\ge1}$ enumerate the bands $D_{31}\cap\{r<\sigma<r'\}$ with $r$
rational, $r'$ rational or $\infty$, and $0\le r<r'$. Choose $0<\varepsilon'_n<\varepsilon_n$
with $\varepsilon'_n<\varepsilon'_{n-1}/2$ for $n\ge2$. The set $P_3\cap
b_n\cap\{\varepsilon'_n/2<\tau<\varepsilon'_n\}$ is non-null; choose a non-null $W_n$ in it with
$\nu_{21}(W_n)<\delta_n$. The $W_n$ are pairwise disjoint (disjoint ranges
of $\tau$), they lie in $P_3$ and so miss $T\subseteq P_2$, and $W_n$ is
$T_n$-light at~$3$ by Lemma~\ref{lem:fulllight}. The family $\Bfam$ of
measurable subsets of the sets $W_n$ is band-dense for $[e_3,e_1]$, since
every band contains some $b_n\supseteq W_n$.

By \ref{R:cut} with $(v,w)=(3,1)$, which uses the completeness of
world~$3$, there is a connecting chain from $D_3$ to $D_1$ with perturbation
sets $B_1,\dots,B_{\ell-1}\in\Bfam$, and $\ell\ge2$. Choose $N$ such that
every $B_i$ lies in some $W_n$ with $n\le N$, and put $F\coloneqq
W_1\cup\dots\cup W_N$ and $T^*\coloneqq T_1\cup\dots\cup T_N$. By
Lemma~\ref{lem:light}(a), $F$ is $T^*$-light at~$3$, and so is every
$B_i\subseteq F$. By Lemma~\ref{lem:light}(b), $F\preceq^3_cT^*$; and
$T^*\preceq^3_cT\preceq^3_cT_0$.

Suppose $\Hset(T_0)$ has $\ell$ distinct members $\Uf^1,\dots,\Uf^\ell$. No
$\Uf^l$ contains $F$, since otherwise
$T_0\prec^3_cF\preceq^3_cT^*\preceq^3_cT_0$; so $D_{31}\setminus F\in\Uf^l$.
For $l<l'$ choose $G_{ll'}\in\Uf^l$ with $H'\setminus G_{ll'}\in\Uf^{l'}$,
and put $G_{l'l}\coloneqq H'\setminus G_{ll'}$. The sets
$K_l\coloneqq(D_{31}\setminus F)\cap\bigcap_{l'\ne l}G_{ll'}\in\Uf^l$ are
non-null and pairwise disjoint, and $T^*\preceq^3_cT_0\prec^3_cK_l$. For
every $i$ and $l$ the sets $B_i\subseteq P_3$, $T^*\subseteq P_2$ and
$K_l\subseteq D_{31}\setminus F$ are pairwise disjoint, so
Lemma~\ref{lem:light}(c) gives
$D_3+\debt_3(B_i)\prec_3D_3+\debt_3(K_l)$. Since $K_l\subseteq
D_{31}\setminus F$, this contradicts the counting lemma \ref{R:count} with
$(v,w)=(3,1)$. Hence $|\Hset(T_0)|\le\ell-1$.

Countability: by \ref{R:yardc} at world~$3$ there is $a_0>0$ such that every
measurable subset of $\Lo32{a_0}$ is $\preceq^3_cT_0$. For integers
$k>1/a_0$ put $L_k\coloneqq\Lo32{1/k}$; then $L_k\preceq^3_cT_0$, so
$T_0\prec^3_cB$ implies $L_k\prec^3_cB$, and
$\Hset(T_0)\subseteq\Hset(L_k)$. Hence $\Hset=\bigcup_k\Hset(L_k)$ is a
countable union of finite sets.
\end{proof}

\begin{proposition}\label{prop:heavy3}
Assume (OW), (EP$_{32}$), and that $V$ is undominated. Let $C$ be a
two-sided selection on $(e_1,e_3)$, with thresholds $(\alpha,\beta)$ in
world~$3$'s frame, and put
\[
 T^C_0\coloneqq P_1\cap\{\sigma>\alpha\},\qquad
 B^C_0\coloneqq P_3\cap\{\sigma<\beta\}.
\]
For every non-null $Q\subseteq T^C_0$ there are non-null $Q'\subseteq Q$ and
$T\subseteq B^C_0$ such that $Q'\notin\Kid^3_{S'}$ for every non-null
$S'\subseteq P_2$ with $\nu_{22}(S')<\nu_{21}(T)$. No completeness is used.
\end{proposition}
\begin{proof}
$C$ gives credit at $3$ on $T^C_0$, which on $P_1$ is $o_{23}$, and debt at
$3$ on $B^C_0$, which on $P_3$ is $o_{12}$. $C$ is competitive on
$\{1,3\}$.

\emph{Step 1: a pair with debt part in $B^C_0$ and credit part in $Q$ that
is not weakly AC at~$3$.} Suppose every pair $(A,T')$ of world~$3$ with
$A\subseteq D_{32}$ non-null and $T'\subseteq Q$ is weakly AC. By
\ref{R:yardc} at world~$3$ there is $a>0$ such that every measurable subset
of $\Lo32a$ is $\preceq^3_cQ$. Then every pair $(A,T'')$ with $A\subseteq
D_{32}$ non-null and $T''\subseteq\Lo32a$ is weakly AC: the context
$V+\debt_3(A)$ is silent on $Q\cup T''$, so
$V\preceq_3V+\debt_3(A)+\cred_3(Q)\preceq_3V+\debt_3(A)+\cred_3(T'')$. But
Lemma~\ref{lem:forced}(c) at $(x,k)=(3,2)$, with $b\coloneqq a/2$ and a
non-null $Z\subseteq D_{32}\cap\{a/2<\tau<a\}$, shows that for $s<a/2$
close to $a/2$ the pair $(D_{32}\cap\{s\le\tau<a/2\},Z)$, whose credit
part lies in $\Lo32a$, is not weakly AC. Hence there are non-null $A\subseteq D_{32}$ and
$Q'\subseteq Q$ such that $(A,Q')$ is not weakly AC; $Q'$ is non-null since
$(A,\varnothing)$ is weakly AC. Choose $s_0>0$ with $A^\circ\coloneqq
A\cap\{\tau\ge s_0\}$ non-null. By weak (G),
$V+\debt_3(A^\circ)+\cred_3(Q')\preceq_3V+\debt_3(A)+\cred_3(Q')$, so
$(A^\circ,Q')$ is not weakly AC either. The set $B^C_0\cap\{\tau<s_0\}$ is
non-null; choose a non-null $T$ in it with $\nu_{21}(T)<\nu_{21}(A^\circ)$
and $\nu_{21}(T)<\nu_{22}(P_2)$. Lemma~\ref{lem:forced}(b), with threshold
$s_0$, gives $T\preceq^3_dA^\circ$. The context $V+\cred_3(Q')$ is silent on
$T\cup A^\circ$, so if $(T,Q')$ were weakly AC, then
$V\preceq_3V+\debt_3(T)+\cred_3(Q')\preceq_3V+\debt_3(A^\circ)+\cred_3(Q')$.
Hence $(T,Q')$ is not weakly AC at~$3$.

\emph{Step 2: the sweep.} Let $S'\subseteq P_2$ be non-null with
$\nu_{22}(S')<\nu_{21}(T)$. Since $\nu_{21}(T)<\nu_{22}(P_2)$, the set
$P_2\setminus S'$ has positive $\nu_{22}$-mass and is non-null; choose a
non-null $S''\subseteq P_2\setminus S'$ with
$\nu_{22}(S'')<\nu_{21}(T)-\nu_{22}(S')$, and put $S\coloneqq S'\cup S''$.
Then $\nu_{22}(S)<\nu_{21}(T)$, and $S'\prec^3_cS$ by strict (G).

We verify the hypotheses of Lemma~\ref{lem:sweep} at $x=3$ for $S$, $T$ and
$Q'$. The sets $S\subseteq P_2$, $T\subseteq P_3$ and $Q'\subseteq P_1$ are
pairwise disjoint, $S$ and $T$ are non-null, and $(T,Q')$ is not weakly AC.
Let $\Sigma\subseteq Q'$ be measurable. In terms of options,
$X_\Sigma=V+o_{13}(S)+o_{12}(T)+o_{23}(\Sigma)$: on $P_2$, $o_{13}$ is
credit at $3$; on $P_3$, $o_{12}$ is debt at $3$; on $P_1$, $o_{23}$ is
credit at $3$. Suppose \eqref{eq:cover} fails for $\Sigma$:
$X_\Sigma\not\preceq_3X'$ and $V\not\preceq_3X_\Sigma$, where
$X'\coloneqq V+o_{13}(S)$. $X_\Sigma$ is $X'$ with $C$'s advice on
$A\coloneqq T\cup\Sigma\subseteq D_{13}$. At world~$2$, $o_{13}$ on $P_2$ is
debt, and $o_{12}$ on $P_3$ and $o_{23}$ on $P_1$ are credit; so
$L_2(X_\Sigma)-L_2(V)=\nu_{22}(S)-\nu_{21}(T)-\nu_{21}(\Sigma)<0$, and
(W$_2$) at the coherent undominated $V$ gives $V\not\preceq_2X_\Sigma$.
Apply Lemma~\ref{lem:chain} with $(u,v,w)=(1,3,2)$: $o_{13}$ is good at
$1$ on $P_2$, world~$1$ satisfies (IB) and strict (G), $V$ is competitive on
$\{1,2,3\}$, $C$ on $\{1,3\}$, and world~$3$ satisfies (IB), so
$X_\Sigma\not\preceq_3X'$ gives $C\not\preceq_3Y$ for $Y\coloneqq C$
silenced on $A$. With $V\not\preceq_3X_\Sigma$ and $V\not\preceq_2X_\Sigma$
this is a contradiction. So \eqref{eq:cover} holds for every $\Sigma$, and
Lemma~\ref{lem:sweep} gives $Q'\notin\Nid^3_S$.

\emph{Step 3.} If $Q'\in\Kid^3_{S'}$, then $Q'$ lies in a finite union of
measurable sets $B$ with $B\preceq^3_cS'$, and each such $B$ has
$S\not\preceq^3_cB$, since otherwise $S\preceq^3_cB\preceq^3_cS'\prec^3_cS$.
So $Q'\in\Nid^3_S$, contrary to Step~2.
\end{proof}

Proposition~\ref{prop:heavy3} is the mirror, under the exchange of worlds
$1$ and $3$, of the corresponding statement at world~$1$; it says that a
two-sided selection on the edge between the two worlds without (W)-bounds
forces world~$3$ to be heavy in credit, densely in $T^C_0\subseteq P_1$,
relative to all credit sets in $P_2$ of small $\nu_{22}$-mass. The choice
$\nu_{21}(T)<\nu_{22}(P_2)$ in Step~1 ensures that $P_2\setminus S'$ is
non-null in Step~2.

\begin{theorem}\label{thm:onecomplete}
Assume (OW) and let world~$3$ be complete. Then $V$ is dominated. The same
holds with worlds $1$ and $3$ exchanged.
\end{theorem}
\begin{proof}
Suppose $V$ is undominated. By Theorem~\ref{thm:EP}, (EP$_{12}$) and
(EP$_{32}$) hold.

\emph{Case I: (EP$_{31}$) holds.} World~$3$ is complete with (IB) and has
(EP) on both of its edges, so by \ref{R:CCF} it is a (W)-world. Worlds $2$
and $3$ are (W)-worlds and world~$1$ is a preorder with (IB), so $V$ is
dominated by \ref{R:TAW}: a contradiction.

\emph{Case II: some selection $C$ on $[e_3,e_1]$ omits both endpoints.}
Then $C$ is two-sided; with the notation of
Proposition~\ref{prop:heavy3}, $T^C_0$ is non-null. By
Proposition~\ref{prop:count3}, $\Hset$ is countable; list its members as
$\Uf^1,\Uf^2,\dots$ (the list may be finite or empty). By \ref{U:small}
choose $B^i\in\Uf^i$ with $\mu(B^i)<2^{-i-1}\mu(T^C_0)$, and put
$Q\coloneqq T^C_0\setminus\bigcup_iB^i$, a non-null set.
Proposition~\ref{prop:heavy3} gives non-null $Q'\subseteq Q$ and
$T\subseteq B^C_0$; choose a non-null $S'_0\subseteq P_2$ with
$\nu_{22}(S'_0)<\nu_{21}(T)$. Then $Q'\notin\Kid^3_{S'_0}$, and by
\ref{U:avoid} some ultrafilter $\Uf$ contains $Q'$ and avoids
$\Kid^3_{S'_0}$. By \ref{U:strict}, which uses the completeness of
$\preceq^3_c$, $S'_0\prec^3_cB$ for every $B\in\Uf$; and $D_{31}\in\Uf$
since $Q'\subseteq P_1$. So $\Uf\in\Hset(S'_0)\subseteq\Hset$, and
$\Uf=\Uf^i$ for some $i$. Then $\Uf$ contains the disjoint sets $Q'$ and
$B^i$, which is impossible.

(OW) is invariant under the exchange of the labels $1$ and $3$, which gives
the last statement.
\end{proof}

The completeness of world~$3$ is used in exactly four places: \ref{R:CCF} in
Case~I, (T1) in Lemma~\ref{lem:SI3} (inside Lemma~\ref{lem:fulllight}),
\ref{R:cut} in Proposition~\ref{prop:count3}, and \ref{U:strict} in
Case~II. Nothing is used about completeness at worlds $1$ and $2$, and
nothing about (IB) at world~$2$. That $V$ is undominated enters through
Theorem~\ref{thm:EP}, through (W$_2$) at $V$ in
Proposition~\ref{prop:heavy3}, and through the competitiveness of $V$ in
Lemma~\ref{lem:chain}.

\begin{corollary}\label{cor:Wforces}
Let $n=3$, assume (G) at every world and (S), let every world be a preorder
satisfying (IB), and suppose $V$ is undominated. If some world is a
(W)-world, then the other two worlds are incomplete, and neither of them is
a (W)-world.
\end{corollary}
\begin{proof}
Relabel so that world~$2$ is the (W)-world; then (OW) holds. If world~$3$ or
world~$1$ were complete, Theorem~\ref{thm:onecomplete} would make $V$
dominated. The second statement is \ref{R:TAW}.
\end{proof}

\subsection{Simple completeness}\label{sec:SC}

\begin{definition}\label{def:SC}
World $v$ satisfies \emph{(SC)} if (i) $\preceq^v_d$ is complete on the
measurable subsets of $H'$, (ii) $\preceq^v_c$ is complete, and (iii) every
pair of world $v$ is SI or weakly AC.
\end{definition}

A complete world satisfies (SC). (SC) concerns only comparisons of $V$ with
models that differ from $V$ by debt and credit at $v$, and, through (IB$'$),
comparisons in any context silent on the sets concerned.

\begin{theorem}\label{thm:SCext}
Assume (G) at every world and (S), and let world $v$ be a preorder with
(IB) and (SC) that has (EP) on both of its edges. Then world $v$ is fine,
and it has a loss $L_v$ of CBOI's form with \eqref{eq:positive} such that
$D\preceq_vE$ implies $L_v(D)\le L_v(E)$ for all regular $D,E$; in
particular it is a (W)-world. Neither completeness nor \To{} is assumed.
\end{theorem}
\begin{proof}
Relabel so that $v=1$. We follow CCF's proofs of its Theorems~7 and~15,
under the hypotheses of CCF's Theorem~7 with completeness of world~$1$
replaced by (SC). Completeness of world~$1$ enters those proofs in
three ways (the three kinds of use are those of CCF, Remark~21; the places are as follows): through (T1), in CCF's Lemmas~5(a), 5(c), 6(a),
6(c) and in the SI half of Lemma~11; in the comparative-probability step
(Lemma~8 and Theorem~9, together with Lemma~10(i), which uses the
completeness of $\preceq_d$ and $\preceq_c$ directly); and in reading a pair
that is not SI (CCF calls it AC) as weakly AC (CCF, Section~6). The rest uses only (G), (IB), transitivity,
(T2), (S) and CCF's Lemmas~3 and~4. We replace the three uses; CCF's
notation $\Band_j[s,s')=D_{1j}\cap\{s\le t^1_j<s'\}$ and similar is used.

\emph{(T1).} In each of the five places, CCF applies its Lemma~4 to an
approach $M_s\to M_{s^*}$ along an edge $[e_1,e_j]$, uses the closedness of
$A_1(Y)$ to get $M_s\not\preceq_1Y$ for the $s$ concerned, and then uses
completeness to get $Y\prec_1M_s$, which it converts by (IB) into a
statement about two sets. In each place $M_s$ and $Y$ differ on exactly two
sets, and by (IB$'$) the non-comparison $M_s\not\preceq_1Y$ already says:
\begin{itemize}[leftmargin=*,itemsep=1pt]
\item Lemma~5(a) ($Y=E_j(b)$ with its debt on $A$ silenced, $M_s=E_j(s)$)
and Lemma~5(c) ($Y=D_j$ with its debt on $A$ silenced, $M_s=E_j(s)$): the
models differ on $A$ (debt in $M_s$, silence in $Y$) and on a band $\Xi$
(silence in $M_s$, debt in $Y$); so $A\not\preceq_d\Xi$, and (SC)(i) gives
$\Xi\prec_dA$, which is CCF's conclusion.
\item Lemma~6(a) ($Y=E^j_1(a')+\cred(T)$, $M_s=E^j_1(s)$) and Lemma~6(c)
($Y=V_{1j}+\cred(T)$, $M_s=E^j_1(s)$): the models differ on a band $\Xi$
(credit in $M_s$, silence in $Y$) and on $T$ (silence in $M_s$, credit in
$Y$); so $V+\cred(\Xi)\not\preceq_1V+\cred(T)$, that is, $T\not\preceq_c\Xi$,
and (SC)(ii) gives $\Xi\prec_cT$, which is CCF's conclusion.
\item Lemma~11, SI half ($Y=E_j(1)+\cred(Z)$, $M_s=E_j(s)$, $s<1$): the
models differ on $\Band_j[s,1)$ (silence in $M_s$, debt in $Y$) and on $Z$
(silence in $M_s$, credit in $Y$); so the pair $(\Band_j[s,1),Z)$ is not
weakly AC, and by (SC)(iii) it is SI, which is CCF's conclusion.
\end{itemize}
\emph{Comparative probability.} CCF's Lemma~8 takes completeness of
$\preceq_d$ and $\preceq_c$ from that of $\preceq_1$; here it is (SC)(i),
(ii). So Theorem~9 (Savage and Wakker \cite{Savage,Wakker}) applies, and
gives $P_d,P_c$. In Lemma~10(i), $P_d(B)<P_d(C)$ gives $C\not\preceq_dB$ by
almost agreement, and (SC)(i) gives $B\prec_dC$; likewise for $P_c$ with
(SC)(ii). Lemma~10(ii)--(iv) use no completeness.

\emph{Not SI means weakly AC.} This is (SC)(iii). It is used in Lemma~12:
there, if all the pieces $(A_{0l},T_{0l})$ of an SI pair were not SI, they
would be weakly AC by (SC)(iii), and pairwise disjoint weakly AC pairs
compose to a weakly AC pair by (IB$'$) and transitivity, which an SI pair is
not. It is used in Lemma~13(b), where a pair that is not SI is read as
weakly AC. The AC pair of Lemma~11 is obtained from (T2) directly as
weakly AC. Theorem~14 uses only Lemmas~10 and~13, (G) and (IB), and
Theorem~15 assembles the loss from $P_d$, $P_c$ and the rate $\kappa$.

With these replacements CCF's Theorem~7 (fineness) and Theorem~15 hold for
world~$1$. The proofs are written in CCF for any $n\ge3$.
\end{proof}

\begin{theorem}\label{thm:SC}
Theorem~\ref{thm:onecomplete} remains true if world~$3$ is assumed, instead
of complete, to satisfy (IB), (SC) and \Tw{} along $[e_3,e_1]$.
\end{theorem}
\begin{proof}
We replace the four uses of completeness listed after
Theorem~\ref{thm:onecomplete}. Case~I: Theorem~\ref{thm:SCext} in place of
\ref{R:CCF}. Lemma~\ref{lem:SI3}: closedness of $A_3(Y)$ and \ref{P:approach}
give $F(s)\not\preceq_3Y$ for $s$ close to $b$, that is, by (IB$'$), the
pair $(D_{32}\cap\{s\le\tau<b\},Z)$ is not weakly AC; by (SC)(iii) it is
SI. Proposition~\ref{prop:count3}: Theorem~\ref{thm:cut} in place of
\ref{R:cut}; it needs \Tw{} at world~$3$ along $[e_3,e_1]$. Case~II:
\ref{U:strict} needs only the completeness of $\preceq^3_c$, which is
(SC)(ii). The remaining steps, including the countability argument in
Proposition~\ref{prop:count3}, use no completeness.
\end{proof}

Consequently, in Corollary~\ref{cor:Wforces}, neither of the two worlds
other than the (W)-world satisfies (SC) together with \Tw{} along the edge
that joins it to the third world.

\section{The condition \texorpdfstring{\To}{(T1o)} and Theorems A and B}\label{sec:T1o}

\subsection{The condition}\label{sec:T1ocond}

\To{} (Definition~\ref{def:T1o}) says that a strict local improvement of a
selected model at world $k$ stays strictly better, at $k$, than the nearby
selections along an edge at $k$. It holds at every complete world. More
generally it holds at world $k$ if $k$ compares every selected model with
every regular model: closedness of $A_k(Y)$ gives $C_p\not\preceq_kY$ near
$p^*$, and comparability gives $Y\preceq_kC_p$. It also holds if the strict
sections $\{p:Y\prec_kC_p\}$ are relatively open along the edges at $k$
(the variant of (S) described in CCF's Remark~21, with strict sections
$\{p:Y\prec_kC_p\}$ open). It is automatic for approaches from
the side where the selections are no better: if $Y\prec_kC_{p^*}\preceq_k
C_p$, then $Y\prec_kC_p$. A failure of \To{} is an incomparability between
$Y$ and nearby selections (Lemma~\ref{lem:Ra}).

The following results of FSE and CCF are stated for complete worlds and use
completeness only through (T1).

\begin{lemma}\label{lem:T1otools}
Assume (G) at every world and (S), and let world $v$ satisfy (IB). Let
$w\ne v$, and let $u$ be the third world.
\begin{enumerate}[label=(\alph*),leftmargin=*,itemsep=2pt]
\item \emph{(FSE, Lemma~7 (E2) and Lemma~11.)} Let $v$ satisfy \To{} along
$[e_v,e_w]$. For every $c<\infty$ and every non-null $A\subseteq
D_{vw}\cap\{t^v_w<c\}$ there is $b<\infty$ with $\Hi vwb\prec^v_dA$.
Moreover, if $C_{q_e}=D_w$ and $q_m\to q_e$ along $[e_v,e_w]$ with
$\alpha^{(v)}_{q_m}<\infty$ for all $m$, then for every $c<\infty$ and every
non-null $A'\subseteq D_{vw}\cap\{t^v_w<c\}$ there is $m_0$ such that for
all $m\ge m_0$ and every non-null $W\subseteq\Hi vw{\alpha^{(v)}_{q_m}}$ the
pair $(A',W)$ is strictly AC at $v$.
\item \emph{(FSE, Lemma~8(a).)} Let $v$ satisfy \To{} along both of its
edges. For every non-null $A\subseteq H'$ there are a non-null
$A_0\subseteq P_v$ with $A_0\preceq^v_dA$ and $b<\infty$ with $\Hi
vwb\prec^v_dA_0$ and $\Hi vub\prec^v_dA_0$. Hence every measurable subset of
these two tails is $\preceq^v_dA$, and for every open interval
$I\subseteq(0,\infty)$ and every $b'\ge b$ the \emph{debt yardstick}
$Z_I(A)\coloneqq P_v\cap\{t^v_u>b'\}\cap\{t^v_w\in I\}$ is non-null and
$\preceq^v_dA$; likewise with $w$ and $u$ exchanged.
\item \emph{(CCF, Lemmas~5 and~6, and Step~1 of the proof of Theorem~7.)}
Let (EP$_{vw}$) hold and let $v$ satisfy \To{} along $[e_v,e_w]$. Then
CCF's Lemmas~5 and~6 hold at $v$ on $[e_v,e_w]$. Moreover, for every
non-null $A_0\subseteq D_{vw}\cap\{t^v_w\ge a_1\}$ with $a_1>0$, $D_{vw}$ is
a finite union of measurable sets each $\preceq^v_dA_0$; and for every
non-null $T_0\subseteq D_{vw}\cap\{t^v_w>c\}$ with $c>0$, $D_{vw}$ is a
finite union of measurable sets each $\preceq^v_cT_0$.
\item \emph{(CCF, Theorem~7.)} If (EP) holds on both edges at $v$ and $v$
satisfies \To{} along both of them, then world $v$ is fine.
\end{enumerate}
\end{lemma}
\begin{proof}
In each source the world in question is assumed complete, and completeness
is used only through (T1): in FSE's proofs of Lemma~7 (E2) and Lemma~11 it
enters only through (T1); FSE's Lemma~8(a) uses it only through (E2); CCF's
proofs of its Lemmas~5 and~6 use completeness only through (T1) (CCF,
Section~3; checked); and CCF's Theorem~7 uses it only through those lemmas. Each application of (T1)
compares a model $Y$ with selections $C_{p}$, $p\to p^*$ along an edge at
$v$, where $Y$ is the following model:
\begin{itemize}[leftmargin=*,itemsep=1pt]
\item (E2) and FSE's Lemma~11: $Y$ is $D_w$ with its debt at $v$ on $A$,
resp.\ $A'$, silenced, and $C_{p^*}=D_w$; $D_w$ gives debt at $v$ on all of
$D_{vw}$.
\item CCF's Lemma~5(a): $Y$ is $E_w(b)$ with its debt on $A$ silenced,
$C_{p^*}=E_w(b)$; Lemma~5(c): $Y$ is $D_w$ with its debt on $A$ silenced;
Lemma~6(a): $Y=E^w_v(a')+\cred_v(T)$, $C_{p^*}=E^w_v(a')$, which is silent
on $T$; Lemma~6(c): $Y=V_{vw}+\cred_v(T)$, $C_{p^*}=V_{vw}$, which is silent
on $D_{vw}\supseteq T$.
\end{itemize}
In every case $Y$ is a strict pointwise $v$-improvement of $C_{p^*}$ and the
approach is along the edge named, so \To{} gives the conclusion
$Y\prec_vC_p$ that (T1) gives. FSE's Lemma~8(a) applies (E2) on both edges
at $v$, and otherwise uses (G) and the joint densities on $P_v$. Step~1 of
CCF's proof of Theorem~7 uses only CCF's Lemmas~4--6, (G), transitivity and
a Lebesgue-number argument; Steps~2 and~3 use Step~1 on both edges at $v$
and the joint densities.
\end{proof}

\subsection{Low-cone heaviness}\label{sec:LH}

\begin{definition}\label{def:sCN}
At world $x$, a measurable $R\subseteq H'$ is \emph{strongly credit-null}
if no pair $(A,T)$ with $A$ non-null and $T\subseteq R$ is weakly SI at $x$,
and \emph{debt-infinitesimal} if every pair $(A,S)$ with $A\subseteq R$
non-null and $S$ non-null is SI at $x$. World $x$ satisfies (LH$^+_x$) if no
non-null set is strongly credit-null at $x$, and (DH$_x$) if no non-null set
is debt-infinitesimal at $x$.
\end{definition}

Both properties are hereditary and invariant under null modifications. By
strict (G), $R$ is strongly credit-null if and only if no such pair is SI:
if $(A,T)$ is weakly SI and $A'\subseteq A$ is non-null with $A\setminus A'$
non-null, then $(A',T)$ is SI. So strongly credit-null is FSE's credit-null
(FSE, Definition~50), and (LH$^+_x$) says that for every non-null $R$ there
is an SI pair $(A,T)$ at $x$ with $A$ non-null and $T\subseteq R$. A pair
$(A,\varnothing)$ with $A$ non-null is strictly AC, which is why the debt
part is required to be non-null.

\begin{lemma}\label{lem:sCN}
Let $n=3$, assume (G) at every world and (S), and let $\{v,w,u\}=
\{1,2,3\}$. Assume (IB) at the worlds named in each part.
\begin{enumerate}[label=(\alph*),leftmargin=*,itemsep=2pt]
\item No world has both a non-null strongly credit-null set and a non-null
debt-infinitesimal set.
\item ((IB) at $v$.) If world $v$ has a non-null strongly credit-null set,
then for some $a>0$ both low cones $\Lo vwa$ and $\Lo vua$ are strongly
credit-null at $v$.
\item ((IB) at $v$.) If (EP$_{vw}$) holds and world $v$ satisfies \To{}
along $[e_v,e_w]$, then no low cone $\Lo vwa$ is strongly credit-null at
$v$.
\item ((IB) at $v$ and $w$.) Let $E$ be a selection on the open edge
$(e_v,e_w)$ with $\beta^{(v)}_E>0$, and let $R\subseteq D_{vw}$ be a set on
which $E$ gives credit at $v$. If $R$ is strongly credit-null at $v$, then
$R$ is debt-infinitesimal at $w$.
\item ((IB) at $v$.) If world $v$ satisfies \To{} along both of its edges
and has a non-null debt-infinitesimal set, then there is $b<\infty$ such
that $\Hi vwb$ and $\Hi vub$ are debt-infinitesimal at $v$.
\item ((IB) at $v$ and $w$.) If worlds $v$ and $w$ satisfy \To{} along both
of their edges and world $w$ has a non-null debt-infinitesimal set, then
world $v$ has a non-null strongly credit-null set.
\end{enumerate}
\end{lemma}
\begin{proof}
(a) Let $L$ be strongly credit-null and $R$ debt-infinitesimal at $x$, both
non-null. If $R\setminus L$ is non-null, put $A\coloneqq R\setminus L$ and
$T\coloneqq L$; otherwise split $R$ into disjoint non-null $A,T$. The pair
$(A,T)$ is SI because $R$ is debt-infinitesimal, and not weakly SI because
$L$ is strongly credit-null and $T\subseteq L$.

(b) Let $Q$ be strongly credit-null and non-null, and split it into
disjoint non-null $Q_1,Q_2$. By \ref{R:yardc} at $v$ there are $a_i>0$ such
that every measurable subset of $\Lo vw{a_i}$ and of $\Lo vu{a_i}$ is
$\preceq^v_cQ_i$; put $a\coloneqq\min(a_1,a_2)$ and let $L$ be one of the two
cones of level $a$. Let $(A,T)$ be a pair with $A$ non-null and $T\subseteq
L$, and suppose it is weakly SI. For some $i$, $A'\coloneqq A\setminus Q_i$
is non-null. The context $V+\debt_v(A')$ is silent on $Q_i\cup T$, so
$T\preceq^v_cQ_i$ and weak (G) give
\[
 V+\debt_v(A')+\cred_v(Q_i)\preceq_vV+\debt_v(A')+\cred_v(T)\preceq_v
 V+\debt_v(A)+\cred_v(T)\preceq_vV.
\]
So $(A',Q_i)$ is weakly SI, with $A'$ non-null and $Q_i\subseteq Q$: a
contradiction.

(c) Fix $0<b<a$ and a non-null $Z\subseteq D_{vw}\cap\{b<t^v_w<a\}$. By
\ref{P:approach} in world $v$'s frame, with $s^*=b$ and $s_n\uparrow b$,
there is $q^*$ with $C_{q^*}=E_w(b)$ such that every neighbourhood of $q^*$
contains points $q$ with $C_q=E_w(s)$, for all $s$ in some $[s_0,b)$. The
model $Y\coloneqq E_w(b)+\cred_v(Z)$ is a strict pointwise $v$-improvement of
$E_w(b)$, which is silent on $Z$. By \To{}, $Y\prec_vE_w(s)$ for some
$s\in[s_0,b)$. Since
$Y=E_w(s)+\debt_v(D_{vw}\cap\{s\le t^v_w<b\})+\cred_v(Z)$, (IB$'$) gives an
SI pair with non-null debt part and credit part $Z\subseteq\Lo vwa$.

(d) Let $A\subseteq R$ and $S$ be non-null and disjoint. Choose a non-null
$S_0\subseteq S$ with $S'\coloneqq S\setminus S_0$ non-null. By
\ref{R:yardc} at $w$ there is $a'>0$ such that every measurable subset of
$\Lo wu{a'}$ is $\preceq^w_cS'$. The set $Q\coloneqq P_w\cap
\{t^v_w<\beta^{(v)}_E\}\cap\Lo wu{a'}$ is non-null, $E$ gives debt at $v$ on
$Q$, and $Q\cap A=\varnothing$. Let $X$ be $E$ silenced on $Q\cup A$, so
that $E=X+\debt_v(Q)+\cred_v(A)$. The pair $(Q,A)$ is not weakly SI at $v$,
since $Q$ is non-null and $A\subseteq R$; by (IB$'$), $E\not\preceq_vX$. As
$E$ is competitive on $\{v,w\}$, $E\preceq_wX$. By \eqref{eq:swap},
$E=X+\debt_w(A)+\cred_w(Q)$, and (IB$'$) gives
$V+\debt_w(A)+\cred_w(Q)\preceq_wV$. The context $V+\debt_w(A)$ is silent on
$Q\cup S'$, so $Q\preceq^w_cS'$ and strict (G) on $S_0$ give
\[
 V+\debt_w(A)+\cred_w(S)\prec_wV+\debt_w(A)+\cred_w(S')\preceq_w
 V+\debt_w(A)+\cred_w(Q)\preceq_wV.
\]
So $(A,S)$ is SI at $w$.

(e) Let $R$ be debt-infinitesimal and non-null, and split it into disjoint
non-null $R_1,R_2$. By Lemma~\ref{lem:T1otools}(b) there are
$A^i_0\preceq^v_dR_i$ and $b_i$ with $\Hi vw{b_i},\Hi vu{b_i}
\prec^v_dA^i_0$; put $b\coloneqq\max(b_1,b_2)$, and let $H_b$ be one of the
two tails of level $b$. Let $B\subseteq H_b$ and $S$ be non-null and
disjoint. If $S\setminus R_1$ is non-null, put $i\coloneqq1$ and
$S''\coloneqq S\setminus R_1$; otherwise $S$ misses $R_2$ up to a null set,
and we put $i\coloneqq2$ and $S''\coloneqq S$. Then $S''$ is non-null and
misses $B\cup R_i$, $B\preceq^v_dR_i$, and $(R_i,S'')$ is SI. Hence
\[
 V+\debt_v(B)+\cred_v(S)\preceq_vV+\debt_v(B)+\cred_v(S'')\preceq_v
 V+\debt_v(R_i)+\cred_v(S'')\prec_vV,
\]
by weak (G), by $B\preceq^v_dR_i$ in the context $V+\cred_v(S'')$, and by
SI. So $(B,S)$ is SI.

(f) By (e) at $w$, $\Hi wvb$ is debt-infinitesimal at $w$ for some
$b<\infty$. By \ref{P:ends}(c) in world $v$'s frame there are a point $q_0$
with $C_{q_0}=D_v$ and points $q_m\to q_0$ of the open edge with
$C_m\coloneqq C_{q_m}\ne D_v$, and by \ref{P:ends}(a)
$\alpha^{(v)}_m\coloneqq\alpha^{(v)}_{q_m}\to0$. In world $w$'s frame,
$\alpha^{(w)}_m=1/\beta^{(v)}_m$ and $\beta^{(w)}_m=1/\alpha^{(v)}_m\to\infty$,
and $C_{q_0}=D_v$ gives debt at $w$ on all of $D_{vw}$.

\emph{Case I: $\beta^{(v)}_m>0$ for infinitely many $m$.} Along these $m$,
$\alpha^{(w)}_m<\infty$, and the second part of
Lemma~\ref{lem:T1otools}(a), at world $w$ on $[e_w,e_v]$ with $q_e\coloneqq
q_0$, applies to $A'\coloneqq D_{vw}\cap\{b<t^w_v<b+1\}$. Since
$\alpha^{(w)}_m\ge\beta^{(w)}_m\to\infty$, there is $m$ for which
$W\coloneqq\Hi wv{\alpha^{(w)}_m}$ is disjoint from $A'$ and $(A',W)$ is
strictly AC at $w$. But $A'\subseteq\Hi wvb$ is debt-infinitesimal, so
$(A',W)$ is SI: a contradiction.

\emph{Case II: $\beta^{(v)}_m=0$ for all large $m$.} Then
$C_m=(\alpha^{(v)}_m,0)$ with $0<\alpha^{(v)}_m<\infty$ for large $m$, and
$C_m$ is silent on $D_{vw}\cap\{t^v_w<\alpha^{(v)}_m\}$. Fix such $m$ with
$1/\alpha^{(v)}_m>b$, put $a_m\coloneqq\alpha^{(v)}_m$, and
$T'\coloneqq P_w\cap\{t^v_w<a_m\}$. $T'$ is non-null, $C_m$ is silent on
it, and $T'\subseteq\Hi wvb$ is debt-infinitesimal at $w$. We show that $T'$
is strongly credit-null at $v$. Let $A$ be non-null and $T\subseteq T'$
disjoint from $A$. If $T$ is null, $(A,T)$ is strictly AC by strict (G).
Let $T$ be non-null, and choose a non-null $A_1\subseteq A$ with
$A\setminus A_1$ non-null. Lemma~\ref{lem:T1otools}(b) at $v$ gives the
debt yardstick $S\coloneqq Z_I(A_1)\subseteq P_v\cap\{t^v_w\in I\}$, with
$I\coloneqq(a_m/4,a_m/2)$ and $S\preceq^v_dA_1$; $C_m$ is silent on $S$,
and $S\cap T=\varnothing$. Let $Y\coloneqq C_m+\debt_w(T)+\cred_w(S)$.
Since $T$ is debt-infinitesimal at $w$, $(T,S)$ is SI at $w$, and (IB$'$)
gives $Y\prec_wC_m$. By competitiveness, $C_m\preceq_vY$. By
\eqref{eq:swap}, $Y=C_m+\cred_v(T)+\debt_v(S)$, so (IB$'$) gives
$V\preceq_vV+\debt_v(S)+\cred_v(T)$. Then
\[
 V\preceq_vV+\debt_v(S)+\cred_v(T)\preceq_vV+\debt_v(A_1)+\cred_v(T)\prec_v
 V+\debt_v(A)+\cred_v(T),
\]
by $S\preceq^v_dA_1$ in the context $V+\cred_v(T)$ and by strict (G) on
$A\setminus A_1$. So $(A,T)$ is strictly AC, hence not weakly SI.
\end{proof}

\begin{theorem}\label{thm:LH}
Let $n=3$, let every ranking be a preorder with (G) and (IB), assume (S)
and \Tos. Then every world satisfies (LH$^+$) and (DH). Nothing is assumed
about $V$, and completeness is not assumed.
\end{theorem}
\begin{proof}
After relabelling, suppose that world~$1$ has a non-null strongly
credit-null set. By Lemma~\ref{lem:sCN}(b) there is $a>0$ such that $\Lo12a$ and
$\Lo13a$ are strongly credit-null, and by Lemma~\ref{lem:sCN}(c) neither
(EP$_{12}$) nor (EP$_{13}$) holds. So there are two-sided selections $E$ on
$(e_1,e_2)$ and $E'$ on $(e_1,e_3)$. The sets
\[
 R\coloneqq P_1\cap\{t^1_2>\alpha^{(1)}_E\}\cap\Lo13a,\qquad
 R'\coloneqq P_1\cap\{t^1_3>\alpha^{(1)}_{E'}\}\cap\Lo12a
\]
are non-null and strongly credit-null at~$1$. $E$ gives credit at $1$ on
$R\subseteq D_{12}$, and $E'$ on $R'\subseteq D_{13}$. By
Lemma~\ref{lem:sCN}(d), $R$ is debt-infinitesimal at~$2$ and $R'$ at~$3$. By
Lemma~\ref{lem:sCN}(f) with $(v,w)=(3,2)$, world~$3$ has a non-null
strongly credit-null set, which contradicts Lemma~\ref{lem:sCN}(a) at
world~$3$. So no world has a non-null strongly credit-null set. If some
world $w$ had a non-null debt-infinitesimal set, Lemma~\ref{lem:sCN}(f)
would give a non-null strongly credit-null set at any $v\ne w$.
\end{proof}

In the proof every passage from a non-comparison to a comparison is the
competitiveness of a selection (Lemma~\ref{lem:sCN}(d) and~(f), Case~II),
and every strict comparison comes from strict (G), from \To{}, or from an
SI or debt-infinitesimal hypothesis. For complete rankings this is FSE's
Theorem~56.

\begin{lemma}\label{lem:exchange}
Under the hypotheses of Theorem~\ref{thm:LH}, let $\{k,m,l\}=\{1,2,3\}$.
For every non-null $Z\subseteq P_m$ there is $c<\infty$ such that
$V+\debt_l(Y)+\cred_l(Z)\prec_lV$ for every non-null $Y\subseteq
P_l\cap\Hi lkc$. For $(k,m,l)=(2,1,3)$ this reads
$V+o_{12}(Y)+o_{23}(Z)\prec_3V$.
\end{lemma}
\begin{proof}
By \ref{R:yardc} at $l$ there is $a'>0$ such that every measurable subset
of $L\coloneqq\Lo lk{a'}$ is $\preceq^l_cZ$. By (LH$^+_l$)
(Theorem~\ref{thm:LH}) there is a weakly SI pair $(A,T'')$ at $l$ with $A$
non-null and $T''\subseteq L$. By Lemma~\ref{lem:T1otools}(b) at $l$ there
are $A_0\preceq^l_dA$ and $b$ with $\Hi lkb\prec^l_dA_0$. Put
$c\coloneqq\max(a',b)$ and let $Y\subseteq P_l\cap\Hi lkc$ be non-null. Then
$Y\prec^l_dA$, and $Y$ misses $T''$ (levels of $t^l_k$ above $c$ and below
$a'$) and $Z$ ($P_l$ and $P_m$). Hence
\[
 V+\debt_l(Y)+\cred_l(Z)\preceq_lV+\debt_l(Y)+\cred_l(T'')\prec_l
 V+\debt_l(A)+\cred_l(T'')\preceq_lV,
\]
by $T''\preceq^l_cZ$ in the context $V+\debt_l(Y)$, by $Y\prec^l_dA$ in the
context $V+\cred_l(T'')$, and because $(A,T'')$ is weakly SI. On $P_3$,
$o_{12}$ is debt at $3$; on $P_1$, $o_{23}$ is credit at $3$.
\end{proof}

\subsection{The endpoint dichotomy}\label{sec:dichotomy}

From here to the end of Section~\ref{sec:AB} the standing hypothesis is:
\begin{enumerate}[label=(A),leftmargin=*]
\item $n=3$; every ranking is a preorder with (G) and (IB); (S); \Tos; and
$V$ is undominated.
\end{enumerate}
So Theorem~\ref{thm:LH} applies, and $V$ is competitive on $\{1,2,3\}$. We
use the notation $T_0(E)$, $B_0(E)$ and $V+\tau_{T,B}$ of \ref{R:XOR} for a
selection $E$ on the open edge $(e_1,e_2)$.

\begin{lemma}\label{lem:fineregion}
Assume (A) and (EP$_{vw}$). For every non-null $A\subseteq H'$, $D_{vw}$ is
a finite union of measurable sets each $\preceq^v_dA$.
\end{lemma}
\begin{proof}
By Lemma~\ref{lem:T1otools}(b), $Z\coloneqq Z_{(1,2)}(A)\subseteq
D_{vw}\cap\{t^v_w>1\}$ is non-null and $Z\preceq^v_dA$. By
Lemma~\ref{lem:T1otools}(c) with $A_0\coloneqq Z$, $D_{vw}$ is a finite union
of sets $\preceq^v_dZ\preceq^v_dA$.
\end{proof}

\begin{lemma}\label{lem:B2}
Assume (A) and (EP$_{23}$). Let $E$ be a selection on the open edge
$(e_1,e_2)$ and let $B\subseteq B_0(E)$ be non-null. Then
$V+\tau_{T,B}\not\preceq_1V$ for every $T\subseteq T_0(E)$.
\end{lemma}
\begin{proof}
By Lemma~\ref{lem:exchange} with $(k,m,l)=(2,1,3)$ there is $c$ such that
$Z\coloneqq P_3\cap\Hi32c$ satisfies $V+o_{12}(Z)+o_{23}(B)\prec_3V$. By
(LH$^+_2$) there is a weakly SI pair $(A,Z_1)$ at~$2$ with $A$ non-null and
$Z_1\subseteq Z$. By Lemma~\ref{lem:fineregion} at world~$2$, $D_{23}$ is a
union of measurable sets $W_1,\dots,W_r$ with $W_j\preceq^2_dA$. Suppose
$V+\tau_{T^*,B}\preceq_1V$ for some $T^*\subseteq T_0(E)\subseteq P_2$, and
put $Q_j\coloneqq(T^*\cap W_j)\setminus(W_1\cup\dots\cup W_{j-1})$ and
$T_j\coloneqq Q_1\cup\dots\cup Q_j$. Then $Q_j\preceq^2_dA$ by weak (G),
and $(Q_j,Z)$ is weakly SI at~$2$:
\[
 V+\debt_2(Q_j)+\cred_2(Z)\preceq_2V+\debt_2(Q_j)+\cred_2(Z_1)\preceq_2
 V+\debt_2(A)+\cred_2(Z_1)\preceq_2V,
\]
by weak (G), by $Q_j\preceq^2_dA$ in the context $V+\cred_2(Z_1)$ (silent
on $Q_j\cup A$, as $Z_1\subseteq P_3$ misses $Q_j\subseteq P_2$), and
because $(A,Z_1)$ is weakly SI. The model $V+\tau_{\varnothing,B}=
V+o_{23}(B)$ has debt at $1$ on the non-null $B$, so it is not $\preceq_1V$.
Let $j^*\ge1$ be least with $V+\tau_{T_{j^*},B}\preceq_1V$. Then
$V+\tau_{T_{j^*-1},B}\not\preceq_1V$, and \ref{R:XOR} gives
$V+\tau_{T_{j^*-1},B}\preceq_2V$. Let $M\coloneqq
V+o_{13}(T_{j^*})+o_{23}(B)+o_{12}(Z)$.
\begin{itemize}[leftmargin=*,itemsep=1pt]
\item World $1$: $o_{12}$ is credit at $1$ on $P_3$, so
$M\prec_1V+\tau_{T_{j^*},B}\preceq_1V$ by strict (G).
\item World $2$: $M$ and $V+\tau_{T_{j^*-1},B}$ differ exactly on $Q_{j^*}$,
where $M$ gives $o_{13}$ (debt at $2$), and on $Z$, where $M$ gives $o_{12}$
(credit at $2$), the other model being silent there; so do
$V+\debt_2(Q_{j^*})+\cred_2(Z)$ and $V$. By (IB$'$),
$M\preceq_2V+\tau_{T_{j^*-1},B}\preceq_2V$.
\item World $3$: $o_{13}$ is credit at $3$ on $P_2$, so
$M\preceq_3V+o_{12}(Z)+o_{23}(B)\prec_3V$.
\end{itemize}
So $M$ dominates $V$: a contradiction.
\end{proof}

\begin{lemma}\label{lem:B3}
Assume (A) and (EP$_{23}$), and let $E$ be a two-sided selection on
$(e_1,e_2)$. Then $T_0(E)$ is strongly credit-null at world~$1$.
\end{lemma}
\begin{proof}
Let $A$ be non-null and $T\subseteq T_0(E)$ disjoint from $A$, and suppose
$(A,T)$ is weakly SI at~$1$. By Lemma~\ref{lem:T1otools}(b) at world~$1$
there are $A_0\preceq^1_dA$ and $b$ with $\Hi13b\prec^1_dA_0$. The set
$Z\coloneqq B_0(E)\cap\Hi13b=P_1\cap\{t^1_2<\beta^{(1)}_E\}\cap\{t^1_3>b\}$
is non-null, since $\beta^{(1)}_E>0$, and $Z\prec^1_dA$. The context
$V+\cred_1(T)$ is silent on $Z\cup A$, so
$V+\debt_1(Z)+\cred_1(T)\prec_1V+\debt_1(A)+\cred_1(T)\preceq_1V$. On $P_1$,
$o_{23}$ is debt at $1$, and on $P_2$, $o_{13}$ is credit at $1$; so
$V+\tau_{T,Z}\prec_1V$, contrary to Lemma~\ref{lem:B2}.
\end{proof}

\begin{proposition}\label{prop:dichotomy}
Assume (A). Either every edge has the endpoint property, or every edge
carries a two-sided selection.
\end{proposition}
\begin{proof}
Suppose (EP$_{23}$). If (EP$_{12}$) failed, there would be a two-sided
selection $E$ on $(e_1,e_2)$; $T_0(E)$ is non-null and, by
Lemma~\ref{lem:B3}, strongly credit-null at world~$1$, contrary to
(LH$^+_1$). So (EP$_{12}$) holds, and (EP$_{13}$) follows by exchanging the
labels $2$ and $3$, under which (A) is invariant. By relabelling, the
endpoint property on any one edge gives it on the other two. If no edge has
it, every edge has a selection omitting both endpoints, that is, a
two-sided selection.
\end{proof}

\subsection{Non-fineness on the own region}\label{sec:own}

\begin{theorem}\label{thm:own}
Assume (A). Let $\{k,m,l\}=\{1,2,3\}$, let $E$ be a two-sided selection on
$(e_k,e_m)$, and let $B_k(E)\coloneqq P_k\cap\{t^k_m<\beta^{(k)}_E\}$, the
part of $P_k$ on which $E$ gives debt at $k$. For every non-null
$T\subseteq B_k(E)$ there is $c\in(0,\infty)$ such that $T$ is not a finite union
of measurable sets $Q$ for which $(Q,P_l\cap\{t^k_l<1/c\})$ is a weakly SI
pair of world~$k$.
\end{theorem}
\begin{proof}
(A) is invariant under relabelling, so we take $(k,m,l)=(2,1,3)$. Then $E$
lies on $(e_1,e_2)$, $B_2(E)=P_2\cap\{t^1_2>\alpha^{(1)}_E\}=T_0(E)$, and
$Y_c\coloneqq P_3\cap\{t^2_3<1/c\}=P_3\cap\Hi32c$. Suppose that for every
$c$, $T$ is a finite union of measurable sets $Q$ with $(Q,Y_c)$ weakly SI
at~$2$.

\emph{Step 1: $V+\tau_{T,Z}\not\preceq_1V$ for every non-null
$Z\subseteq B_0(E)$.} By Lemma~\ref{lem:exchange} with $(k,m,l)=(2,1,3)$
there is $c$ with $V+o_{12}(Y)+o_{23}(Z)\prec_3V$ for $Y\coloneqq Y_c$. By
assumption $T=Q_1\cup\dots\cup Q_r$ with $(Q_j,Y)$ weakly SI at~$2$;
replacing $Q_j$ by $Q_j\setminus(Q_1\cup\dots\cup Q_{j-1})$ keeps this, by
weak (G). Put $T_j\coloneqq Q_1\cup\dots\cup Q_j$ and suppose
$V+\tau_{T,Z}\preceq_1V$. As in the proof of Lemma~\ref{lem:B2}, let
$j^*\ge1$ be least with $V+\tau_{T_{j^*},Z}\preceq_1V$; then
$V+\tau_{T_{j^*-1},Z}\preceq_2V$ by \ref{R:XOR}. The model
$M\coloneqq V+o_{13}(T_{j^*})+o_{23}(Z)+o_{12}(Y)$ satisfies
$M\prec_1V+\tau_{T_{j^*},Z}\preceq_1V$ (strict (G): credit at $1$ on $Y$);
$M\preceq_2V+\tau_{T_{j^*-1},Z}\preceq_2V$ (by (IB$'$), the first comparison
is the weakly SI pair $(Q_{j^*},Y)$ at~$2$); and
$M\preceq_3V+o_{12}(Y)+o_{23}(Z)\prec_3V$ (weak (G): $o_{13}$ is credit at
$3$). So $M$ dominates $V$: a contradiction.

\emph{Step 2.} By (LH$^+_1$), $T$ is not strongly credit-null at~$1$: there
is a weakly SI pair $(A,T'')$ at~$1$ with $A$ non-null and $T''\subseteq T$.
By Lemma~\ref{lem:T1otools}(b) at~$1$ there are $A_0\preceq^1_dA$ and $b$
with $\Hi13b\prec^1_dA_0$, and $Z\coloneqq B_0(E)\cap\Hi13b$ is non-null
with $Z\prec^1_dA$. The context $V+\cred_1(T'')$ is silent on $Z\cup A$, so
$V+\debt_1(Z)+\cred_1(T'')\prec_1V+\debt_1(A)+\cred_1(T'')\preceq_1V$, and
adding credit at $1$ on $T\setminus T''$, which misses $Z$, gives
$V+\tau_{T,Z}\prec_1V$. This contradicts Step~1.
\end{proof}

For complete rankings this is FSE's Theorem~64.

\subsection{Rich far ends}\label{sec:rich}

\begin{theorem}\label{thm:rich}
Assume (A), and suppose every edge carries a two-sided selection. Then for
every frame $(k,l)$ the far end of $(k,l)$ is rich.
\end{theorem}
\begin{proof}
We follow the proof of the three-world note's Proposition~6, which treats
the frame $(1,2)$ for complete rankings; (A) is invariant under relabelling,
so we treat the same frame, with third world~$3$. Write $t\coloneqq t^1_2$
on $D_{12}$, $q(s)\coloneqq(1-s)e_1+se_2$ for $s\in[0,1]$, and
$C(s)\coloneqq C_{q(s)}=(\alpha(s),\beta(s))$ in world~$1$'s frame. Let
$s_e\in(0,1]$ be the parameter of the far end; for $0<s<s_e$, $q(s)$ lies in
the open edge and $C(s)\ne D_2$. Suppose the far end is not rich. Then there
is $s_0\in(0,s_e)$ with $\alpha(s)=\infty$ for $s\in(s_0,s_e)$; for such $s$,
$\beta(s)<\infty$ (since $C(s)\ne D_2$) and $C(s)=E_2(\beta(s))$.

\emph{Step 1: $\beta$ is continuous on $(s_0,s_e)$, and $\beta(s)\to\infty$
as $s\uparrow s_e$.} Let $0\le b<\infty$ and $s\in(s_0,s_e)$. By weak (G) and
the strict comparisons of Section~\ref{sec:frames},
$C(s)\preceq_1E_2(b)$ holds if and only if $\beta(s)\le b$, and
$C(s)\preceq_2E_2(b)$ if and only if $\beta(s)\ge b$. The sections
$A_1(E_2(b))$, $A_2(E_2(b))$ are closed, so $\{s:\beta(s)\le b\}$ and
$\{s:\beta(s)\ge b\}$ are closed in $(s_0,s_e)$ for every $b$, and $\beta$
is continuous there. Since $C_{q(s_e)}=D_2$, \ref{P:ends}(a) gives
$\beta(s)\to\infty$ as $s\uparrow s_e$.

\emph{Step 2: levels above $z_0$ are crossed from both sides.} Fix
$s_1\in(s_0,s_e)$ and put $z_0\coloneqq\beta(s_1)$. Let $z>z_0$. By Step~1
there is $s_2\in(s_1,s_e)$ with $\beta>z$ on $[s_2,s_e)$. The set
$\{s\in[s_1,s_2]:\beta(s)=z\}$ is nonempty and compact; let $s^-\le s^+$ be
its least and greatest elements. Then $s_1<s^-$, $s^+<s_2$, $\beta<z$ on
$[s_1,s^-)$, $\beta>z$ on $(s^+,s_e)$, and $C(s^-)=C(s^+)=E_2(z)$.

\emph{Step 3: thin bands at levels above $z_0$.} Let $z>z_0$ and let
$A\subseteq H'$ be non-null. We show that there is $\varepsilon>0$ with
$D_{12}\cap\{z-\varepsilon<t<z\}\preceq^1_dA$ and
$D_{12}\cap\{z<t<z+\varepsilon\}\preceq^1_dA$.

(i) By Lemma~\ref{lem:T1otools}(b) at world~$1$ there is a non-null
$Z\subseteq P_1\cap\{z/4<t<z/2\}$ with $Z\preceq^1_dA$. Let
$X\coloneqq E_2(z)+\sil(Z)$; since $E_2(z)$ gives debt at $1$ on $Z$, $X$ is a
strict pointwise $1$-improvement of $E_2(z)=C(s^-)$. By \To{} at world~$1$
along $[e_1,e_2]$, $X\prec_1C_q$ for all $q$ in the edge near $q(s^-)$. By
continuity $\beta(s)\to z$ as $s\uparrow s^-$, so there is $s'\in[s_1,s^-)$
with $X\prec_1C(s')$ and $z'\coloneqq\beta(s')\in(z/2,z)$. Put
$L\coloneqq D_{12}\cap\{z'<t<z\}$ and $X_0\coloneqq E_2(z')+\sil(Z)$. Since
$Z\subseteq\{t<z'\}$, $C(s')=E_2(z')=X_0+\debt_1(Z)$, and, up to the null set
$\{t=z'\}$, $X=X_0+\debt_1(L)$. So $X_0+\debt_1(L)\prec_1X_0+\debt_1(Z)$,
where $X_0$ is silent on the disjoint sets $L$ and $Z$, and (IB$'$) gives
$L\prec^1_dZ\preceq^1_dA$.

(ii) By Lemma~\ref{lem:T1otools}(b) there is a non-null $Z'\subseteq
P_1\cap\{z+1<t<z+2\}$ with $Z'\preceq^1_dA$; $E_2(z)$ is silent on $Z'$. The
model $X'\coloneqq E_2(z)+\debt_1(Z')$ gives credit at $2$ on $Z'$, so
$X'\prec_2E_2(z)=C(s^+)$ by strict (G). By (T2), $C_q\preceq_1X'$ for all
$q$ in the open edge near $q(s^+)$; choose $s''\in(s^+,s_e)$ with
$C(s'')\preceq_1X'$ and $z''\coloneqq\beta(s'')\in(z,z+1)$. With
$L'\coloneqq D_{12}\cap\{z<t<z''\}$, $C(s'')=E_2(z)+\debt_1(L')$ up to a
null set, and $E_2(z)$ is silent on $L'\cup Z'$; so (IB$'$) gives
$L'\preceq^1_dZ'\preceq^1_dA$.

Take $\varepsilon\coloneqq\min(z-z',z''-z)$ and use weak (G).

\emph{Step 4: a band region in every $\Jid^1_A$.} Put $R\coloneqq
D_{12}\cap\{z_0+1<t<z_0+2\}$ and let $A$ be non-null. By Step~3, every
$z\in[z_0+1,z_0+2]$ has $\varepsilon_z>0$ such that
$D_{12}\cap\{|t-z|<\varepsilon_z\}$ is, up to a null set, a union of two
sets $\preceq^1_dA$. Finitely many intervals $(z-\varepsilon_z,
z+\varepsilon_z)$ cover $[z_0+1,z_0+2]$, so $R\in\Jid^1_A$.

\emph{Step 5: a contradiction.} Let $E$ be a two-sided selection on
$(e_1,e_3)$, so $0<\beta^{(1)}_E<\infty$. The set
\[
 T\coloneqq P_1\cap\{t^1_3<\beta^{(1)}_E\}\cap\{z_0+1<t<z_0+2\}
\]
is non-null, $T\subseteq R$, and $T\subseteq B_1(E)$. By
Theorem~\ref{thm:own} with $(k,m,l)=(1,3,2)$ there is $c$ such that $T$ is
not a finite union of measurable sets $Q$ for which $(Q,\Gamma)$ is weakly
SI at~$1$, where $\Gamma\coloneqq P_2\cap\{t^1_2<1/c\}$. By (LH$^+_1$)
there is a weakly SI pair $(A,\Gamma')$ at~$1$ with $A$ non-null and
$\Gamma'\subseteq\Gamma$. By Step~4, $R\in\Jid^1_A$; so, up to a null set,
$T=Q'_1\cup\dots\cup Q'_r$ with $Q'_i\subseteq T$ and $Q'_i\preceq^1_dA$
(intersect a finite cover by sets $\preceq^1_dA$ with $T$, use weak (G),
and add the null remainder to $Q'_1$). The context $V+\cred_1(\Gamma')$ is
silent on $Q'_i\cup A$, since $Q'_i\subseteq P_1$ and $A\cap\Gamma'=
\varnothing$; so
\[
 V+\debt_1(Q'_i)+\cred_1(\Gamma)\preceq_1V+\debt_1(Q'_i)+\cred_1(\Gamma')
 \preceq_1V+\debt_1(A)+\cred_1(\Gamma')\preceq_1V.
\]
So each $(Q'_i,\Gamma)$ is weakly SI, and $T$ is their union: a
contradiction.
\end{proof}

Steps~1, 2 and~4 are those of the three-world note and use no completeness.
In Step~3, the three-world note uses (T1) at a complete world~$1$ and its
debt yardsticks (FSE's Lemma~8(a)); here \To{} and
Lemma~\ref{lem:T1otools}(b) serve. In Step~5, the three-world note uses
FSE's Theorems~56 and~64, which assume completeness; here (LH$^+_1$) and
Theorem~\ref{thm:own} serve, and only weakly SI pairs are needed.

\subsection{Positive heaviness from sweeps}\label{sec:sweepheavy}

\begin{theorem}\label{thm:sweepheavy}
Assume (A) and let $E$ be a two-sided selection on $(e_1,e_2)$. For every
non-null $Q\subseteq B_0(E)$ there is a non-null $S\subseteq P_3$ with
$Q\notin\Nid^2_S$. Hence some ultrafilter $\Uf$ contains $Q$, and therefore
$P_1$, and satisfies $S\preceq^2_cP$ for every $P\in\Uf$.
\end{theorem}
\begin{proof}
On $P_1$, $o_{23}$ is credit at $2$ and at $3$; on $P_2$, $o_{13}$ is debt at
$2$ and credit at $3$; on $P_3$, $o_{12}$ is credit at $2$ and debt at $3$.

\emph{(1) An SI pair at world~$2$.} By (LH$^+_2$) there is a weakly SI pair
$(A,Q_1)$ at~$2$ with $A$ non-null and $Q_1\subseteq Q$. By
Lemma~\ref{lem:T1otools}(b) at~$2$ there are $A_0\preceq^2_dA$ and $b$ with
$\Hi23b\prec^2_dA_0$. The set $T\coloneqq T_0(E)\cap\Hi23b$ is non-null
(two ratios on $P_2$, and $\alpha^{(1)}_E<\infty$), and $T\prec^2_dA$. The
context $V+\cred_2(Q_1)$ is silent on $T\cup A$, so
\[
 V+\debt_2(T)+\cred_2(Q)\preceq_2V+\debt_2(T)+\cred_2(Q_1)\prec_2
 V+\debt_2(A)+\cred_2(Q_1)\preceq_2V.
\]
So $(T,Q)$ is SI at~$2$: $V+o_{13}(T)+o_{23}(Q)\prec_2V$.

\emph{(2) An SI pair at world~$3$.} By (LH$^+_3$) there is a weakly SI pair
$(A',T_1)$ at~$3$ with $A'$ non-null and $T_1\subseteq T$. By
Lemma~\ref{lem:T1otools}(b) at~$3$ there is $b'$ such that
$S\coloneqq P_3\cap\Hi32{b'}$ is non-null and $S\prec^3_dA'$. As in (1),
$V+\debt_3(S)+\cred_3(T)\prec_3V$, that is,
$V+o_{12}(S)+o_{13}(T)\prec_3V$.

\emph{(3) The sweep.} We apply Lemma~\ref{lem:sweep} at $x=2$ to $S$, $T$
and $Q$. The sets $S\subseteq P_3$, $T\subseteq P_2$ and $Q\subseteq P_1$
are pairwise disjoint, $S$ and $T$ are non-null, and $(T,Q)$ is SI, hence
not weakly AC, at~$2$. For measurable $\Sigma\subseteq Q$,
$X_\Sigma=V+\cred_2(S)+\debt_2(T)+\cred_2(\Sigma)=V+o_{12}(S)+o_{13}(T)
+o_{23}(\Sigma)$. Suppose \eqref{eq:cover} fails:
$X_\Sigma\not\preceq_2X'$ and $V\not\preceq_2X_\Sigma$, where
$X'\coloneqq V+o_{12}(S)$. At world~$3$, $o_{23}$ is credit on $\Sigma$, so
$X_\Sigma\preceq_3V+o_{12}(S)+o_{13}(T)\prec_3V$ by (2), and
$V\not\preceq_3X_\Sigma$. $X_\Sigma$ is $X'$ with $E$'s advice on
$A\coloneqq T\cup\Sigma\subseteq D_{12}$, since $E$ gives $o_{13}$ on
$T\subseteq T_0(E)$ and $o_{23}$ on $\Sigma\subseteq B_0(E)$. Apply
Lemma~\ref{lem:chain} with $(u,v,w)=(1,2,3)$ and $C\coloneqq E$: world~$2$
satisfies (IB), so $X_\Sigma\not\preceq_2X'$ gives $E\not\preceq_2Y$ for $Y$
the model $E$ silenced on $A$; with $V\not\preceq_2X_\Sigma$ and
$V\not\preceq_3X_\Sigma$ this contradicts the lemma. So \eqref{eq:cover}
holds, and Lemma~\ref{lem:sweep} gives $Q\notin\Nid^2_S$ and the
ultrafilter.
\end{proof}

FSE's Theorem~65 obtains non-membership in the ideals $\Kid^2_{S'}$,
$S'\prec^2_cS$, and needs completeness to make their union an ideal; here
$\Nid^2_S$ is an ideal without completeness, and its avoiders satisfy the
positive comparison $S\preceq^2_cP$ that the counting lemma needs.

\subsection{Counting heavy ultrafilters}\label{sec:count2}

For a non-null $T_0\subseteq H'$ put
\[
 \mathcal G(T_0)\coloneqq\{\Uf:\ P_1\in\Uf\text{ and }T_0\preceq^2_cP\text{
 for every }P\in\Uf\}.
\]

\begin{proposition}\label{prop:count2}
Assume (G) at every world and (S), let world~$2$ be a preorder with (IB)
that satisfies \Tw{} along $[e_2,e_1]$, and suppose the far end of the frame
$(3,2)$ is rich. Then $\mathcal G(T_0)$ is finite for every non-null $T_0$,
and $\mathcal G\coloneqq\bigcup_{T_0}\mathcal G(T_0)$ is countable.
\end{proposition}
\begin{proof}
Write $t\coloneqq t^2_3$ on $D_{23}$. Fix $T_0$. By \ref{R:yardc} at
world~$2$ there is $a>0$ such that every measurable subset of $\Lo23a$ is
$\preceq^2_cT_0$. Split $T''\coloneqq P_3\cap\{a/2<t<a\}$ into pairwise
disjoint non-null sets $T''_n$, $n\ge1$. Each $T''_n$ lies in
$D_{23}\cap\{t>a/2\}$, so Lemma~\ref{lem:near} with $(v,w)=(2,3)$ (the far
end of $(3,2)$ is rich) gives $b_n>0$ such that every measurable subset of
$D_{23}\cap\{t<b_n\}$ is $T''_n$-light at~$2$. Choose $0<\varepsilon_1<\min(b_1,a/2)$
and $0<\varepsilon_n<\min(b_n,\varepsilon_{n-1}/2)$, let $(I_n)$ enumerate the bounded open
intervals in $(0,\infty)$ with rational endpoints, and put
\[
 J_n\coloneqq P_2\cap\{t^2_1\in I_n\}\cap\{\varepsilon_n/2<t<\varepsilon_n\}.
\]
The $J_n$ are non-null (two ratios on $P_2$), pairwise disjoint (disjoint
ranges of $t$), disjoint from $P_1$ and $P_3$, and $J_n$ is $T''_n$-light at
$2$. The family $\Bfam$ of measurable subsets of the sets $J_n$ is
band-dense for $[e_2,e_1]$: a band $D_{21}\cap\{s<t^2_1<s'\}$ contains every
$J_n$ with $I_n\subseteq(s,s')$.

By Theorem~\ref{thm:cut} with $(v,w)=(2,1)$ there is a connecting chain
from $D_2$ to $D_1$ with perturbation sets $B_1,\dots,B_{\ell-1}\in\Bfam$,
$\ell\ge2$; choose $N$ with every $B_i$ in some $J_n$, $n\le N$, and put
$F\coloneqq J_1\cup\dots\cup J_N$ and $T_*\coloneqq T''_1\cup\dots\cup
T''_N$. By Lemma~\ref{lem:light}(a), $F$ is $T_*$-light at~$2$, and so is
every $B_i$. Since $T_*\subseteq\Lo23a$, $T_*\preceq^2_cT_0$.

Suppose $\mathcal G(T_0)$ has $\ell$ distinct members $\Uf^1,\dots,
\Uf^\ell$. As in the proof of Proposition~\ref{prop:count3}, there are
pairwise disjoint $K_l\in\Uf^l$ with $K_l\subseteq P_1$; they are non-null,
$K_l\subseteq D_{21}\setminus F$ since $F\subseteq P_2$, and
$T_*\preceq^2_cT_0\preceq^2_cK_l$. For all $i,l$ the sets
$B_i\subseteq P_2$, $T_*\subseteq P_3$ and $K_l\subseteq P_1$ are pairwise
disjoint, so Lemma~\ref{lem:light}(c) gives
$D_2+\debt_2(B_i)\prec_2D_2+\debt_2(K_l)$. This contradicts the counting
lemma \ref{R:count} with $(v,w)=(2,1)$. So $|\mathcal G(T_0)|\le\ell-1$.

Countability: with $a$ as above for $T_0$, and $L_k\coloneqq\Lo23{1/k}$ for
integers $k>1/a$, we have $L_k\preceq^2_cT_0$, so
$\mathcal G(T_0)\subseteq\mathcal G(L_k)$, and
$\mathcal G=\bigcup_k\mathcal G(L_k)$.
\end{proof}

\subsection{Theorems A and B}\label{sec:AB}

\begin{theorem}\label{thm:A}
Let $n=3$, let every ranking be a preorder with (G) and (IB), and assume
(S) and \Tos. If $V$ is undominated, then every edge has the endpoint
property and every world is fine.
\end{theorem}
\begin{proof}
Suppose $V$ is undominated; then (A) holds. Suppose some edge lacks the
endpoint property. By Proposition~\ref{prop:dichotomy}, every edge carries a
two-sided selection. By Theorem~\ref{thm:rich}, the far end of the frame
$(3,2)$ is rich, and Proposition~\ref{prop:count2} (\To{} implies \Tw)
shows that $\mathcal G$ is countable; list it as $\Uf^1,\Uf^2,\dots$ (the
list may be finite or empty). Let $E$ be a two-sided selection on
$(e_1,e_2)$; $B_0(E)$ is non-null since $\beta^{(1)}_E>0$. By \ref{U:small}
choose $B^i\in\Uf^i$ with $\mu(B^i)<2^{-i-1}\mu(B_0(E))$, and put
$Q\coloneqq B_0(E)\setminus\bigcup_iB^i$, a non-null set.
Theorem~\ref{thm:sweepheavy} gives a non-null $S$ and an ultrafilter $\Uf$
containing $Q$ and $P_1$ with $S\preceq^2_cP$ for every $P\in\Uf$. So
$\Uf\in\mathcal G(S)\subseteq\mathcal G$, and $\Uf=\Uf^i$ for some $i$; but
$\Uf$ contains the disjoint sets $Q$ and $B^i$. Hence every edge has the
endpoint property, and every world is fine by
Lemma~\ref{lem:T1otools}(d).
\end{proof}

\begin{theorem}\label{thm:B}
Let $n=3$, let every ranking be a preorder with (G) and (IB), and assume
(S). Suppose that at most one world is incomplete, and that this world
satisfies \To{} along both of its edges. Then $V$ is dominated.
\end{theorem}
\begin{proof}
Complete worlds satisfy \To, so \Tos{} holds. Suppose $V$ is undominated.
By Theorem~\ref{thm:A} every edge has the endpoint property. At least two
worlds are complete; by \ref{R:CCF} each of them is a (W)-world. This
contradicts \ref{R:TAW}, since every world is a preorder with (IB).
\end{proof}

\begin{remark}\label{rem:AB}
(a) \emph{Where \To{} is used.} In the proof of Theorem~\ref{thm:A}, \To{}
enters only through Lemma~\ref{lem:T1otools} ((E2), FSE's Lemma~11, the debt
yardsticks, CCF's Lemmas~5 and~6), Lemma~\ref{lem:sCN}(c), Step~3(i) of
Theorem~\ref{thm:rich}, and Theorem~\ref{thm:cut} in
Proposition~\ref{prop:count2}. Completeness is used nowhere. That $V$ is
undominated enters through \ref{R:XOR} (Lemma~\ref{lem:B2},
Theorem~\ref{thm:own}) and the competitiveness of $V$
(Lemma~\ref{lem:chain} in Theorem~\ref{thm:sweepheavy}).

(b) For complete rankings, Theorem~\ref{thm:B} gives the three-world note's
Theorem~12 again. The two proofs share the counting lemma, the
no-saturated-cut theorem and the light sets; they differ in the source of
heaviness. The three-world note uses FSE's Theorem~64 and reads a pair that
is not weakly SI as strictly AC; here the sweep of
Theorem~\ref{thm:sweepheavy} produces positive comparisons directly.
\end{remark}

\section{Sharper reductions}\label{sec:sharp}

Corollary~\ref{cor:C} of Section~\ref{sec:shape} reduces a configuration
with $V$ undominated to two residual cases: (R-a) a failure of \To{} at an
incomplete world, and (R-b) the endpoint property on every edge together
with two fine worlds that are not (W)-worlds. This section sharpens both.
Sections~\ref{sec:NCsetup}--\ref{sec:exact} reduce the extension problem
behind (R-b) to an exact cancellation condition (NC) on a single world.
Section~\ref{sec:conf} shows that (S), (G) and (IB) do not imply \To.
Section~\ref{sec:T1s} identifies the weaker forms of \To{} that the proof of
Theorem~\ref{thm:A} uses, and Section~\ref{sec:crossing} shows that the one
remaining use, at the cut of Theorem~\ref{thm:cut}, can fail only at
crossing approaches. Section~\ref{sec:twoNC} combines the two directions.
The resulting reduction is Corollary~\ref{cor:refined}.

\subsection{Change vectors}\label{sec:NCsetup}

Fix a world $v$. Put $\Omega'\coloneqq H'\times\{1,2\}$. A set
$S\subseteq\Omega'$ is measurable if its sections $S_l\coloneqq\{g:(g,l)\in
S\}$ are measurable, and null if both are null. A function $f$ on $\Omega'$
is a pair $(f_1,f_2)$ of functions on $H'$; functions are identified modulo
null sets, and $f\le f'$ means $f_l\le f'_l$ a.e.\ for $l=1,2$. A function is
\emph{simple} if it is measurable and takes finitely many values. For a
regular model $X$ put
\[
 u^X\coloneqq\bigl(1\{r_v(\cdot,X)\ge1\},\,1\{r_v(\cdot,X)\ge2\}\bigr),
\]
and for regular $X,Y$ let $\delta(X,Y)\coloneqq u^Y-u^X$ be the
\emph{change vector} from $X$ to $Y$. The map $r\mapsto(1\{r\ge1\},1\{r\ge
2\})$ is injective on $\{0,1,2\}$, with values $(0,0)$, $(1,0)$, $(1,1)$. So
at each point $\delta(X,Y)$ takes one of the seven values $(0,0)$,
$\pm(1,0)$, $\pm(0,1)$, $\pm(1,1)$, and where it is not $(0,0)$ it
determines the ranks of $X$ and $Y$ there. Let $\mathcal C$ be the set of
measurable functions on $\Omega'$ with these seven values. Every
$\delta\in\mathcal C$ is $\delta(X,Y)$ for some regular $X,Y$, since every
measurable function $H'\to\{0,1,2\}$ is the rank function of a regular
model. The \emph{acceptance set} of world $v$ is
\[
 \mathcal A_v\coloneqq\{\delta(X,Y):\ X,Y\text{ regular},\ X\preceq_vY\}.
\]
A \emph{charge} on $\Omega'$ is a finitely additive nonnegative function on
its measurable sets that vanishes on null sets. A charge $\ell$ is a pair of
charges $(\nu_1,\nu_2)$ on $H'$, $\nu_l(A)\coloneqq\ell(A\times\{l\})$, and
for simple $f$ we write $\ell(f)\coloneqq\sum_c c\,\ell(\{f=c\})$, which is
linear and monotone in $f$. For regular $X$,
$\ell(u^X)=\nu_1(R_1(X))+\nu_2(R_2(X))$ with $R_l(X)\coloneqq\{g\in
H':r_v(g,X)\ge l\}$; this is how a loss of CBOI's form acts on regular
models (Section~\ref{sec:losses}). Moreover $\ell(u^Y)-\ell(u^X)=
\ell(\delta(X,Y))$.

\begin{lemma}\label{lem:accept}
Let world $v$ be a preorder.
\begin{enumerate}[label=(\alph*),leftmargin=*,itemsep=2pt]
\item If $v$ satisfies (IB), then a function $\delta\in\mathcal C$ lies in
$\mathcal A_v$ if and only if $X\preceq_vY$ for every pair of regular models
with $\delta(X,Y)=\delta$.
\item If $v$ satisfies (G), every $\delta\in\mathcal C$ with $\delta\ge0$
lies in $\mathcal A_v$, and no $\delta\in\mathcal C$ with $\delta\le0$ and
$\delta\ne0$ lies in $\mathcal A_v$.
\item Let $v$ satisfy (IB). If $(0,f)$ and $(0,f')$ lie in $\mathcal A_v$
and $f+f'$ takes values in $\{-1,0,1\}$, then $(0,f+f')\in\mathcal A_v$. The
same holds for $(f,0)$ and $(f',0)$.
\item For measurable $Q,A\subseteq H'$: if $Q\preceq^v_dA$, then
$(0,1_A-1_Q)\in\mathcal A_v$; if $Q\preceq^v_cA$, then
$(1_A-1_Q,0)\in\mathcal A_v$. If $(D,Z)$ is a weakly SI pair at $v$, then
$(1_Z,-1_D)\in\mathcal A_v$; if it is weakly AC, then
$(-1_Z,1_D)\in\mathcal A_v$.
\end{enumerate}
\end{lemma}
\begin{proof}
(a) Regular models agree off $H'$, and on $H'$ they are determined by their
world-$v$ ranks. If $\delta(X,Y)=\delta(X',Y')$, then $\{X\ne Y\}=\{X'\ne
Y'\}=\{\delta\ne0\}$, and on this set $\delta$ determines the world-$v$ ranks,
so $X=X'$ and $Y=Y'$ there. (IB$'$) gives $X\preceq_vY\iff X'\preceq_vY'$.

(b) If $\delta(X,Y)\ge0$, then $X\le^v_\pw Y$ and $X\preceq_vY$ by weak (G).
If $\delta(X,Y)\le0$ and $\delta(X,Y)\ne0$, then $Y$ is a strict pointwise
$v$-improvement of $X$, so $Y\prec_vX$ and $X\not\preceq_vY$.

(c) At each point choose $(x,y,z)\in\{0,1\}^3$ with $y-x=f$ and $z-y=f'$:
if $f=1$, take $(0,1,1+f')$, which is possible since $f'\le0$; if $f=-1$,
take $(1,0,f')$, possible since $f'\ge0$; if $f=0$, take $y=1$ when $f'=-1$
and $y=0$ otherwise, and $x=y$, $z=y+f'$. Let $X,Y,Z$ be the regular models
with world-$v$ rank $1+x$, $1+y$, $1+z$ on $H'$. Then $\delta(X,Y)=(0,f)$
and $\delta(Y,Z)=(0,f')$, so $X\preceq_vY\preceq_vZ$ by (a), and
$\delta(X,Z)=(0,f+f')\in\mathcal A_v$. For the first layer use ranks $x,y,z$.

(d) $V$ has rank $1$ on $H'$, so $\delta(V+\debt_v(Q),V+\debt_v(A))=
(0,1_A-1_Q)$ and $\delta(V+\cred_v(A),V+\cred_v(Q))=((1-1_Q)-(1-1_A),0)$.
Likewise $\delta(V+\debt_v(D)+\cred_v(Z),V)=(1_Z,-1_D)$, and the weakly AC
case is the reverse change.
\end{proof}

For a family $\mathfrak G$ of simple functions on $\Omega'$ consider:
\begin{enumerate}[label=(NC$_{\mathfrak G}$),leftmargin=*,itemsep=2pt]
\item there are no $\gamma_1,\dots,\gamma_m\in\mathfrak G$ and
$\lambda_1,\dots,\lambda_m\ge0$ with $\sum_i\lambda_i\gamma_i\le-1$;
\end{enumerate}
\begin{enumerate}[label=(NC$^+_{\mathfrak G}$),leftmargin=*,itemsep=2pt]
\item there are no $\gamma_1,\dots,\gamma_m\in\mathfrak G$ and
$\lambda_1,\dots,\lambda_m\ge0$ with $\sum_i\lambda_i\gamma_i\le0$ and
$\sum_i\lambda_i\gamma_i\ne0$.
\end{enumerate}
Since $\Omega'$ is non-null, (NC$^+_{\mathfrak G}$) implies
(NC$_{\mathfrak G}$). World $v$ satisfies \emph{(NC)}, resp.\
\emph{(NC$^+$)}, if (NC$_{\mathcal A_v}$), resp.\ (NC$^+_{\mathcal A_v}$),
holds. If $\mathfrak G\subseteq\mathcal C$, the coefficients in a violation
may be taken rational: for fixed $\gamma_i$, the admissible $\lambda$ form a
nonempty polyhedron given by finitely many linear inequalities with integer
coefficients (one for each non-null atom of the finite algebra generated by
the level sets of the $\gamma_i$, the signs of the $\lambda_i$, and, for
(NC$^+_{\mathfrak G}$), the value $\le-1$ on one atom), and such a
polyhedron has a rational point. So a violation is a finite family of
elements of $\mathfrak G$, repetitions allowed, whose sum is $\le-1$, resp.\
$\le0$ and $\ne0$: an \emph{exact finite cancellation}. Finally, world $v$
satisfies \emph{(GC)} if every element of $\mathcal C$ that lies in the
convex cone generated by $\mathcal A_v$ lies in $\mathcal A_v$.

\subsection{The extension theorem}\label{sec:NCthm}

A charge $\ell=(\nu_1,\nu_2)$ on $\Omega'$ is \emph{strictly positive} if
$\nu_1(A)>0$ and $\nu_2(A)>0$ for every non-null measurable $A\subseteq H'$.
World $v$ has a \emph{full extension} if it has a loss $L_v$ of CBOI's form
with \eqref{eq:positive} such that $X\preceq_vY$ implies $L_v(X)\le L_v(Y)$
for all regular $X,Y$; then (W$_v$) holds, and $v$ is a (W)-world.
\ref{R:CCF} and Theorem~\ref{thm:SCext} produce full extensions.

For non-null $A\subseteq H'$, a \emph{debt cover with yardstick $A$} is a
finite family of measurable sets $Q_1,\dots,Q_q$ with $Q_k\preceq^v_dA$ and
$H'=Q_1\cup\dots\cup Q_q$ up to a null set; a \emph{credit cover} is defined
likewise with $\preceq^v_c$. World $v$ is fine if and only if every non-null
$A$ is the yardstick of a debt cover and of a credit cover. An
\emph{admissible family} $\mathfrak F_v$ at $v$ consists of the functions
$(0,1_A-1_{Q_k})$ of one debt cover and $(1_A-1_{Q'_k},0)$ of one credit
cover with yardstick $A$, for every non-null $A$, together with
$(1_Z,-1_D)$ for one SI pair $(D,Z)$ with $D$ non-null and
$(-1_{Z'},1_{D'})$ for one weakly AC pair $(D',Z')$ with $Z'$ non-null. By
Lemma~\ref{lem:accept}(d), $\mathfrak F_v\subseteq\mathcal A_v$; an
admissible family exists if and only if world $v$ is fine and has such
pairs. Put
\[
 \mathfrak W_v\coloneqq\{\delta(B,X):\ B\text{ coherent undominated},\
 X\text{ regular},\ B\preceq_vX\}\subseteq\mathcal A_v.
\]

\begin{theorem}\label{thm:NC}
Let $\mathfrak G$ be a family of simple functions on $\Omega'$.
\begin{enumerate}[label=(\alph*),leftmargin=*,itemsep=2pt]
\item If (NC$_{\mathfrak G}$) holds, there is a finitely additive
probability $\ell$ on the measurable subsets of $\Omega'$, vanishing on null
sets, with $\ell(\gamma)\ge0$ for every $\gamma\in\mathfrak G$.
\item Suppose that $\mathfrak G$ contains, for every non-null $A\subseteq
H'$, the functions $(0,1_A-1_{Q_k})$, $k\le q$, for some measurable sets
$Q_1,\dots,Q_q$ covering $H'$ up to a null set, and the functions
$(1_A-1_{Q'_k},0)$, $k\le q'$, for some measurable sets
$Q'_1,\dots,Q'_{q'}$ covering $H'$ up to a null set, and
also functions $(1_Z,-1_D)$ and $(-1_{Z'},1_{D'})$ with $D$ and $Z'$
non-null. Then every $\ell$ as in (a) is strictly positive.
\item If some strictly positive charge $\ell$ has $\ell(\gamma)\ge0$ for
every $\gamma\in\mathfrak G$, then (NC$^+_{\mathfrak G}$) holds.
\end{enumerate}
\end{theorem}
\begin{proof}
(a) Let $\Sigma$ be the family of measurable subsets of $\Omega'$ and
$\mathfrak M$ the set of finitely additive probabilities on $\Sigma$ that
vanish on null sets, a subset of $[0,1]^\Sigma$ with the product topology.
Each defining condition ($\ell(S\cup S')=\ell(S)+\ell(S')$ for disjoint $S,
S'$; $\ell(\Omega')=1$; $\ell(N)=0$ for null $N$) involves finitely many
coordinates and is closed, so $\mathfrak M$ is closed in the compact space
$[0,1]^\Sigma$ (Tychonoff~\cite{Tychonoff}) and hence compact. It is
nonempty: it contains the normalisation $P$ of the product of
$\mu|_{H'}$ and counting measure on $\{1,2\}$. For simple $\gamma$, the
value $\ell(\gamma)$ is a finite linear combination of coordinates of
$\ell$, so $C_\gamma\coloneqq\{\ell\in\mathfrak M:\ell(\gamma)\ge0\}$ is
closed. It suffices to show that every finite subfamily of these sets has
non-empty intersection.

Let $\gamma_1,\dots,\gamma_m\in\mathfrak G$, let $a_1,\dots,a_N$ be the
non-null atoms of the finite algebra generated by their level sets ($N\ge1$),
and let $M_{ik}$ be the value of $\gamma_i$ on $a_k$. Suppose first that
some $p$ in the simplex $\Delta_N$ has $\sum_kM_{ik}p_k\ge0$ for every $i$.
Then $\ell\coloneqq\sum_kp_kP(\cdot\cap a_k)/P(a_k)$ lies in $\mathfrak M$,
the null atoms have $\ell$-measure $0$, and
$\ell(\gamma_i)=\sum_kM_{ik}p_k\ge0$ for every $i$. Otherwise the continuous
function $p\mapsto\min_i(Mp)_i$ is negative on the compact set $\Delta_N$,
so $\max_{p\in\Delta_N}\min_{\lambda\in\Delta_m}\lambda^{\mathsf T}Mp<0$.
By von Neumann's minimax theorem~\cite{vN} this equals
$\min_{\lambda\in\Delta_m}\max_{p\in\Delta_N}\lambda^{\mathsf T}Mp$, so
there is $\lambda\in\Delta_m$ with $\max_k(\lambda^{\mathsf T}M)_k=-\eta<0$.
Then $\sum_i\lambda_i\gamma_i\le-\eta$ on every non-null atom, hence a.e.,
and $\sum_i(\lambda_i/\eta)\gamma_i\le-1$, contrary to (NC$_{\mathfrak G}$).

(b) Write $\ell=(\nu_1,\nu_2)$, so $\nu_1(H')+\nu_2(H')=1$. For non-null $A$,
the first cover gives $\nu_2(Q_k)\le\nu_2(A)$ for every $k$, hence
$\nu_2(H')\le q(A)\,\nu_2(A)$, where $q(A)$ is the size of the cover; the
second gives $\nu_1(H')\le q'(A)\,\nu_1(A)$. Further $\nu_1(Z)\ge\nu_2(D)$
and $\nu_2(D')\ge\nu_1(Z')$. If $\nu_2(H')=0$, then $\nu_1(Z')\le
\nu_2(D')=0$ and $\nu_1(H')\le q'(Z')\nu_1(Z')=0$: the total mass would be
$0$. If $\nu_1(H')=0$, then $\nu_2(D)\le\nu_1(Z)=0$ and
$\nu_2(H')\le q(D)\nu_2(D)=0$, again impossible. So $\nu_1(H')$ and
$\nu_2(H')$ are positive, and $\nu_l(A)\ge\nu_l(H')/q_l(A)>0$ for every
non-null $A$, where $q_1\coloneqq q'$ and $q_2\coloneqq q$.

(c) Suppose $\sum_i\lambda_i\gamma_i\le0$ and $\ne0$. This function is
simple, so it is $\le-\varepsilon$ on $S\times\{l\}$ for some $\varepsilon>0$,
some $l$ and some non-null $S\subseteq H'$. Then $0\le\sum_i\lambda_i
\ell(\gamma_i)=\ell(\sum_i\lambda_i\gamma_i)\le-\varepsilon\nu_l(S)<0$.
\end{proof}

No property of a ranking enters Theorem~\ref{thm:NC}. In the following
corollary neither completeness nor (IB) is used.

\begin{corollary}\label{cor:NCext}
Let world $v$ be fine, and let it have an SI pair whose debt part is
non-null and a weakly AC pair whose credit part is non-null.
\begin{enumerate}[label=(\roman*),leftmargin=*,itemsep=2pt]
\item World $v$ has a full extension if and only if it satisfies (NC).
\item (NC) and (NC$^+$) are equivalent at $v$.
\item If (NC$_{\mathfrak W_v\cup\mathfrak F_v}$) holds for some admissible
family $\mathfrak F_v$, then world $v$ is a (W)-world.
\end{enumerate}
\end{corollary}
\begin{proof}
(i) Suppose (NC). An admissible family $\mathfrak F_v$ exists and lies in
$\mathcal A_v$, so Theorem~\ref{thm:NC}(a),(b) with $\mathfrak
G\coloneqq\mathcal A_v$ gives a strictly positive $\ell=(\nu_1,\nu_2)$ with
$\ell\ge0$ on $\mathcal A_v$. Extend $\nu_1,\nu_2$ by $0$ to $H$ and put
$\nu_{v1}\coloneqq\nu_1+\mu|_{H\setminus H'}$, $\nu_{v2}\coloneqq\nu_2+
\mu|_{H\setminus H'}$ and $\nu_{v3}\coloneqq\mu$. These charges satisfy
\eqref{eq:positive}: they vanish on null sets, and a non-null set meets $H'$
or $H\setminus H'$ in a non-null set. The loss $L_v$ they define has
$L_v(X)=\ell(u^X)$ for regular $X$, so $X\preceq_vY$ gives
$L_v(Y)-L_v(X)=\ell(\delta(X,Y))\ge0$. Conversely, if $L_v$ is a full
extension with charges $\nu_{vl}$, the restrictions of $\nu_{v1},\nu_{v2}$
to $H'$, divided by their total mass, form a strictly positive $\ell$ with
$\ell(\delta(X,Y))\ge0$ whenever $X\preceq_vY$; Theorem~\ref{thm:NC}(c)
gives (NC$^+$), hence (NC).

(ii) (NC$^+$) implies (NC). Conversely, let $c\coloneqq\sum_i\lambda_i
\delta_i\le0$ with $\delta_i\in\mathcal A_v$, $\lambda_i\ge0$, and $c\le
-\varepsilon1_{S\times\{l\}}$ for some $\varepsilon>0$, $l$ and non-null
$S$. Put $e_1\coloneqq(1_{H'},0)$ and $e_2\coloneqq(0,1_{H'})$.

Case $l=2$. A debt cover $Q_1,\dots,Q_q$ with yardstick $S$ gives
$\sum_k(0,1_S-1_{Q_k})\le(0,q1_S-1_{H'})$, so
$c_2\coloneqq(q/\varepsilon)c+\sum_k(0,1_S-1_{Q_k})\le-e_2$, as the first
layer of $c$ is $\le0$. Let $(D',Z')$ be the weakly AC pair and
$Q'_1,\dots,Q'_{q'}$ a credit cover with yardstick $Z'$. Then
$c_1\coloneqq q'(-1_{Z'},1_{D'})+\sum_k(1_{Z'}-1_{Q'_k},0)+q'c_2\le
(-1_{H'},q'1_{D'}-q'1_{H'})\le-e_1$, and $c_1+c_2\le-1$.

Case $l=1$. With a credit cover $Q'_1,\dots,Q'_{q'}$ with yardstick $S$, in
the same way $c_1\coloneqq(q'/\varepsilon)c+\sum_k(1_S-1_{Q'_k},0)\le-e_1$. Let $(D,Z)$ be the SI pair and $Q_1,\dots,Q_q$ a debt cover
with yardstick $D$. Then $c_2\coloneqq q(1_Z,-1_D)+\sum_k(0,1_D-1_{Q_k})
+qc_1\le(q1_Z-q1_{H'},-1_{H'})\le-e_2$, and $c_1+c_2\le-1$.

In both cases $c_1+c_2$ is a nonnegative combination of elements of
$\mathcal A_v$ (Lemma~\ref{lem:accept}(d)), so (NC) fails.

(iii) Theorem~\ref{thm:NC}(a),(b) with $\mathfrak G\coloneqq\mathfrak
W_v\cup\mathfrak F_v$ gives a strictly positive $\ell$ with $\ell\ge0$ on
$\mathfrak W_v$; define $L_v$ as in (i). Then $B\preceq_vX$ with $B$
coherent undominated gives $L_v(B)\le L_v(X)$.
\end{proof}

So at a fine world with an SI pair and a weakly AC pair as above, the only
obstruction to a full extension is an exact finite cancellation among
accepted change vectors. Neither Savage's theorem nor an exchange rate
between debt and credit enters: Theorem~\ref{thm:NC}(a) produces the two
charges jointly, and the two pairs only exclude a degenerate layer.

\begin{corollary}\label{cor:NCclasses}
Let world $v$ be a preorder with (G), let it be fine, and let it have an SI
pair whose debt part is non-null and a weakly AC pair whose credit part is
non-null. In each of the following cases world $v$ satisfies (NC$^+$) and
has a full extension.
\begin{enumerate}[label=(\roman*),leftmargin=*,itemsep=2pt]
\item World $v$ belongs to a configuration with (G) at every world and (S),
satisfies (IB) and (SC), and has (EP) on both of its edges.
\item World $v$ satisfies (GC).
\item $\preceq_v$ is a \emph{filter-Pareto order}: there are a set
$\mathbb K$, a filter $\mathcal F$ on $\mathbb K$ and linear functionals
$\ell_k$, $k\in\mathbb K$, on the simple functions on $\Omega'$ such that
$X\preceq_vY$ if and only if $\{k:\ell_k(\delta(X,Y))\ge0\}\in\mathcal F$.
Examples: Pareto orders of finitely many losses of CBOI's form
($\mathbb K$ finite, $\mathcal F=\{\mathbb K\}$), and ``eventual'' orders
($\mathbb K=\mathbb N$, $\mathcal F$ the cofinite filter).
\item $\preceq_v$ is contained in a complete preorder $\preceq'$ on regular
models with strict (G) and (IB); for instance, $\preceq_v$ is an
intersection of such preorders.
\end{enumerate}
In cases (ii)--(iv) completeness of $\preceq_v$ is not assumed, and it fails
in general.
\end{corollary}
\begin{proof}
By Corollary~\ref{cor:NCext}(i) and Theorem~\ref{thm:NC}(c), in each case it
suffices to show either (NC$^+$) or the existence of a full extension.

(i) Theorem~\ref{thm:SCext} gives a full extension.

(ii) Let $c\coloneqq\sum_i\lambda_i\delta_i\le0$, $c\ne0$, with
$\delta_i\in\mathcal A_v$ and $\lambda_i\ge0$. As $c$ is simple, $c\le
-\varepsilon1_{S\times\{l\}}$ for some $\varepsilon>0$, $l$ and non-null
$S$. The function $f\coloneqq-1_{S\times\{l\}}-c/\varepsilon\ge0$ is simple,
so it is a nonnegative combination of functions $1_{A\times\{l'\}}$, which
lie in $\mathcal C$, are $\ge0$, and so lie in $\mathcal A_v$
(Lemma~\ref{lem:accept}(b)). Hence $-1_{S\times\{l\}}\in\mathcal C$ lies in
the convex cone generated by $\mathcal A_v$, and by (GC) in $\mathcal A_v$,
contrary to Lemma~\ref{lem:accept}(b).

(iii) We show (GC). Let $\delta\in\mathcal C$ with
$\delta=\sum_i\lambda_i\delta_i$, $\lambda_i\ge0$ and
$\delta_i=\delta(X_i,Y_i)$, $X_i\preceq_vY_i$. Then
$K_i\coloneqq\{k:\ell_k(\delta_i)\ge0\}\in\mathcal F$, and
$\{k:\ell_k(\delta)\ge0\}$ contains $\bigcap_iK_i\in\mathcal F$, so it lies
in $\mathcal F$. Taking regular $X,Y$ with $\delta(X,Y)=\delta$ gives
$X\preceq_vY$, so $\delta\in\mathcal A_v$. For a Pareto order of losses $L^k$
of CBOI's form, $\ell_k(\delta(X,Y))\coloneqq L^k(Y)-L^k(X)$ is linear in
$\delta(X,Y)$ by Section~\ref{sec:NCsetup}.

(iv) The debt and credit orders of $\preceq'$ contain those of $\preceq_v$,
so they are fine; they are complete, and by strict (G) and (IB) they are
fine atomless comparative probabilities in CCF's sense (CCF, proof of
Lemma~8). Split the debt part $D$ of the SI pair $(D,Z)$ of $v$ into disjoint
non-null $D_1,D_2$; then $V+\debt(D_1)+\cred(Z)\prec'V+\debt(D)+\cred(Z)
\preceq'V$, so $(D_1,Z)$ is SI for $\preceq'$, with non-null debt part. The
weakly AC pair $(D',Z')$ of $v$ has $V\preceq'V+\debt(D')+\cred(Z')$, so it
is not SI for $\preceq'$ (it is AC in CCF's sense), with non-null credit
part. CCF's proofs of its Lemma~8, Theorem~9, Lemmas~10, 12 and~13,
Theorem~14, and of the passage from $\Lambda$ to a loss of CBOI's form with
\eqref{eq:positive} in its Theorem~15, use about world~$1$ only that its
ranking is a complete preorder on regular models with (G) and (IB), that its
debt and credit orders are fine (CCF's Theorem~7), and that it has an SI pair
with non-null debt part and an AC pair with non-null credit part, which give
$0<\sigma_-$ and $\sigma_+<\infty$ (CCF's Lemma~11). So $\preceq'$ has a full
extension $L'$, and $X\preceq_vY$ implies $X\preceq'Y$ and
$L'(X)\le L'(Y)$.
\end{proof}

\subsection{Exact cancellations}\label{sec:exact}

A function $(0,f)$ or $(f,0)$ on $\Omega'$ with $f\colon H'\to\{-1,0,1\}$
measurable is a \emph{single-layer} vector; by
Lemma~\ref{lem:accept}(c), at a world with (IB) the single-layer part of
$\mathcal A_v$ is closed under every sum that stays $\{-1,0,1\}$-valued.

\begin{lemma}\label{lem:smallcanc}
Let world $v$ be a preorder with (G) and (IB). Let $f_1,\dots,f_m$ be
measurable $\{-1,0,1\}$-valued functions on $H'$, repetitions allowed, with
$(0,f_j)\in\mathcal A_v$ for every $j$, $c\coloneqq\sum_jf_j\le0$ and
$c\ne0$. Then the set of points at which at least two of the $f_j$ equal
$+1$ is non-null; in particular $m\ge4$. The same holds with $(f_j,0)$ in
place of $(0,f_j)$.
\end{lemma}
\begin{proof}
Suppose that at a.e.\ point at most one $f_j$ equals $+1$. The $f_j$ are
simple, so $H'$ is a finite union of measurable cells on which
$(f_1,\dots,f_m)$ is constant; choose a non-null cell $S'\subseteq\{c<0\}$
and renumber so that the $f_j$ equal to $+1$ on $S'$, if there is one, is
$f_1$. Put $F_0\coloneqq0$ and, for $k=1,\dots,m$,
\[
 F'_k\coloneqq F_{k-1}+1_{\{F_{k-1}=-1\}\cap\{f_k=-1\}},\qquad
 F_k\coloneqq F'_k+f_k.
\]
At each point, $F$ stays in $\{-1,0\}$ until the index $k^*$ at which $f_k$
is $+1$, if there is one: a value $-1$ of $f_k$ is preceded by a raise from
$-1$ to $0$ where necessary. So $F_{k^*-1}\le0$ and $F_{k^*}\in\{0,1\}$; after
$k^*$ a value $-1$ again moves $F$ down by one, after a raise if $F=-1$.
Hence every $F'_k$ and $F_k$ is $\{-1,0,1\}$-valued. The added functions
$(0,1_{\{F_{k-1}=-1\}\cap\{f_k=-1\}})$ are $\ge0$ and lie in $\mathcal A_v$
by Lemma~\ref{lem:accept}(b), so by Lemma~\ref{lem:accept}(c) every
$(0,F'_k)$ and $(0,F_k)$ lies in $\mathcal A_v$. At a point without a $+1$,
$F_m\le0$, and $F_m=-1$ if some $f_k$ equals $-1$. At a point with $+1$ at
$k^*$, since $c\le0$ some $f_k$ equals $-1$: if one comes after $k^*$, then
$F_m\le0$; if all come before, then $F_{k^*-1}=-1$ and $F_m=0$. On $S'$,
either no $f_k$ is $+1$, or $f_1=+1$ and, as $c\le-1$ there, at least two
later $f_k$ equal $-1$; in both cases $F_m=-1$. So $(0,F_m)\in\mathcal A_v$
is $\le0$ and $\ne0$, contrary to Lemma~\ref{lem:accept}(b). If $m\le3$, two
values $+1$ at a point would force $c\ge1$ there; so the set in question is
null, which we have excluded. The credit layer is the same.
\end{proof}

So at a world with (G) and (IB), an exact cancellation inside one layer
needs at least four accepted changes, counted with multiplicity, two of
which add debt (or remove credit) on a common non-null set. The following
example has exactly this size; it violates (NC$^+$) but not (NC). Its world
is not fine. Example~\ref{ex:half5} violates (NC) itself, with five
accepted changes in each layer; it is not fine either.

\begin{example}\label{ex:K4}
Let $R_{ij}$ ($1\le i<j\le4$) and $h$ be seven pairwise disjoint non-null
measurable subsets of $H'$, and write $R_{ji}\coloneqq R_{ij}$. For
$i=1,\dots,4$ put
\[
 f_i\coloneqq\sum_{j\ne i}1_{R_{ij}}-\sum_{k<l,\ i\notin\{k,l\}}1_{R_{kl}}
 -1_h .
\]
So $f_i=+1$ on the three sets $R_{ij}$, $j\ne i$, and $f_i=-1$ on the other
three and on $h$. On $R_{kl}$ two of the $f_i$ are $+1$ and two are $-1$, so
$\sum_if_i=-4\cdot1_h$. Let $\mathcal Q$ be the set of measurable
$\{-1,0,1\}$-valued $f$ on $H'$ with $f\ge0$ or $f\ge f_i$ for some $i$, and
for regular $X,Y$ define
\[
 X\preceq^*Y\ :\Longleftrightarrow\ \delta(X,Y)_1\ge0\ \text{ and }\
 \delta(X,Y)_2\in\mathcal Q .
\]
Then $\preceq^*$ is a preorder on regular models with strict (G) and (IB).
The vectors $(0,f_i)$ lie in its acceptance set, and their sum
$(0,-4\cdot1_h)$ is $\le0$ and $\ne0$; so (NC$^+$) fails, and $\preceq^*$
has no full extension, by Theorem~\ref{thm:NC}(c). Every element of its
acceptance set has a nonnegative first layer, so it satisfies (NC): without
fineness (NC) and (NC$^+$) differ. Its credit order is inclusion
($B\preceq^*_cC$ if and only if $B\subseteq C$ up to a null set), so it is not
fine.
\end{example}
\begin{proof}
$\preceq^*$ depends only on $\delta(X,Y)$, so (IB$'$) holds, and it is
reflexive. Transitivity: let $X\preceq^*Y\preceq^*Z$. The first layers add
up to $\delta(X,Z)_1\ge0$. Put $f\coloneqq\delta(X,Y)_2$ and
$g\coloneqq\delta(Y,Z)_2$; then $f+g=\delta(X,Z)_2$ is $\{-1,0,1\}$-valued.
If $f,g\ge0$, then $f+g\ge0$; if $f\ge f_i$ and $g\ge0$, or $f\ge0$ and
$g\ge f_i$, then $f+g\ge f_i$. If $f\ge f_i$ and $g\ge f_j$, then $f=g=1$ on
$R_{ij}$ if $i\ne j$, and on $R_{ik}$, $k\ne i$, if $i=j$; so $f+g=2$ on a
non-null set, which is impossible. Weak (G): $X\le^v_\pw Y$ gives
$\delta(X,Y)\ge0$. Strict (G): let $X$ be a strict pointwise
$v$-improvement of $Y$. Then $X\preceq^*Y$ by weak (G), and
$\delta(Y,X)\le0$, $\delta(Y,X)\ne0$; either its first layer is negative on
a non-null set, or its second layer is $\le0$ and $\ne0$, hence neither
$\ge0$ nor $\ge f_i$ (each $f_i$ is $+1$ on a non-null set); so
$Y\not\preceq^*X$. The vectors $(0,f_i)$ are change
vectors with first layer $0$ and second layer $f_i\in\mathcal Q$. Credit
order: $\delta(V+\cred(C),V+\cred(B))=(1_C-1_B,0)$, which is accepted if and
only if $1_C\ge1_B$ a.e.
\end{proof}

\begin{example}\label{ex:half5}
Let $R_S$, where $S$ ranges over the ten $2$-element subsets of
$\{1,\dots,5\}$, be pairwise disjoint non-null measurable sets whose union is
$H'$ up to a null set. For $i=1,\dots,5$ put
\[
 g_i\coloneqq\sum_{S\ni i}1_{R_S}-\sum_{S\not\ni i}1_{R_S}.
\]
So $g_i=+1$ on the four sets $R_S$ with $i\in S$ and $g_i=-1$ on the other
six. On $R_S$ two of the $g_i$ are $+1$ and three are $-1$, so
$\sum_ig_i=-1_{H'}$. Let $\mathcal Q'$ be the set of measurable
$\{-1,0,1\}$-valued $g$ on $H'$ with $g\ge0$ or $g\ge g_i$ for some $i$, and
for regular $X,Y$ define
\[
 X\preceq^\circ Y\ :\Longleftrightarrow\ \delta(X,Y)_1\in\mathcal Q'\
 \text{ and }\ \delta(X,Y)_2\in\mathcal Q' .
\]
Then $\preceq^\circ$ is a preorder on regular models with strict (G) and
(IB). The ten vectors $(g_i,0)$ and $(0,g_i)$, $i=1,\dots,5$, lie in its
acceptance set, and
\[
 \sum_{i=1}^5(g_i,0)+\sum_{i=1}^5(0,g_i)=(-1_{H'},-1_{H'})\le-1 :
\]
(NC) fails. The world is not fine: if a non-null $T$ lies inside one set
$R_S$, then $B\preceq^\circ_cT$ if and only if $B\subseteq T$ up to a null
set.
\end{example}
\begin{proof}
First, if $g,g'\in\mathcal Q'$ and $g+g'$ is $\{-1,0,1\}$-valued, then
$g+g'\in\mathcal Q'$. If $g,g'\ge0$, then $g+g'\ge0$; if $g\ge g_i$ and
$g'\ge0$, or $g\ge0$ and $g'\ge g_i$, then $g+g'\ge g_i$. If $g\ge g_i$ and
$g'\ge g_j$, then $g=g'=1$ on $R_{\{i,j\}}$ if $i\ne j$, and on every $R_S$
with $i\in S$ if $i=j$; so $g+g'=2$ on a non-null set, which is impossible.
The argument is pointwise and applies to arbitrary measurable $g,g'$.

$\preceq^\circ$ depends only on $\delta(X,Y)$, so (IB$'$) holds, and it is
reflexive since $0\in\mathcal Q'$. Transitivity: if $X\preceq^\circ
Y\preceq^\circ Z$, each layer of $\delta(X,Z)=\delta(X,Y)+\delta(Y,Z)$ is
the sum of two elements of $\mathcal Q'$ and is $\{-1,0,1\}$-valued, as a
layer of a change vector; so it lies in $\mathcal Q'$. Weak (G):
$X\le^v_\pw Y$ gives $\delta(X,Y)\ge0$, and both layers lie in
$\mathcal Q'$. Strict (G): let $X$ be a strict pointwise $v$-improvement of
$Y$. Then $X\preceq^\circ Y$, and $\delta(Y,X)\le0$, $\delta(Y,X)\ne0$, so
one of its layers is $\le0$ and $\ne0$; that layer is neither $\ge0$ nor
$\ge g_i$ for any $i$, since $g_i=+1$ on a non-null set; so
$Y\not\preceq^\circ X$. The vectors $(g_i,0)$ and $(0,g_i)$ lie in
$\mathcal C$, and both of their layers lie in $\mathcal Q'$. Credit order:
$B\preceq^\circ_cT$ means $\delta(V+\cred(T),V+\cred(B))=(1_T-1_B,0)$ is
accepted, that is, $1_T-1_B\ge0$ or $1_T-1_B\ge g_i$ for some $i$. The
second alternative requires $T$ to contain the four sets $R_S$ with $i\in
S$, which is impossible for $T$ inside one $R_S$. So $B\preceq^\circ_cT$ if
and only if $B\subseteq T$ up to a null set, and $H'$ is not a finite union
of such sets.
\end{proof}

Example~\ref{ex:half5} is the case $m=5$ of a ``half-subsets'' pattern; it
respects Lemma~\ref{lem:smallcanc}: each layer uses five accepted changes,
and any two of them add debt (or remove credit) together on a set $R_S$.

\subsection{Silence penalties: (S) does not imply \texorpdfstring{\To}{(T1o)}}\label{sec:conf}

In this subsection rankings are defined on all models, regular or not, as in
TAW, Section~1.1. A model $M$ gives at each $g\in H$ an advice
$a_M(g)\in\{+,0,-,\mathrm c\}$ (accept $g$, silence, accept $-g$,
contradiction), and $r_i(g,M)\in\{0,1,2,3\}$ is its rank at world $i$: $0$ if
the accepted option is good at $i$, $1$ for silence, $2$ if it is bad at $i$,
$3$ for contradiction. (G) at world $i$ is read with $r_i(\cdot,X)\le
r_i(\cdot,Y)$ on all of $H$, (IB$'$) for all models, and (S) as in
Section~\ref{sec:axioms}, with competitiveness and closed sections for every
model. For $p\in\Delta$ let $\Dscr_p\coloneqq\Dscr_{\{p\}}$; it accepts $g$
if $p\cdot g>0$ and $-g$ if $p\cdot g<0$, and is silent on the null set
$\{p\cdot g=0\}$; $\Dscr_{e_i}=D_i$. Put
\[
 L_i(M)\coloneqq\int_H|g_i|\,r_i(g,M)\,d\mu(g)\qquad(i=1,2,3),
\]
a loss of CBOI's form with $\nu_{il}\coloneqq|g_i|\mu$ for $l=1,2,3$; these
satisfy \eqref{eq:positive}, since $|g_i|>0$ a.e.\ and $|g_i|$ is
integrable. Fix a measurable $\theta\colon H\to[0,1]$ and put
\[
 \sigma_\theta(M)\coloneqq\int_H\theta(g)\,g_1\,1\{r_1(g,M)=1\}\,d\mu(g),
 \qquad U_\theta\coloneqq L_1+\sigma_\theta .
\]
The pointwise world-$1$ costs of the ranks $0,1,2,3$ are $0$, $g_1$, $2g_1$,
$3g_1$ for $L_1$ and $0$, $(1+\theta)g_1$, $2g_1$, $3g_1$ for $U_\theta$.

\emph{The configuration $\mathrm{Conf}(\theta)$.} World~$1$: $X\preceq_1Y$
if and only if $L_1(X)\le L_1(Y)$ and $U_\theta(X)\le U_\theta(Y)$. Worlds
$i=2,3$: $X\preceq_iY$ if and only if $L_i(X)\le L_i(Y)$. Selections:
$C_p\coloneqq\Dscr_p$ for $p\in\Delta$. For $\theta=0$ this is the
configuration of three complete additive worlds with precise selections. For
$\theta\equiv1$, $U_1(M)=2\int_Hg_1\,1\{r_1(g,M)\ge1\}\,d\mu$ for every
non-contradictory $M$: world~$1$ is the Pareto order of the expected loss
$L_1$ and of the $g_1\mu$-mass of the problems on which the model does not
give the option good at $1$. We call $\mathrm{Conf}(1)$ the \emph{K-world}.

\begin{theorem}\label{thm:conf}
For every measurable $\theta\colon H\to[0,1]$, in $\mathrm{Conf}(\theta)$
every ranking is a preorder with strict (G) and (IB), and (S) holds with
$C_p=\Dscr_p$. Worlds $2$ and $3$ are complete. If $\theta$ is not a.e.\
$0$, world~$1$ is incomplete; if $\{\theta>0\}\cap H'$ is non-null, it is
incomplete on regular models. $V$ is dominated, and every world has a full
extension (world~$1$ the loss $L_1$).
\end{theorem}
\begin{proof}
(a) \emph{Preorders.} Each ranking is a conjunction of relations
$F(X)\le F(Y)$ with real-valued $F$; all values are finite, since ranks are
at most $3$ and $|g_i|$ is integrable.

(b) \emph{(G).} If $r_i(\cdot,X)\le r_i(\cdot,Y)$ on $H$, then
$L_i(X)\le L_i(Y)$, and $U_\theta(X)\le U_\theta(Y)$ because the pointwise
costs of $U_\theta$ are non-decreasing in the rank ($\theta\le1$). If
moreover the inequality of ranks is strict on a non-null set, then
$L_i(X)<L_i(Y)$, as $|g_i|>0$ a.e.; so $Y\not\preceq_iX$ and $X\prec_iY$.

(c) \emph{(IB).} Each $F\in\{L_1,L_2,L_3,U_\theta\}$ has the form
$F(M)=\int f(g,a_M(g))\,d\mu$ for a pointwise cost $f$. Hence
$F(X)-F(Y)=\int_{\{X\ne Y\}}[f(g,a_X(g))-f(g,a_Y(g))]\,d\mu$ depends only on
the set $\{X\ne Y\}$ and the advice of $X$ and $Y$ on it. So every
comparison $F(X)\le F(Y)$, and every Boolean combination of such
comparisons, satisfies (IB$'$).

(d) \emph{Incompleteness.} Let $\{\theta>0\}$ be non-null. Choose a non-null
$A\subseteq\{\theta>0\}$ and a measurable $W\subseteq H\setminus A$ with
$\int_Ag_1\,d\mu<2\int_Wg_1\,d\mu<\int_A(1+\theta)g_1\,d\mu$; this is
possible because $g_1\mu$ is a finite atomless measure, $\int_A\theta
g_1\,d\mu>0$, and $A$ may be taken of small $g_1\mu$-mass. Let $Y$ be $D_1$
silenced on $A$, and $X$ the model $D_1$ with the option bad at $1$ on $W$
($X$ need not be regular). $D_1$ has rank $0$ at world~$1$ everywhere, so
$L_1(Y)=\int_Ag_1<2\int_Wg_1=L_1(X)$ and
$U_\theta(Y)=\int_A(1+\theta)g_1>2\int_Wg_1=U_\theta(X)$. So
$X\parallel_1Y$. If $\{\theta>0\}\cap H'$ is non-null, $A$ and $W$ can be
chosen inside $H'$, and then $X$ and $Y$ are regular.

(e) \emph{Pointwise Bayes property.} For $p\in\Delta$ and $g\in H$ put
$P\coloneqq\sum_{i:g_i>0}p_i|g_i|$ and $N\coloneqq\sum_{i:g_i<0}p_i|g_i|$, so
$p\cdot g=P-N$, and $\varphi_p(g,a)\coloneqq\sum_ip_i|g_i|r_i(g,a)$ for an
advice $a$. Then $\varphi_p(g,+)=2N$, $\varphi_p(g,0)=P+N$,
$\varphi_p(g,-)=2P$ and $\varphi_p(g,\mathrm c)=3(P+N)$. If $p\cdot g>0$,
then $P>N\ge0$ and $2N<P+N<2P<3(P+N)$, so $+$ is the unique minimiser; if
$p\cdot g<0$, $-$ is. Hence $\Dscr_p$ minimises
$\Phi_p(M)\coloneqq\sum_ip_iL_i(M)=\int\varphi_p(g,a_M(g))\,d\mu$, and every
$M$ with $\Phi_p(M)=\Phi_p(\Dscr_p)$ equals $\Dscr_p$ a.e. Further
$\Psi_p(M)\coloneqq p_1U_\theta(M)+p_2L_2(M)+p_3L_3(M)$ has pointwise cost
$\varphi_p(g,a)+p_1\theta(g)g_1\,1\{a=0\}$, which equals $\varphi_p$ at
$a\ne0$ and is $\ge\varphi_p$ at $a=0$; so $\Dscr_p$ is also the unique
minimiser of $\Psi_p$, up to null sets.

(f) \emph{Competitiveness on $S\coloneqq\supp p$.} Let $X$ be any model and
suppose $\Dscr_p\not\preceq_iX$ for every $i\in S$. For $i\in S\setminus
\{1\}$, $L_i(X)<L_i(\Dscr_p)$. If $1\notin S$, then
$\Phi_p(X)<\Phi_p(\Dscr_p)$. If $1\in S$, then either
$L_1(X)<L_1(\Dscr_p)$, and $\Phi_p(X)<\Phi_p(\Dscr_p)$, or
$U_\theta(X)<U_\theta(\Dscr_p)$, and $\Psi_p(X)<\Psi_p(\Dscr_p)$; the terms
with $p_i=0$ vanish. Both contradict (e).

(g) \emph{Undominated.} If $X\preceq_i\Dscr_p$ for every $i$, then
$L_i(X)\le L_i(\Dscr_p)$ for every $i$, since $\preceq_1$ includes the
$L_1$-inequality; so $\Phi_p(X)\le\Phi_p(\Dscr_p)$ and $X=\Dscr_p$ a.e.\ by
(e). No model is strictly better than $\Dscr_p$ at any world.

(h) \emph{Closed sections.} If $p_n\to p$ in $\Delta$, then
$a_{\Dscr_{p_n}}(g)\to a_{\Dscr_p}(g)$ for every $g$ with $p\cdot g\ne0$,
that is, a.e.; the integrands of $L_i(\Dscr_{p_n})$ are bounded by
$2|g_i|$, so $p\mapsto L_i(\Dscr_p)$ is continuous by dominated
convergence. Moreover $\sigma_\theta(\Dscr_p)=0$, as $\Dscr_p$ is silent
only on a null set, so $U_\theta(\Dscr_p)=L_1(\Dscr_p)$. Hence for every
model $X$ the sections $A_i(X)=\{p:L_i(\Dscr_p)\le L_i(X)\}$, $i=2,3$, and
$A_1(X)=\{p:L_1(\Dscr_p)\le L_1(X)\}\cap\{p:L_1(\Dscr_p)\le U_\theta(X)\}$
are closed.

(i) \emph{Coherence.} $\Dscr_p$ is coherent with credal set $\{p\}$, in the
face spanned by $\supp p$.

(j) \emph{$V$ dominated; extensions.} Worlds $2$ and $3$ are ordered by the
losses $L_2$, $L_3$ of CBOI's form, so they have full extensions and are
(W)-worlds; world~$1$ is a preorder with (IB). By \ref{R:TAW}, $V$ is
dominated. Finally $X\preceq_1Y$ implies $L_1(X)\le L_1(Y)$.
\end{proof}

For $j\in\{2,3\}$ use world~$1$'s frame on $[e_1,e_j]$, with $t\coloneqq
t^1_j$ on $D_{1j}$. For $p=(1-\lambda)e_1+\lambda e_j$ put
$\rho_p\coloneqq\lambda/(1-\lambda)\in[0,\infty]$. On $D_{1j}$, where
$g_j<0<g_1$, $p\cdot g>0$ if and only if $t>\rho_p$, and off $D_{1j}$,
$p\cdot g>0$; so $\Dscr_p=(\rho_p,\rho_p)$ in world~$1$'s frame. The map
$p\mapsto\rho_p$ is continuous and increasing from $e_1$ to $e_j$, and for
$p'$ between $e_1$ and $p$ the models $\Dscr_{p'}$ and $\Dscr_p$ differ
exactly on the band $B(p',p)\coloneqq D_{1j}\cap\{\rho_{p'}<t<\rho_p\}$,
where the first gives credit and the second debt at $1$.

\begin{proposition}\label{prop:confT1}
In $\mathrm{Conf}(\theta)$, worlds $2$ and $3$ satisfy \To{} along both of
their edges. Let $j\in\{2,3\}$.
\begin{enumerate}[label=(\alph*),leftmargin=*,itemsep=2pt]
\item World~$1$ satisfies \To{} along $[e_1,e_j]$ if and only if
$D_{1j}\cap\{\theta=1\}$ is null.
\item Let $p^*\in(e_1,e_j]$, $\rho^*\coloneqq\rho_{p^*}$, and let
$A\subseteq D_{1j}\cap\{\theta=1\}\cap\{t<\rho^*\}$ be non-null. Then
$Y\coloneqq\Dscr_{p^*}+\sil(A)$ is a strict pointwise $1$-improvement of
$C_{p^*}$, $Y\not\preceq_1\Dscr_p$ for every $p\in[e_1,p^*)$, and
$Y\parallel_1\Dscr_p$ for all $p\in[e_1,p^*)$ near $p^*$.
\end{enumerate}
In particular, in the K-world \To{} fails at world~$1$ at every point of
$(e_1,e_2]$ and of $(e_1,e_3]$. So (S), strict (G) and (IB) do not imply
\To{}, even with two complete additive worlds, full extensions at all three
worlds and precise selections.
\end{proposition}
\begin{proof}
Worlds $2$ and $3$ are complete, and (T1) gives \To{} there. Let
$p^*\in[e_1,e_j]$, $\Dscr^*\coloneqq\Dscr_{p^*}$, $\rho^*\coloneqq
\rho_{p^*}$, and let $Y$ be a strict pointwise $1$-improvement of
$\Dscr^*$. $\Dscr^*$ has rank $0$ at $1$ off $D_{1j}\cap\{t<\rho^*\}$ and
rank $2$ on it; so $Y$ differs from $\Dscr^*$ exactly on a non-null
$I\subseteq D_{1j}\cap\{t<\rho^*\}$, on which it gives silence (on $I_s$) or
credit (on $I_c$). Hence
\[
 U_\theta(Y)=U_\theta(\Dscr^*)-\int_{I_s}(1-\theta)g_1-2\int_{I_c}g_1,
 \qquad L_1(Y)<L_1(\Dscr^*),
\]
and $U_\theta(Y)=U_\theta(\Dscr^*)$ if and only if $I_c$ is null and
$\theta=1$ a.e.\ on $I_s$. Recall that $p\mapsto L_1(\Dscr_p)=
U_\theta(\Dscr_p)$ is continuous.

If $U_\theta(Y)<U_\theta(\Dscr^*)$, both strict inequalities persist near
$p^*$, and $Y\prec_1\Dscr_p$ for all $p$ near $p^*$. If
$U_\theta(Y)=U_\theta(\Dscr^*)$, then $Y=\Dscr^*+\sil(A)$ with $A\coloneqq
I_s\subseteq\{\theta=1\}$; for $p\in[e_1,p^*)$ the band $B(p,p^*)$ is
non-null and
\[
 U_\theta(\Dscr_p)=L_1(\Dscr_p)=L_1(\Dscr^*)-2\int_{B(p,p^*)}g_1<
 L_1(\Dscr^*)=U_\theta(\Dscr^*)=U_\theta(Y),
\]
so $Y\not\preceq_1\Dscr_p$; and $L_1(\Dscr_p)\to L_1(\Dscr^*)>L_1(Y)$ as
$p\to p^*$, so $\Dscr_p\not\preceq_1Y$ for $p$ near $p^*$. For $p$ between
$p^*$ and $e_j$, $\Dscr^*\le^1_\pw\Dscr_p$, and $Y\prec_1\Dscr^*
\preceq_1\Dscr_p$. This proves (b). For (a): if $D_{1j}\cap\{\theta=1\}$ is
null, only the first case occurs, and \To{} holds. If it is non-null, (b)
with $p^*=e_j$ ($\rho^*=\infty$) and $A\coloneqq D_{1j}\cap\{\theta=1\}$
shows that \To{} fails. In the K-world every
non-null $A\subseteq D_{1j}\cap\{t<\rho^*\}$ qualifies.
\end{proof}

\begin{remark}\label{rem:conf}
(a) In $\mathrm{Conf}(\theta)$ the world violating \To{} is a (W)-world, and
$V$ is dominated. By Corollary~\ref{cor:Wforces}, in a configuration in
which every world is a preorder with (IB) and $V$ is undominated, a
(W)-world can occur only if both other worlds are incomplete and are not
(W)-worlds. So a failure of \To{} at a (W)-world, as in
$\mathrm{Conf}(\theta)$, can occur in such a configuration only next to two
incomplete worlds without (W).

(b) In the K-world, $\preceq^1_d$ and $\preceq^1_c$ are the orders of the
measure $g_1\mu$: $V+\debt_1(B)$ and $V$ have the same $U_1$, and
$L_1(V+\debt_1(B))=L_1(V)+\int_Bg_1$; $L_1$ and $U_1$ both decrease by a
multiple of $\int_Cg_1$ under credit on $C$. So world~$1$ is fine, with
complete debt and credit orders. For a pair $(A,T)$ and
$X\coloneqq V+\debt_1(A)+\cred_1(T)$, $L_1(X)-L_1(V)=\int_Ag_1-\int_Tg_1$ and
$U_1(X)-U_1(V)=-2\int_Tg_1$. So $(A,T)$ is weakly SI if and only if
$\int_Ag_1\le\int_Tg_1$, and weakly AC if and only if $T$ is null. A pair
with $T$ non-null and $\int_Ag_1>\int_Tg_1$ is neither SI nor weakly AC:
(SC)(iii) fails at world~$1$, while world~$1$ has the full extension $L_1$
and satisfies (NC$^+$) (Theorem~\ref{thm:NC}(c)). So, for worlds of
configurations satisfying (S), (NC$^+$) does not imply (SC). The world has
no weakly AC pair with non-null credit part; (EP) fails on its edges, so
CCF's Lemma~11 does not apply. The same computation shows that the
conclusion of FSE's Lemma~11 (the second part of
Lemma~\ref{lem:T1otools}(a)), which asserts strictly AC pairs with non-null
credit part, fails at world~$1$ on both of its edges, although all of FSE's hypotheses
other than completeness hold there (and \To{} fails there, by
Proposition~\ref{prop:confT1}).
\end{remark}

\subsection{What the proof of Theorem~A uses}\label{sec:T1s}

For a regular model $M$ and a world $k$ let $\mathrm{Cr}_k(M)\coloneqq\{g\in
H':r_k(g,M)=0\}$ be its credit set at $k$. For a selection $C_p$ on an edge
$[e_k,e_j]$ and a regular model $Y$, let $C_p\langle Y\rangle$ be $C_p$
silenced on $\mathrm{Cr}_k(C_p)\setminus\mathrm{Cr}_k(Y)$: $C_p$ stripped of
the credit at $k$ that $Y$ does not have. Then $C_p\le^k_\pw C_p\langle
Y\rangle$, so $C_p\preceq_kC_p\langle Y\rangle$ by weak (G).

\begin{definition}\label{def:Ts}
Let $k\ne j$ be worlds and $t\coloneqq t^k_j$ on $D_{kj}$. In (i) and (ii),
$p^*$ ranges over $[e_k,e_j]$, $Y$ over the strict pointwise
$k$-improvements of $C_{p^*}$, and neighbourhoods are taken in $[e_k,e_j]$.
World $k$ satisfies along $[e_k,e_j]$:
\begin{enumerate}[label=(\roman*),leftmargin=*,itemsep=2pt]
\item \Ts{} if for all $p^*,Y$ there is a neighbourhood $N$ of $p^*$ with
$Y\prec_kC_p\langle Y\rangle$ for every $p\in N$;
\item \Tcc{} if for all $p^*,Y$ there is a neighbourhood $N$ of $p^*$ with
$Y\prec_kC_p$ for every $p\in N$ such that $\mathrm{Cr}_k(C_p)\cap D_{kj}
\subseteq\mathrm{Cr}_k(Y)$ up to a null set;
\item (E2) if for every $c<\infty$ and every non-null $A\subseteq
D_{kj}\cap\{t<c\}$ there is $b<\infty$ with $\Hi kjb\prec^k_dA$;
\item \Tcut{} if for every $q\in[e_k,e_j]$ with $C_q\ne D_k$ and every
non-null measurable $B\subseteq D_{kj}\cap\{t<\alpha^{(k)}_q\}$, the model
$E\coloneqq C_q+\cred_k(B)$ satisfies $E\preceq_kC_{q'}$ for all $q'$ in some
neighbourhood of $q$.
\end{enumerate}
\end{definition}

\Ts{} is \To{} with the extra credit of the approaching selections removed.
\Tcc{} is \To{} for \emph{credit-contained} approaches: inside $D_{kj}$ the
approaching selections give credit at $k$ only where $Y$ does. (E2) is a
property of $\preceq^k_d$ alone; it is the first statement of
Lemma~\ref{lem:T1otools}(a), from FSE's Lemma~7 (E2). \Tcut{} is the
instance of \Tw{} used in the proof of Theorem~\ref{thm:cut}.

\begin{lemma}\label{lem:Tsrel}
Assume (G) at every world and (S), and let $k\ne j$.
\begin{enumerate}[label=(\alph*),leftmargin=*,itemsep=2pt]
\item Along $[e_k,e_j]$, \To{} implies \Ts{}, \Ts{} implies \Tcc{}, and
\Tw{} implies \Tcut{}. A complete world satisfies \To, hence \Ts, \Tcc{}
and \Tcut, along both of its edges.
\item If world $k$ satisfies (IB) and \Ts{} along $[e_k,e_j]$, it satisfies
(E2) along $[e_k,e_j]$.
\item Let world $k$ satisfy \Tcc{} along $[e_k,e_j]$, and let $q$, $B$ and
$E$ be as in \Tcut. Then $E\prec_kC_{q'}$ for all $q'$ near $q$ with
$\mathrm{Cr}_k(C_{q'})\cap D_{kj}\subseteq\mathrm{Cr}_k(C_q)\cup B$, in
particular for all $q'$ near $q$ with $\alpha^{(k)}_{q'}\ge
\alpha^{(k)}_q$. So under \Tcc{}, \Tcut{} can fail at $q$ only along
selections $C_{q'}$ with $\alpha^{(k)}_{q'}<\alpha^{(k)}_q$ that give credit
at $k$ on a non-null subset of $D_{kj}\cap\{\alpha^{(k)}_{q'}<t^k_j\le
\alpha^{(k)}_q\}\setminus B$.
\item In $\mathrm{Conf}(\theta)$ every world satisfies \Tcc{} and \Tcut{}
along both of its edges. In the K-world every world satisfies \Ts{}, hence \Tcc{} and
(E2), along both of its edges, while world~$1$ violates \To{}
(Proposition~\ref{prop:confT1}). In $\mathrm{Conf}(1_{P_2})$ world~$1$
violates (E2), hence \Ts{}, along $[e_1,e_2]$, and its debt order is not
fine; it satisfies \Tcc.
\end{enumerate}
\end{lemma}
\begin{proof}
(a) $Y\prec_kC_p$ gives $Y\prec_kC_p\preceq_kC_p\langle Y\rangle$. A strict
pointwise $k$-improvement $Y$ of $C_{p^*}$ gives credit at $k$ off
$D_{kj}$, as every selection on the edge does; so if
$\mathrm{Cr}_k(C_p)\cap D_{kj}\subseteq\mathrm{Cr}_k(Y)$, then
$\mathrm{Cr}_k(C_p)\subseteq\mathrm{Cr}_k(Y)$ and $C_p\langle Y\rangle=C_p$.
In \Tcut, $C_q$ has rank $1$ or $2$ at $k$ on $D_{kj}\cap\{t<\alpha_q\}$, so
$E$ is a strict pointwise $k$-improvement of $C_q$. Complete worlds satisfy
\To{} by (T1).

(b) Put $D\coloneqq D_{kj}$, and let $c<\infty$ and $A\subseteq
D\cap\{t<c\}$ be non-null. Let $q_e$ be the far end of the frame $(k,j)$. Then
$C_{q_e}=D_j$, $q_e\ne e_k$ (\ref{P:ends}(b)), and $Y\coloneqq D_j+\sil(A)$
is a strict pointwise $k$-improvement of $D_j$, which gives debt at $k$ on
all of $D$. Let $N$ be the neighbourhood given by \Ts{} for $q_e$ and $Y$.
Choose $q_m\to q_e$ strictly between $e_k$ and $q_e$; by \ref{P:ends}(a),
$\beta_m\coloneqq\beta^{(k)}_{q_m}\to\infty$. Fix $m$ with $q_m\in N$ and
$\beta_m>c$, and let $C\coloneqq C_{q_m}=(\alpha_m,\beta_m)$. As
$\mathrm{Cr}_k(Y)\cap D=\varnothing$, $C\langle Y\rangle$ gives debt at $k$
on $D\cap\{t<\beta_m\}$, silence on the rest of $D$, and credit off $D$. Let
$N_0\coloneqq D_j+\sil(A\cup(D\cap\{t\ge\beta_m\}))$. Up to null sets,
$C\langle Y\rangle=N_0+\debt_k(A)$ and $Y=N_0+\debt_k(\Hi kj{\beta_m})$, and
$N_0$ is silent on $A\cup\Hi kj{\beta_m}$. \Ts{} gives $Y\prec_kC\langle
Y\rangle$, and (IB$'$) gives $\Hi kj{\beta_m}\prec^k_dA$.

(c) $\mathrm{Cr}_k(E)\cap D_{kj}=(D_{kj}\cap\{t>\alpha_q\})\cup B$. Apply
\Tcc{} with $p^*\coloneqq q$ and $Y\coloneqq E$. If $\alpha_{q'}\ge
\alpha_q$, then $\mathrm{Cr}_k(C_{q'})\cap D_{kj}=D_{kj}\cap\{t>
\alpha_{q'}\}\subseteq\mathrm{Cr}_k(E)$. If the inclusion fails, then
$\alpha_{q'}<\alpha_q$ and $C_{q'}$ gives credit on the non-null set
$(D_{kj}\cap\{\alpha_{q'}<t\le\alpha_q\})\setminus B$.

(d) \emph{\Tcc.} Worlds $2$ and $3$ are complete. At world~$1$ on
$[e_1,e_j]$, let $p^*$, $\Dscr^*$, $\rho^*$, $Y$ and $I=I_s\cup I_c$ be as in
the proof of Proposition~\ref{prop:confT1}. If $U_\theta(Y)<
U_\theta(\Dscr^*)$, that proof gives $Y\prec_1\Dscr_p$ for all $p$ near
$p^*$. Otherwise $I_c$ is null, so $\mathrm{Cr}_1(Y)\cap D_{1j}=D_{1j}\cap
\{t>\rho^*\}$ up to a null set; for $p\in[e_1,p^*)$ the selection $\Dscr_p$
gives credit at~$1$ on the non-null band $B(p,p^*)$, so these approaches are
not credit-contained; and for $p$ between $p^*$ and $e_j$,
$Y\prec_1\Dscr^*\preceq_1\Dscr_p$. So \Tcc{} holds.

\emph{\Tcut.} Worlds $2$ and $3$ are complete. At world~$1$ on
$[e_1,e_j]$, $C_q=\Dscr_q=(\rho_q,\rho_q)$, so $B$ lies in the debt region
of $\Dscr_q$; $E$ is silent only on a null set, so
$U_\theta(E)=L_1(E)=L_1(\Dscr_q)-2\int_Bg_1$. As $p\mapsto L_1(\Dscr_p)=
U_\theta(\Dscr_p)$ is continuous, $E\prec_1\Dscr_{q'}$ for all $q'$ near
$q$.

\emph{The K-world.} Worlds $2$ and $3$ are complete. Let $p^*$, $\Dscr^*$,
$\rho^*$, $Y$ and $I=I_s\cup I_c$ be as in the proof of
Proposition~\ref{prop:confT1}. For $p$ between $p^*$ and $e_j$,
$Y\prec_1\Dscr^*\preceq_1\Dscr_p\preceq_1\Dscr_p\langle Y\rangle$. Let
$p\in[e_1,p^*)$ and $B\coloneqq B(p,p^*)$. Up to null sets,
$\mathrm{Cr}_1(\Dscr_p)\setminus\mathrm{Cr}_1(Y)=B\setminus I_c$, so
$\Dscr_p\langle Y\rangle=\Dscr^*+\sil(B\setminus I_c)+\cred_1(B\cap I_c)$,
while $Y=\Dscr^*+\sil(I_s)+\cred_1(I_c)$. For $\theta\equiv1$, silence and
debt cost $2g_1$ in $U_1$, so
\[
 U_1(\Dscr_p\langle Y\rangle)=U_1(\Dscr^*)-2\int_{B\cap I_c}g_1\ \ge\
 U_1(\Dscr^*)-2\int_{I_c}g_1=U_1(Y),
\]
and $L_1(\Dscr_p\langle Y\rangle)-L_1(Y)\ge\int_{I_s}g_1+2\int_{I_c}g_1
-2\int_Bg_1$. As $p\to p^*$, $\int_{B(p,p^*)}g_1\to0$, while
$\int_{I_s}g_1+2\int_{I_c}g_1>0$. So $Y\prec_1\Dscr_p\langle Y\rangle$ for
all $p\in[e_1,p^*)$ near $p^*$: \Ts{} holds, and \Tcc{} and (E2) follow by
(a) and (b).

\emph{$\theta=1_{P_2}$.} Here $\sigma_\theta(V+\debt_1(B))=
\int_{P_2\setminus B}g_1$, so $U_\theta(V+\debt_1(B))=L_1(V)+\int_Bg_1+
\int_{P_2\setminus B}g_1=\mathrm{const}+\int_{B\setminus P_2}g_1$. Hence
$C\preceq^1_dA$ implies $\int_{C\setminus P_2}g_1\le\int_{A\setminus
P_2}g_1$. For non-null $A\subseteq P_2\cap\{t^1_2<1\}$, every
$C\preceq^1_dA$ lies in $P_2$ up to a null set. Every high tail $\Hi12b$
contains the non-null set $P_1\cap\{t^1_2>b\}$, so $\Hi12b\not\preceq^1_dA$,
and (E2) fails; by (b), \Ts{} fails. $H'$ is not a finite union of sets
$\preceq^1_dA$.
\end{proof}

The following lemma replaces the second part of
Lemma~\ref{lem:T1otools}(a), FSE's Lemma~11, whose conclusion (strictly AC
pairs) can fail without completeness (Remark~\ref{rem:conf}(b)).

\begin{lemma}\label{lem:SACo}
Let $v\ne w$, assume (G) at worlds $v$ and $w$, (IB) at $w$, and (S). Use
world $w$'s frame on $[e_w,e_v]$, with $t\coloneqq t^w_v$ on $D\coloneqq
D_{vw}$ and $C_q=(\alpha_q,\beta_q)\coloneqq(\alpha^{(w)}_q,\beta^{(w)}_q)$.
Let $C_{q_0}=D_v$, and let $q_m\to q_0$ along the edge with
$\alpha_{q_m}<\infty$ for all $m$. Then for every $c<\infty$ and every
non-null $A'\subseteq D\cap\{t<c\}$ there is $m_0$ such that for every
$m\ge m_0$ the set $W_m\coloneqq D\cap\{t>\alpha_{q_m}\}$ is non-null and
disjoint from $A'$, and the pair $(A',W_m)$ is not SI at $w$. In particular
no non-null subset of $D\cap\{t<c\}$ is debt-infinitesimal at $w$.
\end{lemma}
\begin{proof}
$D_v$ gives credit at $v$, that is, debt at $w$, on all of $D$; so
\ref{P:ends}(a) in world $w$'s frame gives $\beta_m\coloneqq\beta_{q_m}\to
\infty$. $W_m$ is non-null since $\alpha_{q_m}<\infty$, and once
$\beta_m>c$ it lies in $\{t\ge\beta_m\}$ and misses $A'$. Let $Y\coloneqq
D_v+\sil(A')$; then $Y\prec_wD_v$ by strict (G). Suppose $(A',W_m)$ is SI at
$w$ for infinitely many $m$ with $\beta_m>c$. For such $m$ put
$N_m\coloneqq D_v+\sil(A'\cup(D\cap\{t\ge\beta_m\}))$. Up to null sets,
\[
 C_{q_m}=N_m+\debt_w(A')+\cred_w(W_m),\qquad
 Y=N_m+\debt_w(D\cap\{t\ge\beta_m\}):
\]
both $C_{q_m}$ and $N_m+\debt_w(A')+\cred_w(W_m)$ give debt at $w$ on
$D\cap\{t<\beta_m\}$, silence on $D\cap\{\beta_m\le t\le\alpha_{q_m}\}$,
credit at $w$ on $W_m$, and the option good at $v$ and $w$ off $D$. $N_m$ is
silent on $A'\cup W_m$, so (IB$'$) turns the SI pair into
$C_{q_m}\prec_wN_m$, and $N_m\preceq_wY$ by weak (G). Hence $q_m\in A_w(Y)$
for infinitely many $m$, and, $A_w(Y)$ being closed, $q_0\in A_w(Y)$:
$D_v\preceq_wY$, a contradiction. A debt-infinitesimal $A'$ would make every
$(A',W_m)$ SI.
\end{proof}

The lemma uses closedness of one section, \ref{P:ends}(a), and (G) and (IB)
at $w$: no completeness and no form of (T1). At a complete world it is
weaker than FSE's Lemma~11, which gives strictly AC pairs.

\begin{proposition}\label{prop:audit}
Assume (G) at every world and (S). Let $\{v,w,u\}=\{1,2,3\}$, and let world
$v$ satisfy (IB).
\begin{enumerate}[label=(\alph*),leftmargin=*,itemsep=2pt]
\item If $v$ satisfies (E2) along both of its edges,
Lemma~\ref{lem:T1otools}(b) holds at $v$.
\item If (EP$_{vw}$) holds and $v$ satisfies \Tcc{} along $[e_v,e_w]$, then
CCF's Lemma~5 holds at $v$ on $[e_v,e_w]$, and so does the statement about
$\preceq^v_d$ in Lemma~\ref{lem:T1otools}(c). If (EP) holds on both edges at
$v$ and $v$ satisfies \Tcc{} along both, then the debt order of $v$ is fine:
$H'\in\Jid^v_A$ for every non-null $A$.
\item Lemma~\ref{lem:sCN}(c) holds with \Tcc{} in place of \To.
\item Lemma~\ref{lem:sCN}(e) holds if $v$ satisfies (E2) along both of its
edges. Lemma~\ref{lem:sCN}(f) holds if $w$ also satisfies (IB) and $v$ and
$w$ satisfy (E2) along both of their edges.
\item Theorem~\ref{thm:LH} holds with \Tos{} replaced by: every world
satisfies (E2) and \Tcc{} along both of its edges.
\item Theorem~\ref{thm:cut} holds with \Tcut{} at $v$ along $[e_v,e_w]$ in
place of \Tw. Hence Theorem~\ref{thm:SC} holds with \Tcut{} at world~$3$
along $[e_3,e_1]$, and Proposition~\ref{prop:count2} with \Tcut{} at
world~$2$ along $[e_2,e_1]$, in place of \Tw.
\item Step~3(i) of the proof of Theorem~\ref{thm:rich} holds with \Tcc{} at
world~$1$ along $[e_1,e_2]$ in place of \To.
\end{enumerate}
\end{proposition}
\begin{proof}
(a) By the proof of Lemma~\ref{lem:T1otools}, FSE's Lemma~8(a) uses
completeness only through (E2) on both edges at $v$, and otherwise (G) and
the joint densities on $P_v$.

(b) CCF's Lemma~5(b) uses only (T2). In Lemma~5(a), resp.\ 5(c), the model
$Y=E_w(b)+\sil(A)$, resp.\ $D_w+\sil(A)$, is compared with the selections
$E_w(s)$ that \ref{P:approach} provides in every neighbourhood of the point
$q^*$ with $C_{q^*}=E_w(b)$, resp.\ $D_w$, for all $s$ in an interval
$[s_{n_0},b)$, resp.\ $[s_{n_0},\infty)$. These selections give no credit
at $v$ on $D_{vw}$, so \Tcc{} gives $Y\prec_vE_w(s)$ for these $s$, which is
what CCF derives from (T1). The statement about $\preceq^v_d$ in
Lemma~\ref{lem:T1otools}(c) is the debt half of Step~1 of CCF's proof of
Theorem~7, which uses only CCF's Lemmas~4 and~5, (G), transitivity and a
Lebesgue-number argument; Step~0 uses only (G), and the debt halves of
Steps~2 and~3 use Step~1 on both edges and the joint densities.

(c) In the proof of Lemma~\ref{lem:sCN}(c), $Y=E_w(b)+\cred_v(Z)$ is
compared with the selections $E_w(s)$, $s\in[s_0,b)$, which give no credit
at $v$ on $D_{vw}$.

(d) The proof of Lemma~\ref{lem:sCN}(e) uses Lemma~\ref{lem:T1otools}(b) at
$v$ only. In the proof of Lemma~\ref{lem:sCN}(f), the first step is (e) at
$w$, and Case~II uses Lemma~\ref{lem:T1otools}(b) at $v$. Case~I uses the
second part of Lemma~\ref{lem:T1otools}(a) at $w$; we replace it by
Lemma~\ref{lem:SACo} at $w$. Along the $m$ with $\beta^{(v)}_m>0$ we have
$\alpha^{(w)}_m=1/\beta^{(v)}_m<\infty$, and Lemma~\ref{lem:SACo}, with
$q_0$ and these $q_m$, $c\coloneqq b+1$ and $A'\coloneqq D_{vw}\cap\{b<
t^w_v<b+1\}\subseteq\Hi wvb$, gives $m$ for which $W_m$ is non-null and
disjoint from $A'$ and $(A',W_m)$ is not SI at $w$. But $A'$ is
debt-infinitesimal at $w$, so $(A',W_m)$ is SI: a contradiction.

(e) The proof of Theorem~\ref{thm:LH} uses Lemma~\ref{lem:sCN}(a), (b) and
(d), which do not involve \To, and Lemma~\ref{lem:sCN}(c) and (f), which
hold under the new hypothesis by parts (c) and (d) above.

(f) In the proof of Theorem~\ref{thm:cut}, \Tw{} is applied only to
$M^-=C_q+\cred_v(B)$ with $C_q\ne D_v$ and $B\subseteq D_{vw}\cap\{0<
t^v_w<\alpha_q\}$, which is an instance of \Tcut. In the proofs of
Theorem~\ref{thm:SC} and Proposition~\ref{prop:count2}, \Tw{} enters only
through Theorem~\ref{thm:cut}.

(g) In Step~3(i), $X=E_2(z)+\sil(Z)$ is a strict pointwise $1$-improvement of
$C(s^-)=E_2(z)$, and it is compared only with selections $C(s')$,
$s'\in[s_1,s^-)\subseteq(s_0,s_e)$. There $\alpha(s')=\infty$, so
$C(s')=E_2(\beta(s'))$ gives no credit at~$1$ on $D_{12}$, and these
approaches are credit-contained. \Tcc{} gives a neighbourhood $N$ of
$q(s^-)$ with $X\prec_1C(s')$ whenever $q(s')\in N$, and the rest of Step~3(i)
is unchanged.
\end{proof}

\begin{remark}\label{rem:audit}
\emph{The uses of \To{} in the proof of Theorem~\ref{thm:A}}, as listed in
Remark~\ref{rem:AB}(a), and the weakest hypothesis known to suffice for each:
\begin{enumerate}[label=(\arabic*),leftmargin=*,itemsep=2pt]
\item The first part of Lemma~\ref{lem:T1otools}(a), and through it the
debt yardsticks of Lemma~\ref{lem:T1otools}(b), which enter
Lemma~\ref{lem:sCN}(e),(f), Lemmas~\ref{lem:exchange}--\ref{lem:B3} and
Theorems~\ref{thm:own}--\ref{thm:sweepheavy}: (E2) along both edges, which
\Ts{} implies (Lemma~\ref{lem:Tsrel}(b)).
\item The second part of Lemma~\ref{lem:T1otools}(a) (FSE's Lemma~11), used
only in Case~I of Lemma~\ref{lem:sCN}(f): nothing
(Proposition~\ref{prop:audit}(d)).
\item Lemma~\ref{lem:sCN}(c): \Tcc.
\item Step~3(i) of Theorem~\ref{thm:rich}: \Tcc{}
(Proposition~\ref{prop:audit}(g)). The approaching selections $C(s)$,
$s\in(s_0,s^-)$, have $\alpha(s)=\infty$ and give no credit at~$1$ on
$D_{12}$.
\item CCF's Lemmas~5(a) and~5(c), in Lemma~\ref{lem:T1otools}(c),(d) and
Lemma~\ref{lem:fineregion}: \Tcc.
\item Theorem~\ref{thm:cut} in Proposition~\ref{prop:count2}, which closes a
cut: \Tcut{} at world~$2$ along $[e_2,e_1]$. By Lemma~\ref{lem:Tsrel}(c),
\Tcc{} gives the instances of \Tcut{} approached by selections without extra
credit; whether \Ts{} suffices for \Tcut{} is not known.
\item CCF's Lemmas~6(a) and~6(c), used only for the fineness of the credit
orders in the conclusion of Theorem~\ref{thm:A}: \To{} for these
instances. \Tcc{} and \Ts{} give nothing here. For
$Y=E^w_v(a')+\cred_v(T)$ with $T\subseteq D_{vw}\cap\{t^v_w<a\}$, approached
by $E^w_v(s)$ with $a\le s<a'$, the approach is not credit-contained, and
$E^w_v(s)\langle Y\rangle=E^w_v(a')=C_{p^*}$, which $Y$ beats by strict (G)
in any case; likewise for $Y=V_{vw}+\cred_v(T)$.
\end{enumerate}
No other step of the proof of Theorem~\ref{thm:A} uses \To.
\end{remark}

\begin{theorem}\label{thm:Astar}
Let $n=3$, let every ranking be a preorder with (G) and (IB), and assume
(S). Suppose that every world satisfies (E2) and \Tcc{} along both of its
edges (for instance \Ts), and that some world satisfies \Tcut{} along one of
its edges. If $V$ is undominated, then every edge has the endpoint property
and every world's debt order is fine. If in addition every world satisfies
the instances of \To{} in item~(7) of Remark~\ref{rem:audit}, every world
is fine.
\end{theorem}
\begin{proof}
Suppose $V$ is undominated. The hypotheses other than \Tcut{} are invariant
under relabelling; relabel so that world~$2$ satisfies \Tcut{} along
$[e_2,e_1]$. Let (A$^-$) be (A) with \Tos{} replaced by ``every world
satisfies (E2) and \Tcc{} along both of its edges''. By
Proposition~\ref{prop:audit}(e), Theorem~\ref{thm:LH} holds under (A$^-$).
By Remark~\ref{rem:audit}, Lemmas~\ref{lem:exchange}--\ref{lem:B3},
Proposition~\ref{prop:dichotomy} and
Theorems~\ref{thm:own}--\ref{thm:sweepheavy} use \To{} only through
Lemma~\ref{lem:T1otools}(b) (item~1), the statement about debt in
Lemma~\ref{lem:T1otools}(c) (item~5), Theorem~\ref{thm:LH}, and Step~3(i)
of Theorem~\ref{thm:rich} (item~4); so they hold under (A$^-$), by
Proposition~\ref{prop:audit}(a), (b) and~(g). Proposition~\ref{prop:count2} holds
with \Tcut{} at world~$2$ along $[e_2,e_1]$
(Proposition~\ref{prop:audit}(f)). The proof of Theorem~\ref{thm:A} now
shows that every edge has the endpoint property, and
Proposition~\ref{prop:audit}(b) gives the fineness of the debt orders. With
the instances of item~(7), the credit half of Lemma~\ref{lem:T1otools}(d)
applies as well.
\end{proof}

\begin{theorem}\label{thm:Bstar}
Let $n=3$, let every ranking be a preorder with (G) and (IB), and assume
(S). Suppose that at most one world is incomplete, and that this world
satisfies \Ts{} along both of its edges, or merely (E2) and \Tcc. Then $V$
is dominated.
\end{theorem}
\begin{proof}
By Lemma~\ref{lem:Tsrel}(a),(b), the complete worlds satisfy (E2), \Tcc{}
and \Tcut{} along both of their edges, and \Ts{} implies (E2) and \Tcc{}.
At least two worlds are complete, so Theorem~\ref{thm:Astar} applies. If
$V$ were undominated, every edge would have the endpoint property, and by
\ref{R:CCF} both complete worlds would be (W)-worlds, contrary to
\ref{R:TAW}.
\end{proof}

Theorem~\ref{thm:Bstar} is stronger than Theorem~\ref{thm:B}: world~$1$ of
the K-world satisfies \Ts{} but not \To{} (Lemma~\ref{lem:Tsrel}(d)).

\subsection{Cuts and crossing approaches}\label{sec:crossing}

Let $(v,w)$ be an ordered pair of distinct worlds, and use world $v$'s frame
on $[e_v,e_w]$, with $C_q=(\alpha_q,\beta_q)$. A sequence $q_m\to q^\circ$ in
the edge is \emph{crossing} if $\alpha_{q_m}<\alpha_{q^\circ}$ and
$\beta_{q_m}>\beta_{q^\circ}$ for every $m$, and \emph{eventually crossing}
if this holds for all large $m$. Then $C_{q_m}$ gives more credit at $v$ than
$C_{q^\circ}$ on $D_{vw}\cap\{\alpha_{q_m}<t^v_w<\alpha_{q^\circ}\}$ and more
debt on $D_{vw}\cap\{\beta_{q^\circ}<t^v_w<\beta_{q_m}\}$: the two
selections are not comparable pointwise at $v$. A family $\Bfam$ of
measurable sets is \emph{hereditary} if it contains every measurable subset
of each of its members.

\begin{lemma}\label{lem:reach}
Let $v\ne w$ and assume (G) at worlds $v$ and $w$ and (S). Let $\Bfam$ be a
band-dense family of measurable subsets of $D_{vw}$, $\mathcal R$ the set of
models reachable at world $v$ from $D_v$ by connecting chains with
perturbation sets in $\Bfam$, and $\Sigma\coloneqq\{q\in[e_v,e_w]:C_q\in
\mathcal R\}$.
\begin{enumerate}[label=(\alph*),leftmargin=*,itemsep=2pt]
\item If $D_w\notin\mathcal R$, then $\Sigma$ is relatively open in
$[e_v,e_w]$, contains $e_v$ and does not contain $e_w$.
\item If $q^\circ\in(e_v,e_w)$ and $C_{q^\circ}\notin\mathcal R$, then
$C_{q^\circ}\preceq_wX$ for every $X\in\mathcal R$.
\item If $q^\circ\in(e_v,e_w)$, $F$ is regular with $F\prec_vC_{q^\circ}$,
and $q_m\to q^\circ$, then $C_{q_m}\preceq_wF$ for all large $m$.
\end{enumerate}
\end{lemma}
\begin{proof}
(a) The paragraph ``$\Sigma$ is open'' of the proof of
Theorem~\ref{thm:cut} uses only strict (G) at $w$, saturation, down-closure
and (T2); $C_{e_v}=D_v\in\mathcal R$ and $C_{e_w}=D_w\notin\mathcal R$.
(b) If $C_{q^\circ}\not\preceq_wX$, competitiveness on $\{v,w\}$ gives
$C_{q^\circ}\preceq_vX$, and down-closure gives $C_{q^\circ}\in\mathcal R$.
(c) This is (T2) with $v$ and $w$ exchanged; $q_m$ lies in the open edge near
$q^\circ$ for large $m$.
\end{proof}

\begin{proposition}\label{prop:bcount}
Assume (G) at every world and (S), let world~$2$ be a preorder with (IB),
and suppose the far end of the frame $(3,2)$ is rich. Let $T_0$ be non-null,
let $\Bfam$ be the family of measurable subsets of the sets $J_n$
constructed from $T_0$ in the proof of Proposition~\ref{prop:count2}, and
let $\mathcal R$ be the set of models reachable at world~$2$ from $D_2$ by
connecting chains with perturbation sets in $\Bfam$. For $0<M<\infty$ put
\[
 \mathcal G_M(T_0)\coloneqq\{\Uf:\ P_1\cap\{t^2_1<M\}\in\Uf\text{ and }
 T_0\preceq^2_cP\text{ for every }P\in\Uf\}.
\]
If $\mathcal R$ contains a model $X$ with debt at~$2$ on
$D_{21}\cap\{t^2_1<M\}$, then $\mathcal G_M(T_0)$ is finite. The
construction of $\Bfam$ does not use \Tw, and $\Bfam$ is hereditary,
band-dense for $[e_2,e_1]$, and consists of subsets of $P_2$.
\end{proposition}
\begin{proof}
The construction in the proof of Proposition~\ref{prop:count2} uses
\ref{R:yardc}, Lemma~\ref{lem:near} and Lemma~\ref{lem:light}(a), and no
form of (T1). $X$ has debt at~$2$ on a non-null set, so $D_2\prec_2X$ by
strict (G), and a connecting chain from $D_2$ to $X$ with perturbation sets
$B_1,\dots,B_{\ell-1}\in\Bfam$ has length $\ell\ge2$. Choose $N$, $F$ and
$T_*$ as in the proof of Proposition~\ref{prop:count2}: every $B_i$ is
$T_*$-light at~$2$, $F\subseteq P_2$, $T_*\subseteq P_3$ and
$T_*\preceq^2_cT_0$. Suppose $\mathcal G_M(T_0)$ has $\ell$ distinct members.
As in the proof of Proposition~\ref{prop:count3}, there are pairwise disjoint
$K_l$ in these ultrafilters; intersecting with $P_1\cap\{t^2_1<M\}$, we may
take $K_l\subseteq P_1\cap\{t^2_1<M\}\subseteq\{r_2(\cdot,X)=2\}\setminus
F$. They are non-null and $T_*\preceq^2_cT_0\preceq^2_cK_l$, so
Lemma~\ref{lem:light}(c) gives $D_2+\debt_2(B_i)\prec_2D_2+\debt_2(K_l)$ for
all $i,l$. This contradicts Remark~\ref{rem:bounded}. So
$|\mathcal G_M(T_0)|\le\ell-1$.
\end{proof}

\begin{proposition}\label{prop:nocross}
Let $(v,w)$ be an ordered pair of distinct worlds. Assume (G) at every
world, (S), and (IB) at $v$ and $w$. Let world $v$ satisfy \Tcc{} along
$[e_v,e_w]$, let world $w$ satisfy (E2) along both of its edges and
(LH$^+_w$), and let the far end of the frame $(w,v)$ be rich. Let $\Bfam$ be
a hereditary band-dense family of measurable subsets of $P_v$, and let
$\mathcal R$ and $\Sigma$ be as in Lemma~\ref{lem:reach}. Let $q^\circ$ be a
point of the open edge $(e_v,e_w)$ with $C_{q^\circ}\notin\mathcal R$ and
$C_{q^\circ}\ne D_w$, and let $q_m\in\Sigma$ converge to $q^\circ$. Then
$(q_m)$ is eventually crossing, and $\beta_{q^\circ}<\alpha_{q^\circ}$.
\end{proposition}
\begin{proof}
Put $t\coloneqq t^v_w$, $D\coloneqq D_{vw}$,
$(\alpha^\circ,\beta^\circ)\coloneqq C_{q^\circ}$ and
$(\alpha_m,\beta_m)\coloneqq C_{q_m}$. Since $D_v\in\mathcal R$,
$C_{q^\circ}\ne D_v$ and $\alpha^\circ>0$. On $D$ the world-$w$ rank of a
regular model is $2$ minus its world-$v$ rank \eqref{eq:swap}.

(i) \emph{$\alpha_m<\alpha^\circ$ for all large $m$.} Otherwise
$\alpha_m\ge\alpha^\circ$ along a subsequence. Take a non-null $B\in\Bfam$
in the band $D\cap\{0<t<\alpha^\circ\}$ and $E\coloneqq
C_{q^\circ}+\cred_v(B)$, a strict pointwise $v$-improvement of
$C_{q^\circ}$. Along the subsequence $\mathrm{Cr}_v(C_{q_m})\cap D=
D\cap\{t>\alpha_m\}\subseteq\mathrm{Cr}_v(E)$, so \Tcc{} gives
$E\prec_vC_{q_m}$ for large $m$ in it. As $\mathcal R$ is down-closed,
$E\in\mathcal R$, and its $\Bfam$-neighbour $C_{q^\circ}$ lies in
$\mathcal R$: a contradiction.

Suppose now that $\beta_m\le\beta^\circ$ for infinitely many $m$, and pass
to these $m$. We derive a contradiction, which shows that $(q_m)$ is
eventually crossing.

(ii) \emph{$\alpha_m\to\alpha^\circ$.} For $a<\alpha^\circ$ let $X_a\coloneqq
C_{q^\circ}+\cred_v(D\cap\{a<t<\alpha^\circ\})$; $X_a\prec_vC_{q^\circ}$ by
strict (G). If $\alpha_m\le a$ for infinitely many $m$, then for these $m$,
$C_{q_m}\le^v_\pw X_a$: on $D\cap\{t<\beta_m\}$ both have rank $2$, since
$\beta_m\le\min(\beta^\circ,a)$; on $D\cap\{\beta_m\le t\le\alpha_m\}$,
$C_{q_m}$ has rank $1$ and $X_a$ rank at least $1$; elsewhere $C_{q_m}$ has
rank $0$. So $q_m\in A_v(X_a)$, and closedness gives $C_{q^\circ}
\preceq_vX_a$: a contradiction. With (i), $\alpha_m<\alpha^\circ$ and
$\alpha_m\to\alpha^\circ$.

(iii) \emph{For $a'<\alpha^\circ$, $Y_{a'}\coloneqq C_{q^\circ}+
\cred_v(D\cap\{a'<t<\alpha^\circ\})$ lies in $\mathcal R$.} $Y_{a'}$ is a
strict pointwise $v$-improvement of $C_{q^\circ}$. By (ii), $\alpha_m>a'$
for large $m$, and then $\mathrm{Cr}_v(C_{q_m})\cap D\subseteq\mathrm{Cr}_v
(Y_{a'})$ up to a null set. \Tcc{} gives $Y_{a'}\prec_vC_{q_m}$ for large
$m$, and $\mathcal R$ is down-closed.

\emph{Case 1: $\beta^\circ<\alpha^\circ$.} Fix
$\beta^\circ<a'''<a''<\alpha^\circ$ and a non-null $B'_0\in\Bfam$ in the band
$D\cap\{\beta^\circ<t<a''\}$. By (LH$^+_w$), $B'_0$ is not strongly
credit-null at $w$: there is a weakly SI pair $(A,T)$ at $w$ with $A$
non-null and $T\subseteq B'_0$. $T$ is non-null, since $(A,\varnothing)$ is
strictly AC; $T\in\Bfam$, and $T\subseteq P_v$. By
Proposition~\ref{prop:audit}(a), Lemma~\ref{lem:T1otools}(b) holds at $w$;
it gives a non-null debt yardstick $Z\subseteq P_w\cap\{t^w_v\in I\}$,
$I\coloneqq(1/a'',1/a''')$, with $Z\preceq^w_dA$. As $t^w_v=1/t$ on $D$,
$Z\subseteq P_w\cap\{a'''<t<a''\}$. Split $Z$ into disjoint non-null $Z_1,
Z_2$. $T$ misses $Z\cup A$, so
\[
 V+\debt_w(Z_1)+\cred_w(T)\prec_wV+\debt_w(Z)+\cred_w(T)\preceq_w
 V+\debt_w(A)+\cred_w(T)\preceq_wV,
\]
by strict (G), by $Z\preceq^w_dA$ in the context $V+\cred_w(T)$, and because
$(A,T)$ is weakly SI. So $(Z_1,T)$ is SI at $w$.

Let $F\coloneqq C_{q^\circ}+\cred_v(Z_1)$; $Z_1$ lies in the silence of
$C_{q^\circ}$, so $F\prec_vC_{q^\circ}$. By Lemma~\ref{lem:reach}(c) and
(ii), fix $m$ with $C_{q_m}\preceq_wF$, $\alpha_m>a''$ and
$\beta_m\le\beta^\circ$, and put $G_m\coloneqq D\cap\{\alpha_m<t<
\alpha^\circ\}$. $C_{q_m}$ and $F$ agree except on $D\cap\{\beta_m<t<
\beta^\circ\}$ (world-$w$ ranks $1$ and $0$), on $Z_1$ ($1$ and $2$) and on
$G_m$ ($2$ and $1$). Let $C'$ be $C_{q_m}$ with debt at $v$ on
$D\cap\{\beta_m<t<\beta^\circ\}$. Then $C'\le^w_\pw C_{q_m}$, so
$C'\preceq_wC_{q_m}\preceq_wF$, and $C',F$ differ exactly on $Z_1$ and
$G_m$, with world-$w$ ranks $(1,2)$ and $(2,1)$. By (IB$'$),
\[
 V+\debt_w(G_m)\preceq_wV+\debt_w(Z_1).\tag{P}
\]
By (iii) with $a'\coloneqq\alpha_m$, $Y\coloneqq Y_{\alpha_m}\in\mathcal R$.
$T\subseteq D\cap\{\beta^\circ<t<a''\}$ lies in the silence of $Y$, as
$a''<\alpha_m$; so $X\coloneqq Y+\debt_v(T)$ is a $\Bfam$-neighbour of $Y$
and lies in $\mathcal R$. By Lemma~\ref{lem:reach}(b),
$C_{q^\circ}\preceq_wX$. These models differ exactly on $G_m$
(world-$w$ ranks $1$ and $2$) and on $T$ ($1$ and $0$); by (IB$'$),
\[
 V\preceq_wV+\debt_w(G_m)+\cred_w(T).\tag{Q}
\]
The sets $G_m$, $Z_1$ and $T$ are pairwise disjoint. (P) in the context
$V+\cred_w(T)$ and (Q) give $V\preceq_wV+\debt_w(G_m)+\cred_w(T)
\preceq_wV+\debt_w(Z_1)+\cred_w(T)\prec_wV$: a contradiction.

\emph{Case 2: $\beta^\circ=\alpha^\circ\eqqcolon\gamma$.} Then
$0<\gamma<\infty$, since $C_{q^\circ}\notin\{D_v,D_w\}$, and
$\beta_m\le\alpha_m<\gamma$ for large $m$ by (i). First, $\beta_m\to\gamma$:
for $b<\gamma$ the model $X_b\coloneqq C_{q^\circ}+\sil(D\cap\{b<t<\gamma\})$
satisfies $X_b\prec_vC_{q^\circ}$; if $\beta_m\le b$ for infinitely many
$m$, then for these $m$, $C_{q_m}\le^v_\pw X_b$ (both have rank $2$ on
$\{t<\beta_m\}$, $C_{q_m}$ has rank $1$ on $\{\beta_m\le t\le\alpha_m\}$,
where $X_b$ has rank at least $1$, and rank $0$ elsewhere), and closedness
of $A_v(X_b)$ gives $C_{q^\circ}\preceq_vX_b$, a contradiction. So
$\beta_m\to\gamma$ and $\alpha_m\to\gamma$.

Choose a non-null $B'\in\Bfam$ in the band $D\cap\{\gamma<t<\gamma+1\}$,
where $C_{q^\circ}$ gives credit at $v$, and a non-null $B'_1\subseteq B'$
with $B'\setminus B'_1$ non-null; then $B'_1\subseteq D\cap\{t^w_v<
1/\gamma\}$. Let $q_e$ be the far end of the frame $(w,v)$, and let
$q'_n\to q_e$, strictly between $e_w$ and $q_e$, with
$\alpha^{(w)}_{q'_n}<\infty$. As $C_{q_e}=D_v$, \ref{P:ends}(a) in world
$w$'s frame gives $\beta^{(w)}_{q'_n}\to\infty$, so
$\alpha^{(w)}_{q'_n}\to\infty$. Fix $b''\in(0,\gamma)$. By
Lemma~\ref{lem:SACo} at $w$, with $q_0\coloneqq q_e$, the sequence $(q'_n)$,
$c\coloneqq1/\gamma$ and $A'\coloneqq B'_1$, there is $n$ such that
$B\coloneqq D\cap\{t^w_v>\alpha^{(w)}_{q'_n}\}$ is non-null, disjoint from
$B'_1$, contained in $D\cap\{t<b''\}$, and $(B'_1,B)$ is not SI at $w$.

Let $F\coloneqq C_{q^\circ}+\sil(B)$; $B$ lies where $C_{q^\circ}$ gives debt
at $v$, so $F\prec_vC_{q^\circ}$. By Lemma~\ref{lem:reach}(c), fix $m$ with
$C_{q_m}\preceq_wF$ and $\beta_m>b''$, and put $G_m\coloneqq D\cap
\{\alpha_m<t<\gamma\}$. $C_{q_m}$ and $F$ agree except on $B$ (world-$w$
ranks $0$ and $1$), on $D\cap\{\beta_m<t<\alpha_m\}$ ($1$ and $0$) and on
$G_m$ ($2$ and $0$). Let $C'$ be $C_{q_m}$ with debt at $v$ on
$D\cap\{\beta_m<t<\alpha_m\}$; then $C'\preceq_wC_{q_m}\preceq_wF$, and
$C',F$ differ exactly on $G_m$ and $B$, with world-$w$ ranks $(2,0)$ and
$(0,1)$. By (IB$'$),
\[
 V+\debt_w(G_m)+\cred_w(B)\preceq_wV+\cred_w(G_m).\tag{P$'$}
\]
By (iii), $Y_{\alpha_m}\in\mathcal R$, and so is its $\Bfam$-neighbour
$X\coloneqq Y_{\alpha_m}+\sil(B')$. By Lemma~\ref{lem:reach}(b),
$C_{q^\circ}\preceq_wX$; they differ exactly on $G_m$ (world-$w$ ranks
$0$ and $2$) and on $B'$ ($2$ and $1$). By (IB$'$),
\[
 V+\cred_w(G_m)+\debt_w(B')\preceq_wV+\debt_w(G_m).\tag{Q$'$}
\]
The sets $G_m$, $B$ and $B'$ are pairwise disjoint ($B\subseteq\{t<
b''\}$, $b''<\beta_m\le\alpha_m$, $B'\subseteq\{t>\gamma\}$). (Q$'$) in the
context $\cred_w(B)$ and (P$'$) give
\begin{align*}
 &V+\cred_w(G_m)+\debt_w(B')+\cred_w(B)\\
 &\qquad\preceq_wV+\debt_w(G_m)+\cred_w(B)\preceq_wV+\cred_w(G_m).
\end{align*}
The outer models differ exactly on $B'$ (world-$w$ ranks $2$ and $1$) and on
$B$ ($0$ and $1$), so (IB$'$) gives $V+\debt_w(B')+\cred_w(B)\preceq_wV$,
and strict (G) on $B'\setminus B'_1$ makes $(B'_1,B)$ SI at $w$: a
contradiction.

In Case~2 every sequence is non-crossing, so Case~2 cannot occur, and
$\beta^\circ<\alpha^\circ$.
\end{proof}

\begin{theorem}\label{thm:crossing}
Let $n=3$, let every ranking be a preorder with (G) and (IB), assume (S),
suppose that every world satisfies (E2) and \Tcc{} along both of its edges,
and suppose $V$ is undominated. If some edge lacks the endpoint property,
then for every ordered pair $(v,w)$ of distinct worlds there are a
hereditary band-dense family $\Bfam$ of measurable subsets of $P_v$, with
reachable set $\mathcal R$ at world $v$ as in Lemma~\ref{lem:reach}, and a
point $q^\circ$ of the open edge $(e_v,e_w)$ such that:
\begin{enumerate}[label=(\roman*),leftmargin=*,itemsep=2pt]
\item $C_q\in\mathcal R$ for every $q\in[e_v,q^\circ)$, and
$C_{q^\circ}\notin\mathcal R$ (a \emph{cut});
\item $C_{q^\circ}\notin\{D_v,D_w\}$ and
$\beta^{(v)}_{q^\circ}<\alpha^{(v)}_{q^\circ}$;
\item every sequence in $[e_v,q^\circ)$ that converges to $q^\circ$ is
eventually crossing.
\end{enumerate}
Consequently, if for some ordered pair $(v,w)$ no point of $(e_v,e_w)$ is
the limit of an eventually crossing sequence of points between $e_v$ and
that point (for instance, if any two selections on $[e_v,e_w]$ are
comparable pointwise at $v$), then every edge has the endpoint property and
every world's debt order is fine.
\end{theorem}
\begin{proof}
Let (A$^-$) be as in the proof of Theorem~\ref{thm:Astar}; it holds. By
Proposition~\ref{prop:dichotomy}, every edge carries a two-sided selection;
by Theorem~\ref{thm:rich}, every far end is rich; and by
Theorem~\ref{thm:LH}, every world satisfies (LH$^+$) and (DH). All of these
hold under (A$^-$) (proof of Theorem~\ref{thm:Astar}). Since (A$^-$) is
invariant under relabelling, it suffices to treat $(v,w)=(2,1)$, with third
world~$3$.

Let $E$ be a two-sided selection on $(e_1,e_2)$. Then $B_0(E)=P_1\cap
\{t^2_1>1/\beta^{(1)}_E\}$; fix $M_0>1/\beta^{(1)}_E$, so that
$Q_0\coloneqq B_0(E)\cap\{t^2_1<M_0\}$ is non-null. For integers $k\ge1$ let
$L_k\coloneqq\Lo23{1/k}$, let $\Bfam_k$ be the family of
Proposition~\ref{prop:bcount} for $T_0\coloneqq L_k$, and $\mathcal R_k$ its
reachable set. Suppose that every $\mathcal R_k$ contains a model with debt
at~$2$ on $D_{21}\cap\{t^2_1<M_0\}$. By Proposition~\ref{prop:bcount} every
$\mathcal G_{M_0}(L_k)$ is finite, so $\mathcal G'\coloneqq\bigcup_k
\mathcal G_{M_0}(L_k)$ is countable; list it as $\Uf^1,\Uf^2,\dots$. By
\ref{U:small} choose $B^i\in\Uf^i$ with $\mu(B^i)<2^{-i-1}\mu(Q_0)$, and
put $Q\coloneqq Q_0\setminus\bigcup_iB^i$, a non-null subset of $B_0(E)$.
Theorem~\ref{thm:sweepheavy} gives a non-null $S$ and an ultrafilter $\Uf$
containing $Q$, hence $P_1\cap\{t^2_1<M_0\}$, with $S\preceq^2_cP$ for every
$P\in\Uf$. By \ref{R:yardc} at world~$2$, $L_k\preceq^2_cS$ for large $k$,
so $\Uf\in\mathcal G_{M_0}(S)\subseteq\mathcal G_{M_0}(L_k)\subseteq
\mathcal G'$; but $\Uf$ contains the disjoint sets $Q$ and $B^i$ for the
index $i$ with $\Uf=\Uf^i$. This contradiction shows that for some $k$ no
model in $\mathcal R\coloneqq\mathcal R_k$ has debt at~$2$ on
$D_{21}\cap\{t^2_1<M_0\}$; put $\Bfam\coloneqq\Bfam_k$.

In particular $D_1\notin\mathcal R$, and every $q\in\Sigma$ has
$\beta^{(2)}_q<M_0$, since $C_q$ gives debt at~$2$ on $D_{21}\cap\{t^2_1<
\beta^{(2)}_q\}$. By Lemma~\ref{lem:reach}(a), $\Sigma$ is relatively open,
$e_2\in\Sigma$ and $e_1\notin\Sigma$. Let $q^\circ$ be the point of the
closed set $[e_2,e_1]\setminus\Sigma$ nearest to $e_2$. Then
$[e_2,q^\circ)\subseteq\Sigma$ and $q^\circ\ne e_2$. If $C_{q^\circ}=D_1$,
\ref{P:ends}(a) would give $\beta^{(2)}\to\infty$ along a sequence in
$[e_2,q^\circ)$ converging to $q^\circ$, contrary to $\beta^{(2)}<M_0$ on
$\Sigma$. So $C_{q^\circ}\ne D_1$, hence $q^\circ\ne e_1$; and
$C_{q^\circ}\notin\mathcal R\ni D_2$. So $q^\circ$ lies in the open edge, and
(i) and the first part of (ii) hold. Proposition~\ref{prop:nocross} at
$(v,w)=(2,1)$ applies: \Tcc{} at~$2$ along $[e_2,e_1]$, (E2) at~$1$ along
both edges, (LH$^+_1$), the rich far end of the frame $(1,2)$, and $\Bfam$
hereditary, band-dense and contained in $P_2$
(Proposition~\ref{prop:bcount}). It gives (iii) and
$\beta^{(2)}_{q^\circ}<\alpha^{(2)}_{q^\circ}$.

For the last statement: crossing sequences consist of selections that are
not pointwise comparable at $v$ with the limit selection; so the first part
of the theorem cannot hold for that pair, and every edge has the endpoint
property. Proposition~\ref{prop:audit}(b) gives the fineness of the debt
orders.
\end{proof}

\subsection{Two worlds that violate (NC)}\label{sec:twoNC}

\begin{theorem}\label{thm:twoNC}
Let $n=3$, let every ranking be a preorder with (G) and (IB), assume (S),
and suppose $V$ is undominated. Suppose that every edge has the endpoint
property and that every world satisfies \Tcc{} along both of its edges.
Then every world's debt order is fine, and every world has an SI pair with
non-null debt part and a weakly AC pair with non-null credit part. At most
one world $x$ whose credit order is fine satisfies (NC$_{\mathfrak W_x\cup
\mathfrak F_x}$) for some admissible family $\mathfrak F_x$; in particular
at most one such world satisfies (NC). So if every world's credit order is
fine, at least two worlds $x$ violate (NC$_{\mathfrak W_x\cup\mathfrak
F_x}$) for every admissible $\mathfrak F_x$, and hence violate (NC) and
(NC$^+$).
\end{theorem}
\begin{proof}
The debt orders are fine by Proposition~\ref{prop:audit}(b). Let $x$ be a
world and $[e_x,e_y]$ one of its edges, and use world $x$'s frame, with
$t\coloneqq t^x_y$ and $D\coloneqq D_{xy}$.

\emph{SI pair.} By Proposition~\ref{prop:audit}(c), the proof of
Lemma~\ref{lem:sCN}(c) gives an SI pair at $x$ with non-null debt part
$D\cap\{s\le t<b\}$.

\emph{Weakly AC pair} (CCF, Lemma~11, second half). By \ref{P:approach}
in the parametrisation of \ref{P:param}, with $s^*=-1$ (so that
$M_{s^*}=E^y_x(1)$) and $s_n\downarrow-1$, there is $q^*$ with
$C_{q^*}=E^y_x(1)$, in the open edge, such that every neighbourhood of $q^*$
contains points $q$ with $C_q=E^y_x(a)$ for all $a\in(1,a_0]$, for some
$a_0>1$ depending on the neighbourhood. Let $A\subseteq D\cap\{t<1\}$ be
non-null and $F\coloneqq E^y_x(1)+\debt_x(A)$. $E^y_x(1)$ is silent on $A$,
and debt at $x$ is credit at $y$ there, so $F\prec_yE^y_x(1)$ by strict (G).
By (T2), $C_q\preceq_xF$ for all $q$ in the open edge near $q^*$; so
$E^y_x(a)\preceq_xF$ for some $a>1$. $E^y_x(a)$ and $F$ differ exactly on $A$ (silence
and debt at $x$) and on $D\cap\{1<t\le a\}$ (silence and credit), so (IB$'$)
gives $V\preceq_xV+\debt_x(A)+\cred_x(D\cap\{1<t\le a\})$.

So a world whose credit order is fine is fine and has admissible families.
By Corollary~\ref{cor:NCext}(iii), such a world $x$ that satisfies
(NC$_{\mathfrak W_x\cup\mathfrak F_x}$) for some admissible
$\mathfrak F_x$ is a (W)-world, and by \ref{R:TAW} at most one world is a
(W)-world. Since $\mathfrak W_x\cup\mathfrak F_x\subseteq\mathcal A_x$,
(NC) implies (NC$_{\mathfrak W_x\cup\mathfrak F_x}$), and (NC$^+$) implies
(NC).
\end{proof}

\begin{corollary}\label{cor:twoNC}
Let $n=3$, let every ranking be a preorder with (G) and (IB), and assume (S)
and \Tos. If $V$ is undominated, then every edge has the endpoint property,
every world is fine and has an SI pair with non-null debt part and a weakly
AC pair with non-null credit part, and at least two worlds $x$ violate
(NC$_{\mathfrak W_x\cup\mathfrak F_x}$) for every admissible family
$\mathfrak F_x$. Equivalently, if two worlds satisfy (NC), then $V$ is
dominated; this applies to any two worlds of the kinds (i)--(iv) of
Corollary~\ref{cor:NCclasses}.
\end{corollary}
\begin{proof}
By Theorem~\ref{thm:A}, every edge has the endpoint property and every
world is fine. By Lemma~\ref{lem:Tsrel}(a), \To{} implies \Tcc. So
Theorem~\ref{thm:twoNC} applies. Since $\mathfrak W_x\cup\mathfrak
F_x\subseteq\mathcal A_x$, a world that satisfies (NC) satisfies the
condition for every admissible family. If $V$ were undominated, every world
would be fine with the two pairs, and the worlds of
Corollary~\ref{cor:NCclasses}, cases (i)--(iv), would satisfy (NC$^+$).
\end{proof}

By Theorem~\ref{thm:SCext} and Theorem~\ref{thm:NC}(c), a world with (IB),
(SC) and the endpoint property on both of its edges satisfies (NC$^+$); and
world~$1$ of the K-world satisfies (NC$^+$) but not (SC)
(Remark~\ref{rem:conf}(b)). So the conclusion of
Corollary~\ref{cor:twoNC}, two worlds violating (NC), is stronger than that
of Corollary~\ref{cor:C}, two worlds violating (SC). Without fineness of the
credit orders Theorem~\ref{thm:NC}(b) does not apply: under (NC) a charge
$\ell\ge0$ on $\mathcal A_x$ still exists, but neither layer is shown to be
strictly positive (if $\nu_2(H')=0$, the weakly AC pair only gives
$\nu_1(Z')=0$). So at a world whose debt order is fine and whose credit
order may not be, (NC) is not known to make it a (W)-world.

\section{The shape of a configuration with \texorpdfstring{$V$}{V} undominated}\label{sec:shape}

\subsection{Failures of edge-openness}\label{sec:Ra}

\begin{lemma}\label{lem:Ra}
Assume (G) at worlds $k$ and $j$ and (S). Let $p^*\in[e_k,e_j]$, let $Y$ be
a strict pointwise $k$-improvement of $C_{p^*}$, and let $p_m\to p^*$ along
the edge with $Y\not\prec_kC_{p_m}$ for every $m$. Then $p_m$ lies in the
open edge for large $m$, and for all large $m$:
\begin{enumerate}[label=(\alph*),leftmargin=*,itemsep=1pt]
\item $Y\parallel_kC_{p_m}$;
\item $C_{p_m}\preceq_jY$;
\item $C_{p^*}\not\preceq_kC_{p_m}$; in particular $C_{p_m}$ is not
pointwise worse than $C_{p^*}$ at $k$.
\end{enumerate}
In particular world $k$ is incomplete.
\end{lemma}
\begin{proof}
$Y\prec_kC_{p^*}$ by strict (G), so $p_m\ne p^*$; as $p_m\to p^*$, $p_m$ is
an endpoint of the edge for at most finitely many $m$. Since $p^*\notin
A_k(Y)$ and $A_k(Y)$ is closed, $C_{p_m}\not\preceq_kY$ for large $m$;
together with $Y\not\prec_kC_{p_m}$ this gives $Y\not\preceq_kC_{p_m}$, which
is (a). (b) is the competitiveness of $C_{p_m}$ on $\{k,j\}$. (c): if
$C_{p^*}\preceq_kC_{p_m}$, then $Y\prec_kC_{p^*}\preceq_kC_{p_m}$ gives
$Y\prec_kC_{p_m}$; and $C_{p^*}\le^k_\pw C_{p_m}$ would give
$C_{p^*}\preceq_kC_{p_m}$ by weak (G).
\end{proof}

So a failure of \To{} at world $k$ consists of a strict local improvement
$Y$ of a selected model that becomes incomparable, at $k$, with the nearby
selections, while those selections are weakly better than $Y$ at the
partner world $j$. In the proofs above \To{} is applied to the following
$Y$ (Remark~\ref{rem:AB}(a)): $C_q+\cred_v(B)$ at the boundary of a
reachable family (Theorem~\ref{thm:cut}); $D_w$ with its debt on a bounded
set silenced, at the far end of a frame (Lemma~\ref{lem:T1otools}(a),(b));
the endpoint-family models of CCF's Lemmas~5 and~6 and of
Lemma~\ref{lem:sCN}(c); and $E_2(z)$ with its debt on a yardstick silenced
(Theorem~\ref{thm:rich}). In each case a small change near a selected model
would have to be incomparable with a fixed strict improvement.
Remark~\ref{rem:audit} shows that, except at the far ends, where only the
conclusion (E2) is used, at the cut of Theorem~\ref{thm:cut}, and in CCF's
Lemma~6, the approaching selections give no credit at $k$ beyond that of
$Y$, so that (E2) and \Tcc{} suffice there; and
Proposition~\ref{prop:confT1} shows that (S) allows failures of \To.

\subsection{The shape}\label{sec:C}

\begin{corollary}\label{cor:C}
Let $n=3$, let every ranking be a preorder with (G) and (IB), assume (S),
and suppose $V$ is undominated. Then at least one of the following holds.
\begin{enumerate}[label=(R-\alph*),leftmargin=*,itemsep=2pt]
\item Some world $k$ violates \To{} along one of its edges $[e_k,e_j]$:
there are $p^*$ in the edge, a strict pointwise $k$-improvement $Y$ of
$C_{p^*}$, and points $p_m\to p^*$ of the open edge with
$Y\parallel_kC_{p_m}$ and $C_{p_m}\preceq_jY$ for every $m$. World $k$ is
incomplete.
\item \Tos{} holds, every edge has the endpoint property, every world is
fine, and at least two worlds are not (W)-worlds; each of them violates
(SC), is incomplete, and by EDL's Theorem~2 carries a leverage sequence.
\end{enumerate}
\end{corollary}
\begin{proof}
If (R-a) fails, Lemma~\ref{lem:Ra} shows that \Tos{} holds. By
Theorem~\ref{thm:A}, every edge has the endpoint property and every world is
fine. By Theorem~\ref{thm:SCext}, every world satisfying (SC) is a
(W)-world, and by \ref{R:TAW} at most one world is a (W)-world. So at least
two worlds are not (W)-worlds; they violate (SC), and are therefore
incomplete. The last statement is EDL's Theorem~2.
\end{proof}

Equivalently: $V$ is dominated in every configuration with $n=3$, (G), (IB)
and (S) in which \To{} holds at every incomplete world and every world with
(IB), \To{} and the endpoint property on both of its edges is a (W)-world.
The second condition is the question (EXT) below. Compared with the
three-world note's Corollary~13, which says that some world is incomplete,
Corollary~\ref{cor:C} locates what incompleteness has to achieve: either
incomparabilities between strict local improvements and nearby selections,
or two incomplete worlds whose debt order, credit order or SI/AC dichotomy
is incomplete, in a configuration with the endpoint property on every edge.

Section~\ref{sec:sharp} refines both alternatives.

\begin{corollary}\label{cor:refined}
Let $n=3$, let every ranking be a preorder with (G) and (IB), assume (S),
and suppose $V$ is undominated. Then at least one of the following holds.
\begin{enumerate}[label=(R-\alph*),leftmargin=*,itemsep=2pt]
\item[(R-a1)] Some incomplete world violates (E2) or \Tcc{} along one of
its edges (so it violates \Ts).
\item[(R-a2$''$)] Every world satisfies (E2) and \Tcc{} along both of its
edges; every edge carries a two-sided selection, every far end is rich, and
every world satisfies (LH$^+$) and (DH); for every ordered pair $(v,w)$ of
distinct worlds there is a cut $q^\circ$ as in Theorem~\ref{thm:crossing},
approached only by eventually crossing sequences, at which \Tcut{} fails at
world $v$ along $[e_v,e_w]$; and all three worlds are incomplete.
\item[(R-b$'$)] Every world satisfies (E2) and \Tcc{} along both of its
edges, every edge has the endpoint property, every world's debt order is
fine, and at least two worlds are not (W)-worlds; each of them violates
(SC), is incomplete, and carries a leverage sequence. At most one world
whose credit order is fine satisfies (NC); if every world's credit order is
fine, at least two worlds $x$ violate (NC$_{\mathfrak W_x\cup\mathfrak
F_x}$) for every admissible $\mathfrak F_x$, and in particular (NC).
\end{enumerate}
\end{corollary}
\begin{proof}
Complete worlds satisfy (E2) and \Tcc{} along both of their edges
(Lemma~\ref{lem:Tsrel}(a),(b)), and \Ts{} implies both. Suppose (R-a1)
fails; then every world satisfies (E2) and \Tcc{} along both of its edges.

If every edge has the endpoint property, the debt orders are fine by
Proposition~\ref{prop:audit}(b). By Theorem~\ref{thm:SCext} every world
with (SC) is a (W)-world, and by \ref{R:TAW} at most one world is a
(W)-world; so at least two worlds are not (W)-worlds, and each violates (SC)
and is therefore incomplete. EDL's Theorem~2 gives the leverage sequences,
and the last statement is Theorem~\ref{thm:twoNC}. This is (R-b$'$).

Otherwise Theorem~\ref{thm:crossing} gives, for every ordered pair
$(v,w)$, a family $\Bfam$ and a cut $q^\circ$ with (i)--(iii), and its
proof gives the two-sided selections, the rich far ends, (LH$^+$) and (DH).
Suppose \Tcut{} held at $v$ along $[e_v,e_w]$ at $q^\circ$. As
$C_{q^\circ}\ne D_v$, $\alpha^{(v)}_{q^\circ}>0$; take a non-null
$B\in\Bfam$ in the band $D_{vw}\cap\{0<t^v_w<\alpha^{(v)}_{q^\circ}\}$ and
$E\coloneqq C_{q^\circ}+\cred_v(B)$. Then $E\preceq_vC_{q'}$ for some
$q'\in[e_v,q^\circ)$, so $E\in\mathcal R$ by down-closure, and its
$\Bfam$-neighbour $C_{q^\circ}$ lies in $\mathcal R$: a contradiction. A
complete world satisfies \Tcut{} along both of its edges, and every world is
the first member of two ordered pairs; so all three worlds are incomplete.
This is (R-a2$''$).
\end{proof}

Compared with Corollary~\ref{cor:C}, alternative (R-a) splits. Its part
(R-a1) is a failure of (E2) (far-tail debt not lighter than bounded debt) or
of \To{} along an approach whose selections give no credit beyond that of
the improved model. Proposition~\ref{prop:confT1} and
Lemma~\ref{lem:Tsrel}(d) show that (S) allows failures of \To{} of two
different kinds: in the K-world \To{} fails while (E2), \Tcc{} and \Tcut{}
hold, so failures of this kind are not part of (R-a1) or (R-a2$''$); in
$\mathrm{Conf}(1_{P_2})$ (E2) fails, at a (W)-world. \Tcc{} holds at every
world of every $\mathrm{Conf}(\theta)$, so no failure of \Tcc{} under (S) is
known. Whether $V$ undominated excludes (R-a1) is open. In
(R-a2$''$) the selection map is, near each of six cut points, non-monotone in
a strong sense: the nearby selections on the $e_v$-side give more credit and
more debt at $v$ than the selection at the cut. No configuration satisfying
(S) with a crossing cut is known. Alternative (R-b) is sharpened to
(R-b$'$), whose last statement replaces the violation of (SC) by an exact
finite cancellation (Section~\ref{sec:twoNC}).

\subsection{The extension problem}\label{sec:EXT}

\begin{enumerate}[label=(EXT),leftmargin=*]
\item Let $n=3$, assume (G) at every world and (S), and let world $v$ be a
preorder with (IB) that satisfies \To{} along both of its edges, with (EP)
on both of them; so world $v$ is fine (Lemma~\ref{lem:T1otools}(d)). Is
world $v$ a (W)-world?
\end{enumerate}
Under the hypotheses of (EXT), world $v$ is fine, it has an SI pair with
non-null debt part (Lemma~\ref{lem:sCN}(c) and its proof), and it has a
weakly AC pair with non-null credit part (proof of
Theorem~\ref{thm:twoNC}). So, by Corollary~\ref{cor:NCext}, world $v$ has a
full extension if and only if it satisfies (NC), and it is a (W)-world as
soon as (NC$_{\mathfrak W_v\cup\mathfrak F_v}$) holds for some admissible
family $\mathfrak F_v$: the only possible obstruction is an exact finite
cancellation among accepted change vectors. The answer is yes in the cases
of Corollary~\ref{cor:NCclasses}: if world $v$ satisfies (SC)
(Theorem~\ref{thm:SCext}), in particular if it is complete (CCF,
Theorem~15); if it satisfies (GC); if it is a filter-Pareto order; and if it
is contained in a complete preorder with strict (G) and (IB). It is also
yes, when $V$ is
undominated, if world $v$ compares every selected model with every regular
model and its debt and credit orders are complete. Indeed $V$ is then tied
at world $v$ with a selected model (CBOI, Lemma~1; TAW, (K2)), so $V$ is
comparable at $v$ with every regular model, and every pair is SI or weakly
AC; so (SC) holds. (Such a world also satisfies \To, by
Section~\ref{sec:T1ocond}.) It is also yes if
$\preceq_v$ is contained in a preorder $\preceq'$ that has a full extension,
that is, a loss $L$ of CBOI's form with \eqref{eq:positive} such that
$B\preceq'X$ implies $L(B)\le L(X)$ for all regular $B,X$: then
$B\preceq_vX$ gives $B\preceq'X$ and so $L(B)\le L(X)$.

In general (EXT) is open. CCF's route needs a complete comparative
probability for Savage's and Wakker's theorem, and the reading of a pair
that is not SI as weakly AC; (SC) is exactly what this route needs, and
Corollary~\ref{cor:NCext} shows that neither is necessary. The order axioms
alone give neither (NC) (Example~\ref{ex:half5}) nor (NC$^+$)
(Example~\ref{ex:K4}), but the worlds of these examples are not fine; by Lemma~\ref{lem:smallcanc}, an exact cancellation
within one layer needs at least four accepted changes, two of which overlap
in the same direction on a non-null set. We know no fine world with (IB)
that violates (NC). By Corollary~\ref{cor:twoNC}, (NC) at two worlds makes
$V$ dominated under \Tos. By Corollary~\ref{cor:refined} and
Theorem~\ref{thm:twoNC}, if $V$ is undominated, (E2) and \Tcc{} hold at the
incomplete worlds, and two worlds have fine credit orders and satisfy (NC),
then the configuration has the crossing cuts of (R-a2$''$).

\subsection{One incomplete world without \texorpdfstring{\To}{(T1o)}}\label{sec:F4}

\begin{proposition}\label{prop:F4}
Let $n=3$, assume (G) at every world and (S), let every world be a preorder
with (IB), and let worlds $2$ and $3$ be complete. If $V$ is undominated,
then neither world~$2$ nor world~$3$ has the endpoint property on both of
its edges. Hence either $(e_2,e_3)$ carries a two-sided selection, or
(EP$_{23}$) holds and both $(e_1,e_2)$ and $(e_1,e_3)$ carry two-sided
selections.
\end{proposition}
\begin{proof}
If world~$2$ had the endpoint property on both of its edges, it would be a
(W)-world by \ref{R:CCF}, and world~$3$ would be complete, contrary to
Corollary~\ref{cor:Wforces}. The same holds with $2$ and $3$ exchanged. If
(EP$_{23}$) fails, some selection on $[e_2,e_3]$ omits both endpoints and
is two-sided. If (EP$_{23}$) holds, then (EP$_{21}$) and (EP$_{31}$) fail,
and the edges $[e_1,e_2]$ and $[e_1,e_3]$ carry two-sided selections.
\end{proof}

By Theorem~\ref{thm:Bstar}, $V$ is dominated if exactly one world is
incomplete and it satisfies (E2) and \Tcc{} along both of its edges, for
instance \Ts; by Proposition~\ref{prop:confT1}, \To{} itself can fail in a
configuration satisfying (S). Whether $V$ can be undominated when exactly
one world is incomplete and violates (E2) or \Tcc{} remains open. The proof
of Theorem~\ref{thm:Astar} needs (E2) at the incomplete world, world~$1$
say, for its debt yardsticks (Lemma~\ref{lem:T1otools}(b)). These enter
(LH$^+$) at every world: in the proof of Theorem~\ref{thm:LH}, a strongly
credit-null set at world~$2$ is refuted through Lemma~\ref{lem:sCN}(f) with
$(v,w)=(1,3)$, whose Case~II uses debt yardsticks at world~$1$, or with
$(v,w)=(3,1)$, which uses Lemma~\ref{lem:sCN}(e) at world~$1$ (Case~I needs
no form of (T1) by Proposition~\ref{prop:audit}(d)). They also enter
Theorem~\ref{thm:own}, through Lemma~\ref{lem:exchange} at world~$l$ and
Step~2 at world~$m$, whenever $1\in\{m,l\}$. \Tcc{} at world~$1$ enters
through Lemma~\ref{lem:sCN}(c), CCF's Lemma~5 and Step~3(i) of
Theorem~\ref{thm:rich}.

\subsection{A (W)-world with two incomplete partners}\label{sec:OW}

\begin{remark}\label{rem:OWgeneral}
Under (OW) with both worlds $1$ and $3$ incomplete, whether $V$ can be
undominated is open. What is known: (EP$_{12}$) and (EP$_{32}$) hold
(Theorem~\ref{thm:EP}); neither world~$1$ nor world~$3$ is a (W)-world
(\ref{R:TAW}); neither satisfies (SC) together with \Tw{} along $[e_1,e_3]$
(Theorem~\ref{thm:SC}). If (EP$_{13}$) holds, then worlds $1$ and $3$ have
the endpoint property on both of their edges, and neither satisfies (SC),
by Theorem~\ref{thm:SCext} and \ref{R:TAW}. If (EP$_{13}$) fails,
Proposition~\ref{prop:heavy3} and its mirror at world~$1$ give heaviness of
world~$3$ in $P_1$ and of world~$1$ in $P_3$, as non-memberships in ideals,
relative to credit sets in $P_2$ of small $\nu_{22}$-mass. If moreover
world~$2$ satisfies (IB) and \Tos{} holds, Theorem~\ref{thm:A} applies:
every edge has the endpoint property, and the configuration is of type
(R-b), so the residual is (EXT) at worlds $1$ and $3$. Without (IB) at
world~$2$ this reduction is not available, since the proof of
Theorem~\ref{thm:A} uses (IB) at every world.
\end{remark}

\begin{thebibliography}{9}
\bibitem{Savage} L.~J.~Savage, \emph{The Foundations of Statistics}, Wiley,
New York, 1954. Chapter~3, \S3.
\bibitem{Wakker} P.~Wakker, Agreeing probability measures for comparative
probability structures, \emph{Annals of Statistics} 9 (1981), 658--662.
\bibitem{Tychonoff} A.~Tychonoff, \"Uber die topologische Erweiterung von
R\"aumen, \emph{Mathematische Annalen} 102 (1930), 544--561.
\bibitem{vN} J.~von Neumann, Zur Theorie der Gesellschaftsspiele,
\emph{Mathematische Annalen} 100 (1928), 295--320.
\bibitem{YH} K.~Yosida and E.~Hewitt, Finitely additive measures,
\emph{Transactions of the American Mathematical Society} 72 (1952), 46--66.
\end{thebibliography}
\end{document}
